% concatenated from OCP-Nonlinear-Final.tex
\documentclass[11pt,a4paper]{article}

% -------------------- Page + typography --------------------
\usepackage[a4paper,margin=2.2cm]{geometry}

\usepackage[T1]{fontenc}
\usepackage[utf8]{inputenc} % keep for pdfLaTeX (remove for Xe/LuaLaTeX)
\usepackage{microtype}

% A more professional text+math font than Computer Modern
% (Comment these two lines out if you prefer Latin Modern.)
%\usepackage{newtxtext,newtxmath}

% Paragraph look (more readable than cramped no-indent w/ no parskip)
\setlength{\parindent}{0pt}
\setlength{\parskip}{6pt plus 2pt minus 1pt}

% -------------------- Math --------------------
\usepackage{mathtools}
\usepackage{amssymb,amsfonts,amsthm}
\usepackage{bm}
\usepackage{upgreek} % for \upmu

% If you need blackboard-bold 1, choose ONE:
% \usepackage{dsfont}
% \usepackage{bbm}

% -------------------- Figures + tables --------------------
\usepackage{graphicx}
\usepackage{caption}
%\usepackage{subcaption} % modern replacement for subfigure (subfigure is obsolete)
\usepackage{booktabs}
\usepackage{siunitx}

% If you truly need EPS support, uncomment:
% \usepackage{epstopdf}

% -------------------- Lists --------------------
\usepackage{enumitem}
%\setlist[itemize]{topsep=4pt,itemsep=2pt,parsep=0pt,partopsep=0pt}
%\setlist[enumerate]{topsep=4pt,itemsep=2pt,parsep=0pt,partopsep=0pt}

% -------------------- TikZ/PGFPlots (keep only once) --------------------
\usepackage{tikz}
\usetikzlibrary{
	calc,positioning,patterns,arrows.meta,
	decorations.pathreplacing,decorations.pathmorphing,
	angles,quotes,intersections,backgrounds,shadings
}
\usepackage{tikz-3dplot}
\usepackage{pgfplots}
%\pgfplotsset{compat=1.18}
\usepgfplotslibrary{fillbetween}
\tikzset{snake it/.style={decorate, decoration=snake}}

% -------------------- Code listings --------------------
\usepackage{xcolor}
\usepackage{listings}
\lstdefinestyle{clean}{
	basicstyle=\ttfamily\small,
	frame=single,
	framerule=0.3pt,
	breaklines=true,
	columns=fullflexible,
	keepspaces=true,
	showstringspaces=false
}
\lstset{style=clean}

% -------------------- Section styling --------------------
\usepackage{titlesec}
\titleformat{\section}{\Large\sffamily\bfseries}{\thesection}{0.8em}{}
\titleformat{\subsection}{\large\sffamily\bfseries}{\thesubsection}{0.8em}{}
\titlespacing*{\section}{0pt}{2.2ex plus .6ex minus .2ex}{1.0ex plus .2ex}
\titlespacing*{\subsection}{0pt}{1.6ex plus .4ex minus .2ex}{0.8ex plus .2ex}

% -------------------- Headers/footers --------------------
\usepackage{fancyhdr}
\pagestyle{fancy}
\fancyhf{}
\lhead{\nouppercase{\leftmark}}
\rhead{\thepage}
\renewcommand{\headrulewidth}{0.3pt}

% -------------------- Links + smart references --------------------
\usepackage{hyperref}
\usepackage[nameinlink,noabbrev]{cleveref}
\hypersetup{
	colorlinks=true,
	linkcolor=blue,
	citecolor=blue,
	urlcolor=blue,
	pdftitle={Your Article Title},
	pdfauthor={Your Name}
}
\usepackage{bookmark}

% -------------------- Appendix --------------------
\usepackage[toc,page]{appendix}

% -------------------- Draft helpers (toggle) --------------------
\newif\ifdraft
\draftfalse % set \drafttrue while writing

\ifdraft
\usepackage[notref]{showkeys}
\usepackage[colorinlistoftodos]{todonotes}
\else
\usepackage[disable]{todonotes}
\fi

% If you need \sout, keep ulem but avoid changing \emph behavior:
\usepackage[normalem]{ulem}

% -------------------- Bibliography (BibTeX) --------------------
\usepackage{cite}
\bibliographystyle{plain}




\newtheorem{notation}{Notation}[section]
\numberwithin{equation}{section} 
\newcommand\supp{\mathop{\rm supp}}
\newcommand\ds{\displaystyle}

\def\1{{\bf 1}}
\def\wo{\overline}
\def\clb{\color{blue}}
\def\clr{\color{red}}
\def\clm{\color{magenta}}
\def\clv{\color{violet}}
\def\clbr{\color{brown}}

\def\O{\Omega}
\def\o{\omega}
\def\G1{{\bf 1}}
\def\r{{\mathfrak{r}}}
\def\SS{\mathbb{S}}
\def\S{{\GS}}
\def\d{\delta}
\def\sc{\longrightarrow}

\def\e{\varepsilon}
\def\er{{\varepsilon r}}
\def\erd{{\varepsilon \delta r}}

\def\fv{\frak{v}}
\def\fu{\frak{u}}
\def\bE{\text{\bf E}}
\def\C{\mathcal{C}}
\def\id{\text{\bf id}}
\def\sk{\text{skew}}
\def\sy{\text{sym}}
\def\Te{{\cal T}_\varepsilon}
\def\Me{{\cal M}_\varepsilon}
\def\Mes{{\cal M}^{*}_\varepsilon}
\def\Qe{{\cal Q}_\varepsilon}
\def\Qes{{\cal Q}^{*}_\varepsilon}
\def\My{{\cal M}_Y }
\def\MY{{\cal M}_{\cal Y}}
\def\fTe{{\mathfrak{T}_\varepsilon}}


\def\fS{{\frak S}}
\def\fA{\frak A}
\def\fE{\frak E}
\def\fa{\frak a}
\def\fB{\frak B}
\def\fb{\frak b}
\def\fC{\frak C}
\def\fc{\frak c}
\def\fr{\frak r}
\def\fM{\frak M}
\def\a{{\alpha}}

\def\Ga{{\bf a}}
\def\Gb{{\bf b}}
\def\Ge{{\bf e}}
\def\Gf{{\bf f}}
\def\Gl{{\bf l}}
\def\Gr{{\bf r}}
\def\s{{\bf s}}
\def\Gt{{\bf t}}




\def\ds{\displaystyle}

\def\F{{\mathbb{F}}}
\def\K{{\mathbb{K}}}
\def\N{{\mathbb{N}}}
\def\M{{\mathbb{M}}}
\def\R{{\mathbb{R}}}
\def\Z{{\mathbb{Z}}}
\def\V{{\mathbb{V}}}
\def\A{{\mathbb{A}}}
\def\D{{\mathbb{D}}}
\def\R{{\mathbb{R}}}
\def\V{{\mathbb{V}}}
\def\W{{\mathbb{W}}}
\def\P{{\mathbb{P}}}
\def\N{{\mathbb{N}}}
\def\Z{{\mathbb{Z}}}
\def\U{\mathbb{U}}
\def\B{{\mathbb{B}}}

\def\wUU{\widetilde{\mathbb{U}}}


\def\GcA{\bm{\mathcal A}}
\def\GcU{\bm{\mathcal U}}
\def\GcR{\bm{\mathcal R}}
\def\GcD{\bm{\mathcal D}}
\def\GcB{\bm{\mathcal B}}
\def\GcH{\bm{\mathcal H}}
\def\GcV{\bm{\mathcal V}}
\def\GcL{\bm{\mathcal L}}
\def\GcN{\bm{\mathcal N}}
\def\GcZ{\bm{\mathcal Z}}
\def\GcF{\bm{\mathcal F}}
\def\GcM{\bm{\mathcal M}}
\def\GcW{\bm{\mathcal W}}
\def\GcS{\bm{\mathcal S}}


\def\rc{{\mathcal{r}}}
\def\Bc{{\mathcal{B}}}
\def\Ec{{\mathcal{E}}}
\def\Ic{{\mathcal{I}}}
\def\Qc{{\mathcal{Q}}}
\def\Sc{{\mathcal{S}}}
\def\Pc{{\mathcal{P}}}
\def\Kc{{\mathcal{K}}}
\def\Uc{{\mathcal{U}}}
\def\Rc{{\mathcal{R}}}
\def\Zc{{\mathcal{Z}}}
\def\Vc{{\mathcal{V}}}
\def\Tc{{\mathcal{T}}}
\def\Yc{{\mathcal{Y}}}
\def\Pc{{\mathcal{P}}}
\def\Wc{{\mathcal{W}}}
\def\Gc{{\mathcal{G}}}
\def\Jc{{\mathcal{J}}}
\def\Mc{{\mathcal{M}}}



\def\bmgamma{{\bm  \gamma}}
\def\wc{{\;\rightharpoonup\;}}
\def\ed{{\varepsilon\delta}}

\def\wh{\widehat }
\newcommand{\whh}[1]{\ensuremath{\widehat{\widehat{#1}}}}
\def\wt{\widetilde }
\def\X{\times }
\def\div{\,\hbox{div}\,}
\def\hat{\widehat}
\def\Kdot{\overset { \mbox{\huge .}} {{\cal K}}{}}
\def\Ga{{\bf a}}
\def\Gb{{\bf b}}
\def\Gc{{\bf c}}
\def\Ge{{\bf e}}
\def\Gf{{\bf f}}
\def\Gg{{\bf g}}
\def\Gh{{\bf h}}
\def\Gi{{\bf i}}
\def\Gj{{\bf j}}
\def\Gk{{\bf k}}
\def\Gl{{\bf l}}
\def\Gm{{\bf m}}
\def\Gn{{\bf n}}
\def\Go{{\bf o}}
\def\Gp{{\bf p}}
\def\Gq{{\bf q}}
\def\Gr{{\bf r}}
\def\Gs{{\bf s}}
\def\Gt{{\bf t}}
\def\Gu{{\bf u}}
\def\Gv{{\bf v}}
\def\Gw{{\bf w}}
\def\Gx{{\bf x}}
\def\Gy{{\bf y}}
\def\Gz{{\bf z}}

\def\GA{{\bf A}}
\def\GB{{\bf B}}
\def\GC{{\bf C}}
\def\GD{{\bf D}}
\def\GE{\mathbb{E}}
\def\GGE{{\bf E}}
\def\GF{{\bf F}}
\def\GG{{\bf G}}
\def\GH{{\bf H}}
\def\GI{{\bf I}}
\def\GJ{{\bf J}}
\def\GK{{\bf K}}
\def\GL{{\bf L}}
\def\GM{{\bf M}}
\def\GN{{\bf N}}
\def\GO{{\bf O}}
\def\GP{{\bf P}}
\def\GQ{{\bf Q}}
\def\GR{{\bf R}}
\def\GS{{\bf S}}
\def\GT{{\bf T}}
\def\GU{{\bf U}}
\def\GF{{\bf F}}
\def\GV{{\bf V}}
\def\GW{{\bf W}}
\def\GX{{\bf X}}
\def\GY{{\bf Y}}
\def\GZ{{\bf Z}}

\definecolor{skyblue}{RGB}{135,206,235}
\definecolor{deepskyblue}{RGB}{0, 191, 255}


%%%%%%%%%%%%%%%%%%%%%%%%%%%%%%%%%%%%%%%%%%

\usepackage{amsfonts}
\usepackage{graphicx}
\usepackage{epstopdf}
\usepackage{dsfont}
\usepackage{bm}
\usepackage{bbm}
\usepackage{algorithmic}
\usepackage{tikz}
\usepackage{pgfplots}
\usetikzlibrary{patterns}

\usepackage{nicefrac}
\usepackage{mathtools}

\usepackage{ulem}

%%%%%%%%%%%%%%%%%%%%%%%%%%%%%%%%%%%%%%%%%%%%%

\usepackage{url, breakurl}
\usepackage{charter} % for \FloatBarrier
\usepackage{longtable}
\usepackage{float}
\usepackage{supertabular}
\usepackage{graphicx}
\usepackage{color}
\usepackage{subfigure}
\usepackage{overpic}
\usepackage{textcomp}
\usepackage[section]{placeins}

\usepackage{enumitem}

\usepackage{siunitx}
\newtheorem{theorem}{Theorem}[section]
\newtheorem{lemma}[theorem]{Lemma}
\newtheorem{corollary}[theorem]{Corollary}
\newtheorem{definition}{Definition}
\newtheorem{proposition}{Proposition}
\newtheorem{assumption}{Assumption}
\newtheorem{problem}{Problem}
\newtheorem{remark}{Remark}
\newtheorem{note}{Note}


%%%%%%%%%%% End of definitions
\begin{document}
	%\begin{frontmatter}
	
	%\pretitle{}
	\title{Homogenization of an optimal control problem for nonlocal semilinear elasticity with soft inclusions}
	%\runtitle{Simultaneous homogenization and dimension reduction}
	%\subtitle{}
	
	% For one author:
	%\author{\inits{N.}\fnms{Name1} \snm{Surname1}\ead[label=e1]{first@somewhere.com}}
	%\address{Department first, \orgname{University or Company name},
		%Abbreviate US states, \cny{Country}\printead[presep={\\}]{e1}}
	
	% Two or more authors:
	%\begin{aug}
	%
	%\end{aug}
\author{%
	Amartya Chakrabortty\thanks{Corresponding author. Processes and Materials, Fraunhofer ITWM, Kaiserslautern 67663, Germany. \href{mailto:amartya.chakrabortty@itwm.fraunhofer.de}{amartya.chakrabortty@itwm.fraunhofer.de}}\hskip 10mm
%	Julia Orlik\thanks{Processes and Materials, Fraunhofer ITWM, Kaiserslautern 67663, Germany. \href{mailto:julia.orlik@itwm.fraunhofer.de}{julia.orlik@itwm.fraunhofer.de}}\hskip 10mm
	Abu Sufian\thanks{Department of Mathematics, Universidad de Concepcion, 4070409, Chile. \href{mailto:asufian@udec.cl}{asufian@udec.cl}}
}



	\date{}
	
	\maketitle
	{\bf Keywords:}
	Homogenization, Optimal Control Problem, Semi-linear Elasticity, High-Contrast Domain, Unfolding Operator.\\
	
	{\bf { Mathematics Subject Classification (2020):} 35B27, 35B40, 35J20 49J20, 49N20, 74B05.}
%	{\let\thefootnote\relax\footnotetext{
%			Amartya Chakrabortty,: SMS, Fraunhofer ITWM, Kaiserslautern 67663, Germany\\
%			Emails: amartya.chakrabortty@itwm.fraunhofer.de\\[1mm]
%			Abu Sufian: Department of Mathematics,  Universidad de Concepcion, 4070409, Chile \\
%			Email: asufian@udec.cl\\[1mm]
%			Corresponding author Email: amartya.chakrabortty@itwm.fraunhofer.de}}
	
	
	
	%			Corresponding author Email: amartya.chakrabortty@itwm.fraunhofer.de}}
%\begin{abstract}
%	This paper investigates the asymptotic analysis of an optimal control problem (OCP) defined on a domain with soft inclusions (high-contrast domain), governed by the nonlocal semi-linearized elasticity. The domain comprises a connected matrix phase and a soft inclusion region. Our analysis involves two independent parameters: the contrast $\delta>0$ and the periodicity $\varepsilon>0$. The paper addresses an asymptotic analysis of an OCP with $L^2$-tracking cost on the displacements, as $(\varepsilon, \delta) \to (0, 0)$ where the control only acts on the inclusion part.  For that, we first derive the homogenized problem as $(\varepsilon, \delta) \to (0, 0)$ for $\kappa\in(0,+\infty]$ where $\kappa = \lim_{(\e, \delta) \to (0, 0)} \frac{\delta}{\varepsilon}$. Finally, we derive the limit OCP and show that the limit optimal control is an optimal control of the limit OCP, using $\Gamma$-convergence.
%\end{abstract}
\begin{abstract}
	This paper investigates the asymptotic analysis of an optimal control problem (OCP) posed on a high-contrast elastic medium with soft periodic inclusions, governed by a semilinear elasticity system with a nonlocal term. The domain consists of a connected matrix phase and a soft inclusion phase. The model depends on two independent small parameters: the periodicity $\varepsilon>0$ and the contrast $\delta>0$, and the distributed control acts only in the inclusion region. We consider an $L^2$-tracking cost on the displacement and analyze the limit as $(\varepsilon,\delta)\to(0,0)$ in the regime
	\[ 
	 \lim_{(\varepsilon,\delta)\to(0,0)}\frac{\delta}{\varepsilon}=\kappa\in(0,+\infty].
	\]
	First, we derive the homogenized (limit) state system associated with this scaling. We then formulate the limit OCP and prove that the limit of the microscopic optimal controls is an optimal control for the limit problem, using a $\Gamma$-convergence approach.
\end{abstract}

\section*{Introduction}
Smart origami-inspired metamaterials \cite{gries} and advanced textiles \cite{orlik01} offer a platform for designing structures with
programmable deformation patterns. A common engineering mechanism is to combine a connected stiff skeleton
(e.g.\ panels or fibers) with comparatively soft regions (e.g.\ hinges, joints, or inclusions), so that the
overall response is governed by a pronounced contrast of stiffness across the microstructure. In many
applications, one seeks to actively tune the deformation by embedding actuation or pre-strain in the
soft phase, while the stiff phase primarily carries and transmits load. This naturally leads to models in
which the control/actuation is supported only in the soft inclusions.

A widely used modeling viewpoint for such materials is to encode pre-strain (or growth/swelling) through a
multiplicative decomposition of the deformation gradient into an elastic part and a non-elastic part, see \cite{prestrain02}.
Under small-deformation and small-prestrain assumptions, the resulting linearized equations for the
displacement contain an additive term generated by the pre-strain (see \cite{Amartya}), which can be recast as an effective
source term in the equilibrium equations (see \cite{GOW2}). In this work, we adopt this perspective and treat the resulting
effective source term as a distributed control acting only in the soft inclusions. This provides a
mathematically tractable framework for shape/response optimization while remaining consistent with common
linearization principles used in the mechanics of pre-strained or swelling media.

Beyond the contrast and microstructure, many soft active materials exhibit an additional feature: their
mechanical response may depend not only on the local strain, but also on a global measure of the
deformation amplitude. Such global coupling can arise, for instance, from an underlying interaction with a
surrounding medium (e.g.\ solvent redistribution in gels, network effects, or mean-field-type feedback),
or from modeling choices that penalize large overall deformation to reflect saturation/stiffening effects.
To capture this in a minimal way while keeping the elastic operator local, we include a nonlocal semilinear term of the form $\alpha\|u\|_{L^2(\Omega)}^p u$ with $p\geq 2$. Equivalently, this term is generated by adding
to the total energy a convex contribution proportional to $\|u\|_{L^2(\Omega)}^{p+2}$, so that the resulting
Euler--Lagrange equation couples the local elastic equilibrium with a global amplitude-dependent response.

Homogenization for partial differential equations with high-contrast coefficients has attracted sustained attention over the past decades, in particular for diffusion and conductivity-type models. Early contributions include \cite{MR620783,MR926151}, where strongly contrasting coefficients were analyzed in connection with effective properties of composite media. Since then, a broad literature has developed on high-contrast homogenization and the emergence of nonstandard (often nonlocal) effective models.

Beyond diffusion, high-contrast effects also arise in elasticity and more general variational settings. In \cite{MR1952924}, the macroscopic behavior of an elastic composite with a strong component embedded in a weak matrix was studied; the homogenized limit exhibits nonlocal features that depend sensitively on the geometry of the stiff phase. Related high-contrast phenomena for conductivity were investigated in \cite{MR3000657}, which considers a medium consisting of highly conductive vertical fibers embedded in a poorly conductive matrix. Further examples include \cite{MR3342695}, addressing an elliptic variational problem on a pillar-type oscillating domain with an insulating layer surrounding a highly conductive core, and \cite{MR3450748}, where a hyperbolic PDE with contrasting diffusive coefficients leads to a coupled two-scale limit system incorporating both macro- and micro-scale effects. For additional references on high-contrast homogenization we refer to \cite{gaudiello2021limit,briane2007homogenization,MR3251918} and the references therein.

The homogenization of optimal control problems (OCPs) governed by PDE constraints has also been extensively studied. For elliptic OCPs with oscillating coefficients (possibly appearing both in the state equation and in the cost functional), see for instance \cite{MR1466916,MR451153}. The homogenization of OCPs in porous media was investigated in \cite{MR1666365}. For OCPs with high-contrast diffusion coefficients, we refer to \cite{AbuCon,MR4552077}, where the limiting problems may exhibit nonlocal effects. In these works, the principal part of the state equation is given by the Laplacian; see also \cite{MR451153,durante2010asymptotic,bidhan01} and the references therein.

Homogenization of an approximate controllability problem (ACP) constrained by semilinear state systems was initiated in \cite{Conca01} on fixed domains, and later extended to perforated domains in the cases where the holes do not intersect the control region \cite{Conca02} and where the holes intersect the control region \cite{Conca03}. In these works, the semilinearity enters through a local term involving a $\C^1$-function $S_1:\R\to\R$ satisfying
\begin{equation}\label{N1}
S_1(0)=0,
\qquad
\exists\,\gamma>0 \ \text{such that}\ 
0\le \frac{S_1(s)}{s}\le \gamma,
\quad \forall\, s\in\R\setminus\{0\}.
\end{equation}
Homogenization of OCPs in domains with oscillating boundaries was studied in \cite{Nanda01,Nanda02,Nanda03,Nanda04}, where the control acts on the fixed (non-oscillating) part of the boundary/domain. In these works, the semilinearity is also introduced via a local term involving a $\C^2$-function $S_2:\R\to\R$ such that
\begin{equation}\label{N2}
S_2(0)=0,
\qquad
0<C_1\le S_2'(s)\le C_2,
\quad \text{and}\quad S_2'' \ \text{is bounded}.
\end{equation}
In the analysis of above OCPs, the indirect method, which concerns the
asymptotic behavior of the corresponding optimality systems is used. 
%we prefer to stay with
%the original optimal control problem and look for its homogenized limit, for which we
%establish a solution. This is typically referred to as the direct method, which is based
%on the concept of -convergence and its variational properties. From this, we immediately infer that the above mentioned variational principle (convergence of optimal
%solutions) holds true for the ‘directly obtained’ homogenized problems.
In \cite{NL01}, homogenization study of a semilinear OCP is presented the direct method, relying on $\Gamma$-convergence arguments. Their semilinearity is assumed to satisfy a global growth/Lipschitz estimate of at most quadratic type,
together with monotonicity closedness. Moreover, since the leading operator is uniformly coercive/elliptic (i.e., no high-contrast degeneracy), the states admit uniform $H^1$-bounds, hence by the compact embedding strong $L^2$ convergence. This strong $L^2$-compactness is the key ingredient to pass to the limit in such semilinear terms.


This paper studies the asymptotic behavior of an optimal control problem (OCP) posed on a high-contrast domain in $\R^N$, governed by a semilinear elasticity system with a nonlocal (at least cubic) term, using the indirect method and a form of $\Gamma$-convergence. The high-contrast domain is a bounded domain with Lipschitz boundary $\O\subset \R^N$. It is decomposed into two parts: a connected stiff part $\O^2_\e$ with relative stiffness of order $1$, and a soft part $\O^1_\e$ with relative stiffness of order $\d^2$. The microscopic control $\theta_\e\in L^2(\O^1_\e)^N$ acts only on the soft part $\O^1_\e$ (extended by zero to $\O$).

The governing system (understood in the weak sense; see Section~\ref{S02} for details) reads, for $p\ge 2$ and $\alpha>0$,
\begin{equation}\label{Main02}
	\left\{\begin{aligned}
		-\nabla \cdot \sigma_\ed(u_\ed)+\alpha\|u_\ed\|^p_{ L^2(\O)}\, u_\ed&=f(x)+\1_{\O^1_\e}\theta_\e,\quad &&\text{in }\O,\\
		u_\ed&=0,\quad &&\text{on }\Gamma_0,\\
		\sigma_\ed(u_\ed)\Gn&=0,\quad&&\text{on }\Gamma_1,
	\end{aligned}\right.
\end{equation}
where $f=(f_i)$ denotes the body force.

The main aim of this paper is to analyze the asymptotic behavior, as $(\e,\d)\to(0,0)$, of an OCP that seeks to match a prescribed displacement profile using external forces supported in the soft inclusions. Let $u_d\in L^2(\O)^N$ be a given desired displacement. For each admissible control $\theta_\e$, we denote by $u_\ed(\theta_\e)\in\GU$ the unique weak solution of \eqref{Main02} (equivalently, of \eqref{MainW01}). We consider the reduced cost functional
\begin{equation*}
	\Gj_{\ed}(\theta_\e)
	=\frac12\int_\O |u_\ed(\theta_\e)-u_d|^2\,dx
	+\frac{\gamma}{2}\int_{\O^1_\e}|\theta_\e|^2\,dx,
\end{equation*}
where $\gamma>0$ is a regularization parameter. The OCP consists in finding an optimal control $\Theta_\e\in L^2(\O^1_\e)^N$ such that
\begin{equation*}
	\Gj_{\ed}(\Theta_\e)=\inf_{\theta_\e\in L^2(\O^1_\e)^N}\Gj_{\ed}(\theta_\e).
\end{equation*}

The results of the paper can be divided into two parts. First, we derive the limit (homogenized) state system associated with \eqref{Main02} in the asymptotic regime
\begin{equation}\label{R1}
	(\e,\d)\to(0,0),\quad \text{and}\quad \kappa=\lim_{(\e,\d)\to(0,0)}\frac{\d}{\e}\in(0,+\infty].
\end{equation}
Using periodic unfolding, suitable energy arguments and a form of $\Gamma$-convergence, we identify the limit displacement, the associated correctors, the limit total energy and characterize that the limit energy depends on the parameter $\kappa$. Moreover, in the case $\kappa=+\infty$ we obtain a fully scale-separated homogenized system. The main tool for asymptotic analysis is the periodic unfolding operator, which was introduced in \cite{CDG2} and further developed in \cite{CDG1}. For an introduction to homogenization and elasticity, see \cite{Olenik,Donato,twoscale1,twoscale2,CDG}.
 
Second, we analyze the asymptotic behavior of the microscopic OCP as $(\e,\d)\to(0,0)$ with $\d/\e\to\kappa$. We establish existence of microscopic optimal controls and prove convergence of the adjoint state and optimal controls. Moreover, we derive the limit OCP and show that any limit point of a sequence of microscopic optimal controls is an optimal control for the limit problem. The proof relies on a $\Gamma$-convergence approach for the reduced cost functionals, together with stability and strong-convergence properties for the unfolded state variables.

Finally, we emphasize a further novelty of the present work concerns the structure of the nonlocal semilinearity.
In several related contributions on homogenization of OCPs with semilinear state equations, the nonlinearity
enters through a scalar functional of the state combined with a scalar nonlinearity (see, for example,
the functions $S_1$ and $S_2$ in \eqref{N1}--\eqref{N2}), under global growth/monotonicity assumptions. In contrast,
our model involves the Hilbert-space operator
\[
u \mapsto \alpha\|u\|_{L^2(\Omega)}^{p}\,u,\quad p\geq 2
\]
which couples all spatial points through the global $L^2$-amplitude of the displacement.
As a consequence, passing to the limit in this term cannot be achieved from weak convergence alone; it requires
strong $L^2$ convergence (equivalently, convergence of $L^2$ norms), which we obtain through energy convergence
arguments in the unfolding framework. Moreover, the Fr\'echet derivative of this operator contains a
nonlocal rank-one contribution,
\[
D(\|u\|^p u)[h]=\|u\|^p h + p\|u\|^{p-2}(u,h)_{L^2(\Omega)}\,u,
\]
leading to additional global coupling in the linearized and adjoint systems.
These features, combined with the high-contrast elastic operator, make the analysis more delicate than in the
scalar diffusion setting. In particular, uniform $H^1$ bounds typically fail because the elastic energy degenerates
in the soft phase (due to the $\delta^2$ scaling). To overcome this difficulty, we introduce a suitable decomposition
(see Lemmas~\ref{L2}--\ref{L3}) and establish convergence of the total energy via a $\Gamma$-convergence argument.

The paper is organized as follows. General notation and the description of the domain are introduced in
Sections~\ref{N-1} and \ref{S02}, respectively. In Section~\ref{S03}, we prove existence, uniqueness, and
a priori estimates for the weak solution (equivalently, the minimizer) of the $\varepsilon$--$\delta$-dependent
semilinear elasticity system. In Section~\ref{S04}, we establish existence of an optimal control for the
$\varepsilon$--$\delta$-dependent OCP and show that the control-to-state mapping is globally Lipschitz. The main homogenization results for the state equation are presented in
Section~\ref{prototype}.  Finally, in
Section~\ref{S06}, we analyze the asymptotic behavior of the $\varepsilon$--$\delta$ OCP.


%%%%%%%%%%%%%%%%%%%%%%%%%%%%%%%%%%%%%%%%%%%%%%%%%%%%%%%%%%%%%%%%%%%%%%%%%%%%%%%%%%%%%%%%%%%%%%%%%%%%%%%%%%%%%%%%%%%%%%%%%%%%%%%%%%%%%%%%%%%%%%%%%%%%%%%%%%%%%%%%%%%%%%%%%%%%%%%%%%%%%%%%%%%%%%%%%%%%%%%%%%%%%%%%%%%%%%%%%%%%%%%%%%%%
%{\clm  So, we have three different limit regimes depending on the value of $\kappa$.
%\begin{itemize}
%	\item When $\kappa=0$ (see \eqref{L02} for the unfolded limit problem and \eqref{426} for the homogenized problem): The macroscopic behavior is not affected by the soft part. The limit homogenized energy is the same if we start with a perforated plate $\O^2_\e$ i.e. $\d=0$ with $\theta_\e+f=0$ a.e. on $\O^1_\e$ and pass to the limit as $\e\to0$. This limit regime is also attainable if we first pass $\d\to 0$ with $\e$ being fixed then $\e\to0$ (see Remark \ref{RR2}) or first pass $(\e,\d)\to(0,0)$ with $\kappa\in(0,+\infty)$ then pass $\kappa\to 0$ (see Remark \ref{ReM02}). In this regime, the optimal control $\Theta_\e$ is not present at the homogenized problem and using cell problem \eqref{5173}, we have the limit optimal control given by
%	$$\wh\Theta= -f,\quad \text{a.e. in $\O\X Y_1$}.$$
%	\item  When $\kappa=+\infty$ (see \eqref{L03} for the unfolded limit problem and \eqref{425} for the homogenized problem): The macroscopic behavior is not directly affected by the soft part in the sense that the homogenized problem is same if we start with $\d=1$ from the beginning except due to the presence of the limit optimal control $\wh \Theta$ in the homogenized problem. This regime is attainable if we first pass $\e\to0$ with $\d$ being fixed then $\d\to0$ (see Remark \ref{RR2}) or first pass $(\e,\d)\to(0,0)$ with $\kappa\in(0,+\infty)$ then pass $\kappa\to +\infty$ (see Remark \ref{ReM02}). 
%	\item When $\kappa\in(0,+\infty)$ (see \eqref{L01} for the unfolded limit problem and \eqref{425} for the homogenized problem): The macroscopic behavior is the same as in the case when $\kappa=+\infty$, but the unfolded limit problem is different. The limit regime can be achieved if we start with replacing $\d=\kappa\e$ with $\kappa\in (0,+\infty)$, then pass to the limit $\e\to0$. Asymptotic behavior like this is already observed for simultaneous homogenization dimension reduction of non-linear elasticity problem in \cite{Amartya}. %See Remark \ref{ReM02}, for understanding $\kappa\to+\infty$ and $\kappa\to 0$.
%\end{itemize}
%%%%%%%%%%%%%%%%%%%%%%%%%%%%%%%%%%%%%%%%%%%%%%%%%%%%%%%%%%%%%%%%%%%%%%%%%%%%%%%%%%%%%%%%%%%%%%%%%%%%%%%%%%%%%%%%%%%%%%%%%%%%%%%%%%%%%%%%%%%%%%%%%%%%%%%%%%%%%%%%%%%%%%%%%%%%%%%%%%%%%%%%%%%%%%%%%%%%%%%%%%%%%%%%%%%%%%%%%%%%%%%%%%%%%%%%
%Since, from the above conclusion and the homogenized problem \eqref{426}, we have that the limit optimal control is absent in macroscopic problem when $\kappa=0$.}\\

\section{Notation}\label{N-1}
Below, we give some general notations which will be used throughout the paper.
\begin{itemize}
	\item Let $N\geq 2$ and $Y=(0,1)^N$. Let $Y_1\subset Y$ bounded open connected set (domain) with Lipschitz boundary such that $\wo{Y_1}\subset Y$ and $Y_2= Y\setminus\overline{Y_1}$. Moreover, we assume that $Y_2$ is a bounded domain with Lipschitz boundary. We denote the interface by $\Lambda=\wo{Y_1}\cap \wo{Y_2} =\partial Y_1$. 
	\item  $\R^{N\X N}$ is the space of $N\times N$  real matrices, $\R_\sy^{N\X N}$ is the space of real $N\X N$ symmetric matrices.
	\item The strain tensor is given by
	$$e(u)={1\over 2}\big((\nabla u)^T+\nabla u\big),\quad \text{for all $u\in H^1(\O)^N$}.$$
	\item The boundary $\partial\O$ of $\O$ is partitioned into two disjoint parts $\Gamma_0$ and $\Gamma_1$ with the surface measure of $\Gamma_0$ being strictly positive. We set $\Gn(\cdot)$ the unit outward normal on $\partial\O$.
	\item $\alpha,\gamma > 0$ and $(i,j,k,l)\in \{1,2,\ldots, N\}^4$ (if not specified).
	\item We have used the convention ${1\over +\infty}=0$.
	\item By convention in all the estimates we simply write $L^2(\O)$ instead of $L^2(\O)^N$ or $L^2(\O)^{N\times N}$, we write the complete spaces when we give weak or strong convergence.
	\item In this paper, we use the Einstein convention of summation over repeated indices and we denote by $C$ a generic constant independent of $\e$ and $\d$.
\end{itemize}
\section{Domain description and initial setup}\label{S02}
Let $\O\subset\R^N$ be a bounded domain with Lipschitz boundary.
We define the sets $\O_\e^1$ and $\O^2_\e$ as follows:
$$\O_{\e}^1=\bigcup_{k\in \Kc_\e} (\e k+ \e Y_1),\quad \O_\e^2=\O\setminus \overline{\O^1_\e},\qquad  \Kc_\e =\big\{k\in \Z^N\; |\; \e(k+ Y)\subset \O\big\}.$$  Note that the holes $k\e+\e Y_1$, $k\in \Kc_\e$, do not intersect the boundary of $\O$.
We denote the interface by $\Lambda_\e=\wo{\O^1_\e}\cap \wo{\O^2_\e} =\partial \O^1_\e $.\\
Observe that $\O^2_\e$ is bounded domain with Lipschitz boundary and $\partial \O\subset \partial \O^2_\e$. We denote the indicator function by
$$\1_{\O^1_\e}(x)=\left\{\begin{aligned}
1,\quad &x\in \O^1_\e,\\
0,\quad &x\in\GR^N\setminus \O^1_\e.
\end{aligned}\right.\quad 
 \1_{\O\X Y_1}(x,y) =\left\{\begin{aligned}
1,\quad &(x,y)\in \O\X Y_1,\\
0,\quad &(x,y)\in\O\X(\GR^N\setminus Y_1).
\end{aligned}\right.$$
We define the two 4th order Hooke's tensors $A\in [L^\infty(Y)]^{N\X N\X N\X N}$ satisfying the following
\begin{itemize}
	\item Symmetric: $A_{ijkl}= A_{jikl}= A_{klij}$.
	\item Coercivity and  condition ($K>0$)
	\begin{equation}\label{Coe01}
		\begin{aligned}
			A_{ijkl}(y) S_{ij} S_{kl}\geq K| S|^2_F,\quad \text{for a.e. $y\in Y$ and for all $S\in \R^{N\X N}_\sy$}.
		\end{aligned}	
	\end{equation}
	Here $|\cdot|_F$ denotes the standard Frobenius norm. 
	\item We set
	$$A_{\e,ijkl}(x)= A_{ijkl}\left({x\over \e}\right),\quad  \text{for a.e. $x\in\O$}.$$
	\item Let $\GA_\ed=\d^2 A_\e \1_{\O^1_\e}+ A_\e(1-\1_{\O^1_\e})$.
	\item Let us define the stress tensor $\sigma_\ed$ by
	$$\sigma_{\ed,ij}(u)=\GA_{\ed,ijkl}e_{kl}(u),\quad \text{for all $u\in H^1(\O)^N$}.$$
\end{itemize}
\section{Existence of unique solution and a priori estimates}\label{S03}
The problem is to determine the displacement field $u_\ed=(u_{\ed,i})$ satisfying the system \eqref{Main02}. By virtue of the assumptions on $\GA_\ed$, we define the bilinear form $\Ga_\ed$ by
\begin{equation*}
\Ga_\ed(u,w)=\int_{\O}\sigma_{\ed,ij}(u)e_{ij}(w)\,dx,\quad \forall\,u,w\in H^1(\O)^N.
\end{equation*}
We also define the global nonlinear term as
$$g(u,w)=\alpha\|u\|^p_{L^2(\O)}\int_\O u\cdot w\,dx,\quad \forall\,u,w\in H^1(\O)^N.$$
So, the weak form of the problem \eqref{Main02} is given by: For a given function $f\in L^2(\O)^N$ and $\theta_\e\in L^2(\O_\e^1)^N$, find the solution ${u}_\ed\in \GU$ such that
\begin{equation}\label{MainW01}
\Ga_\ed(u_\ed,w)+g(u_\ed,w)=\ell_\e(w)=\int_\O f\cdot w\,dx+\int_{\O^1_\e}\theta_\e\cdot w\,dx,\quad \forall\,w\in \GU,
\end{equation}
where the set of admissible displacement is given by
$$\GU=\left\{u\in H^1(\O)^N\,|\; u=0\;\text{on $\Gamma_0$}\right\}.$$
The associated minimization problem is 
\begin{equation}\label{MainE01}
	\left\{\begin{aligned}
		&\text{Find $u_\ed\in\GU$ such that}\\
		&\hskip 3cm\Gm_\ed=\GT_\ed(u_\ed)=\inf_{w\in\GU}\GT_\ed(w),
	\end{aligned}\right.
\end{equation}
where the total energy to this nonlinear elasticity problem is given by
$$\GT_\ed(w)={1\over 2}\Ga_\ed(w,w)+{1\over p+2}g(w,w)-\ell_\e(w),\quad\text{for all $w\in\GU$}.$$
Since $\GT_{\e\d}$ is Fréchet differentiable and strictly convex on $\GU$ (for $\alpha>0$), its unique minimizer in $\GU$ coincides with the unique weak solution of \eqref{MainW01}.
\begin{theorem}\label{Th1}
	For every $\e,\d>0$, there exist a unique minimizer $u_\ed\in \GU$ to the problem \eqref{MainE01}, i.e.
	$$\GT_\ed(u_\ed)=\Gm_\ed=\inf_{w\in\GU}\GT_\ed(w).$$
\end{theorem}
\begin{proof}
	Observe that
	$$\Gm_\ed\leq \GT_\ed(0)=0,$$
	so, we have $\Gm_\ed\in [-\infty,0]$. 
	
	{\bf Step 1.} We prove $\Gm_\ed\in(-\infty,0]$.
	
	First using the assumption \eqref{Coe01} and Korn's inequality, we have for any $w\in\GU$
	\begin{equation}\label{l2}
		\Ga_\ed(w,w)\geq C(\d)\|w\|^2_{H^1(\O)}.
	\end{equation}
	The positive constant  $C(\d)$  depends on $\d$. 
		
	Using Cauchy-Schwarz and Poincare inequality, we have for $w\in \GU$
	\begin{equation}\label{l1}
		|\ell_\e(w)|\leq C(\|f\|_{L^2(\O)}+\|\theta_\e\|_{L^2(\O^1_\e)})\|w\|_{L^2(\O)}\leq C(\|f\|_{L^2(\O)}+\|\theta_\e\|_{L^2(\O^1_\e)})\|w\|_{H^1(\O)}.
	\end{equation}
	The positive constant $C$ is independent of $\e$ and $\d$.
	
	We also have
	$$g(w,w)=\alpha\|w\|^{p+2}_{L^2(\O)}.$$
	Observe that for all $w\in\GU$ satisfying
	\begin{equation}\label{l3}
		\GT_\ed(w)\leq \GT_\ed(0)=0,
	\end{equation}
	we have using \eqref{l2}--\eqref{l1}
	\begin{multline}\label{AP01}
		{C(\d)\over 2}\|w\|^2_{H^1(\O)}+{\a\over p+2}\|w\|^{p+2}_{L^2(\O)}\leq {1\over 2}\Ga_\ed(w,w)+{1\over p+2}g(w,w)\\
		\leq |\ell_\e(w)|\leq C(\|f\|_{L^2(\O)}+\|\theta_\e\|_{L^2(\O^1_\e)})\|w\|_{H^1(\O)},
	\end{multline}
	which imply
	$$\|w\|_{H^1(\O)}\leq C(\d)(\|f\|_{L^2(\O)}+\|\theta_\e\|_{L^2(\O^1_ \e)}),\quad |\ell_\e(w)|\leq C(\d).$$
	Then, the above inequalities give for all $w\in\GU$ satisfying \eqref{l3}
	\begin{equation}\label{l4}
		\|w\|_{H^1(\O)}\leq C(\d)(\|f\|_{L^2(\O)}+\|\theta_\e\|_{L^2(\O^1_\e)}),\quad	\text{and}\quad
		-k(\e,\d)\leq -\ell_\e(w)\leq \GT_\ed(w)\leq 0.
	\end{equation}
	The positive constant $C(\d)$ depend on $\d$ and $k(\e,\d)$ depends on $\d$ and $\theta_\e$.

	The above inequalities give $\Gm_\ed\in(-\infty,0]$.
	
	{\bf Step 2.} We prove the existence of minimizer.
	
	Note the following: The quadratic elastic part $w\mapsto\Ga_\ed(w,w)$ is convex and continuous on $\GU$, hence weakly lower semicontinuous. The nonlinear term $w\mapsto g(w,w)$ is convex functional on a Hilbert space, so weakly lower semicontinuous. The right hand side term $-\ell_\e(w)$ is weakly continuous. Thus, $\GT_\ed$ is weakly lower semicontinuous on $\GU$. 
	
	Using the estimate \eqref{l4}, there exist a minimizing sequence $\{w_n\}_n\subset \GU$ satisfying \eqref{l3} and
	$$\Gm_\ed=\inf_{w\in\GU}\GT_\ed(w)=\liminf_{n\to+\infty}\GT_\ed(w_n).$$
	We also have the following convergence at-least for a subsequence (denoted by the same subscript)
	$$w_n\rightharpoonup u_\ed,\quad\text{weakly in $H^1(\O)^N$}.$$
	Since, $\GT_\ed$ is weak lower semicontinuous, we obtain
	$$\GT_\ed(u_\ed)\leq \liminf_{n\to+\infty}\GT_\ed(w_n)=\Gm_\ed.$$
	So, we have existence of  minimizer in $\GU$ and the fact that $\Gm_\ed$ is a minimum by sequencial criteria.
	
	{\bf Step 3.} We prove uniquesness of the minimizer.
	
	Let us consider the operator 
	$$\langle F(u),w\rangle:=\Ga_\ed(u,w)+g(u,w),\quad\forall\,u,w\in\GU.$$
	Then, $u\in\GU\mapsto F(u)\in\GU^\ast$, where $\GU^\ast$ is the dual space of $\GU$.
	
	Observe that for $u,w\in\GU$, we have (using \eqref{l2})
	\begin{align*}
		&\langle F(u)-F(w),u-w\rangle\\
		&\hskip1cm=\Ga_\ed(u-w,u-w)+g(u,u-w)-g(w,u-w)\\
		&\hskip 1cm \geq C\|u-w\|^2_{H^1(\O)}.
	\end{align*}
	since
%	\begin{multline*}
%		0\leq g(u,u-w)-g(w,u-w)\\
%		={\alpha_\e\over 2}\left(\|u\|^p_{L^2(\O)}-\|w\|^p_{L^2(\O)}\right)\left(\|u\|^2_{L^2(\O)}-\|w\|^2_{L^2(\O)}\right)+{\alpha_\d\over 2}(\|u\|^p_{L^2(\O)}+\|w\|^p_{L^2(\O)})\|u-w\|^2_{L^2(\O)}.
%	\end{multline*}
	\begin{align*}
		0 &\le g(u,u-w)-g(w,u-w)
		= \alpha\big(\|u\|_{L^2(\Omega)}^{p}u-\|w\|_{L^2(\Omega)}^{p}w,\;u-w\big)_{L^2(\Omega)}\\
		&=\frac{\alpha}{2}\Big(\|u\|_{L^2(\Omega)}^{p}-\|w\|_{L^2(\Omega)}^{p}\Big)
		\Big(\|u\|_{L^2(\Omega)}^{2}-\|w\|_{L^2(\Omega)}^{2}\Big)
		+\frac{\alpha}{2}\Big(\|u\|_{L^2(\Omega)}^{p}+\|w\|_{L^2(\Omega)}^{p}\Big)
		\|u-w\|_{L^2(\Omega)}^{2}.
	\end{align*}
	
	Let $u_\ed,v_\ed$ be two minimizer of \eqref{MainE01} in $\GU$, then, we have from the above expression, we have
	\begin{equation*}
		0=\langle F(u_\ed)-F(v_\ed),u_\ed-v_\ed\rangle\geq C\|u_\ed-v_\ed\|^2_{H^1(\O)},
	\end{equation*}
	since $F(u_\ed)=F(v_\ed)=\ell_\e$ in $\GU^\ast$ using the variational form \eqref{MainW01}. So, we get $u_\ed=v_\ed$ a.e. in $\O$. Hence, there exist a unique minimizer (which also the weak solution to the problem \eqref{MainW01}) to the problem \eqref{MainE01}.
	
	This completes the proof.
\end{proof}	
%Then using the assumptions and Korn's inequality, we have the bilinear form $\Ga_\ed$ is continuous and coercive in $\GU$. Therefore, according to the Lax-Milgram lemma, we have for any $f\in L^2(\O)^N$ and $\theta_\e\in L^2(\O^1_\e)^N$, there exists a unique solution ${u}_\ed\in \GU$ for \eqref{MainW01}. The unique solution $u_\ed$ is also the minimizer for the elastic energy given as
%$$\Gm_\ed=\GT_\ed(u_\ed)=\min_{w\in \GU}\GT_\ed(w),$$
%where
%$$\GT_\ed(w)=\Ga_\ed(w,w)-\int_\O f\cdot w\,dx-\int_{\O^1_\e}\theta_\e\cdot w\,dx,\quad\text{for all $w\in\GU$}.$$
%$$\|u_\e\|_{H^1(\O)}\leq C\{\|f\|_{L^2(\O)}+\|g\|_{L^2(\Gamma_1)}+\|\theta_\e\|_{L^2(\O^2_\e)}\}.$$
The unique solution corresponding to the given $\theta_\e$ will be denoted by $u_\ed(\theta_\e)$.
\begin{lemma}[Preliminary estimates]\label{L2}
	Let $u\in\GU$ be a displacement, then there exist displacements $\Gu\in\GU$ and $\Gw$ in $H^1_0(\O^1_\e)^N$, extended by $0$ in $\O^2_\e$,  such that
	\begin{equation}\label{ExRF}
		u=\Gu+\Gw\quad\text{a.e. in $\O$}
	\end{equation}
	and 
	\begin{equation}\label{46}
		\begin{aligned}
			\|e(\Gu)\|_{L^2(\O)}&\leq C\|e(u)\|_{L^2(\O_\e^2)},\\
			\|e(\Gw)\|_{ L^2(\O)}&\leq C\|e(u)\|_{L^2(\O^1_\e)}+C\|e(u)\|_{L^2(\O_\e^2)},\\
			\|\Gw\|_{L^2(\O^1_\e)}+\e\|\nabla \Gw\|_{L^2(\O^1_\e)} &\leq C\e\|e( \Gw)\|_{L^2(\O)}.
		\end{aligned}
	\end{equation}
	Moreover, we have the following Korn-type inequality
	\begin{equation}\label{NK01}
		\|u\|_{L^2(\O)}\leq C\left(\e\|e(u)\|_{L^2(\O^1_\e)}+\|e(u)\|_{L^2(\O^2_\e)}\right).
	\end{equation}
	The constant(s) do not depend on $\e$ and $\d$.
\end{lemma}
\begin{proof}
	We recall a standard extension result from \cite{Olenik,CDG}, since $\O^2_\e\subset \O$ is bounded connected open set with Lipschitz boundary: There exists an extension operator ${\bm\Pc}_\e$ from $H^1(\O^2_\e)^N$ into $H^1(\O)^N$ satisfying
	\begin{equation}
		\begin{aligned}
			\forall w\in H^1(\O^2_\e)^N,\qquad &{\bm\Pc}_\e(w)\in H^1(\O)^N,\qquad {\bm \Pc}_\e(w)_{|\O^2_\e}=w,\\
			&\big\|e\big({\bm\Pc}_\e(w)\big)\big\|_{L^2(\O)}\leq C\|e(w)\|_{L^2(\O^2_\e)}.
		\end{aligned}
	\end{equation}
	The constant does not depend on $\e$ and $\d$.
	
	Let $u\in \GU$. Then, we denote by $\Gu={\bm\Pc}_\e\left({u}_{|\O^2_\e}\right)$ and $\Gw=u-\Gu$.
	
	The estimate \eqref{46}$_1$ is due to the previous extension result \eqref{ExRF}. Since $u=\Gu+\Gw$ and using \eqref{46}$_1$, we have \eqref{46}$_2$.\\
	We prove \eqref{46}$_2$. We define $\Gw_k(y)=\Gw(\e k+ \e y)$ for all $y\in Y_1$. Since $\Gw$ satisfy \eqref{46}$_2$, we have using Korn's inequality in the fixed domain $Y_1$
	$$\|\Gw_k\|_{L^2(Y_1)}+\|\nabla_y \Gw_k\|_{L^2(Y_1)}\leq C\|e_y(\Gw_k)\|_{L^2(Y_1)}.$$
	The constant does not depend on $k\in \Kc_\e$, it only depends on $Y_1$.
	Since
	$$\nabla_y\Gw_k(y)=\e\nabla\Gw(\e k+ \e y),\quad\text{for a.e. $y\in Y_1$},$$
	we obtain using change of variable
	$$\|\Gw\|^2_{L^2(\e k+\e Y_1)}+ \e^2\|\nabla \Gw\|^2_{L^2(\e k+\e Y_1)}\leq C \e^2\|e(\Gw)\|^2_{L^2(\e k+\e Y_1)}.$$
	Adding for all $k\in\Kc_\e$ give \eqref{46}$_3$. 
	
	Since $ u=\Gu+\Gw$ and using the estimates \eqref{46}, we obtain
	$$\|u\|_{L^2(\O)}\leq C\|e(\Gu)\|_{L^2(\O)}+C\e\|e(\Gw)\|_{L^2(\O)}\leq C\left(\e\|e(u)\|_{L^2(\O^1_\e)}+\|e(u)\|_{L^2(\O^2_\e)}\right).$$
	This completes the proof.
\end{proof}
As a consequence of the above lemma, we have the following Korn-type inequalities:
%\subsubsection{Korn-type inequalities}
% We assume that there exist a $K>0$ independent of $\e$ and $ \d$ such that
%\begin{equation}\label{Ass01}
%	\|f\|_{L^2(\O)}+\|\theta_\e\|_{L^2(\O^1_\e)}\leq K.
%\end{equation}
%{\clm \begin{remark}
		%    The uniform bound on $\theta_\e$ is not an assumption comes from the optimal control problem. Give detail.
		%\end{remark}}
		%As a consequence of the above assumption, we have
		\begin{lemma}\label{L3}
			The unique solution $u_\ed\in\GU$ corresponding to $\theta_\e\in L^2(\O_\e^1)^N$ of \eqref{MainW01} and let $\Gu_\ed(\theta_\e)$, $\Gw_\ed(\theta_\e)$ be the fields from Lemma \ref{L2}. Then, we have the following
			\begin{equation}\label{E01}
				\begin{aligned}
					\d\|e(u_\ed(\theta_\e))\|_{L^2(\O^1_\e)}+\|e(u_\ed(\theta_\e))\|_{L^2(\O^2_\e)}&\leq C\left(\|f\|_{L^2(\O)}+\|\theta_\e\|_{L^2(\O^1_\e)}\right),\\
					\|\Gu_\ed(\theta_\e)\|_{H^1(\O)}+\|e(\Gu_\ed(\theta_\e))\|_{L^2(\O)}&\leq C\left(\|f\|_{L^2(\O)}+\|\theta_\e\|_{L^2(\O^1_\e)}\right),\\
					\|\Gw_\ed(\theta_\e)\|_{L^2(\O)}+\e\|\nabla \Gw_\ed(\theta_\e)\|_{L^2(\O)}+\e\|e(\Gw_\ed(\theta_\e))\|_{L^2(\O)}&\leq {C\e\over \d}\left(\|f\|_{L^2(\O)}+\|\theta_\e\|_{L^2(\O^1_\e)}\right).
				\end{aligned}
			\end{equation}
%			Moreover, we also have
%			\begin{equation}
%				\a\|u_\ed\|^{p+1}_{L^2(\O)}\leq C\left(\|f\|_{L^2(\O)}+\|\theta_\e\|_{L^2(\O^1_\e)}\right)
%			\end{equation}
			The constant(s) do not depend on $\e$ and $\d$.
		\end{lemma}
		\begin{proof}
			Let $\theta_\e\in L^2(\O^1_\e)^N$ be the control. For simplicity, we drop the $(\theta_\e)$ in each displacements.
			
		In the main problem \eqref{MainW01}, taking $w= u_\ed$ and using the coercivity of $\GA_\ed$ in \eqref{Coe01} along with assumption on forces and \eqref{AP01} gives
		$$K\left(\d^2\|e(u_\ed)\|^2_{L^2(\O^1_\e)}+\|e(u_\ed)\|^2_{L^2(\O^2_\e)}\right)+\a\|u_\ed\|^{p+2}_{L^2(\O)}\leq C\left(\|f\|_{L^2(\O)}+\|\theta_\e\|_{L^2(\O^1_\e)}\right)\|u_\ed\|_{L^2(\O)}.$$
		The Korn-type inequality \eqref{NK01}
%		Since $ u_\ed=\Gu_\ed+\Gw_\ed$ and using the estimates \eqref{46}, we obtain
%		$$\|u_\ed\|_{L^2(\O)}\leq C\|e(\Gu_\e)\|_{L^2(\O)}+C\e\|e(\Gw_\e)\|_{L^2(\O)}\leq C\left(\e\|e(\wo{u}_\e)\|_{L^2(\O^1_\e)}+\|e(\wo{u}_\e)\|_{L^2(\O^2_\e)}\right),$$
%		which 
		imply
		\begin{multline*}
			\left(\d^2\|e(u_\ed)\|^2_{L^2(\O^1_\e)}+\|e(u_\ed)\|^2_{L^2(\O^2_\e)}\right)\\
			\leq C\left(\|f\|_{L^2(\O)}+\|\theta_\e\|_{L^2(\O^1_\e)}\right)\left(\e\|e(u_\ed)\|_{L^2(\O^1_\e)}+\|e(u_\ed)\|_{L^2(\O^2_\e)}\right).
		\end{multline*}
		Due to \eqref{R1}, there exist a $C>0$ independent of $\e$ and $\d$ such that
		\begin{equation}\label{R11}
			\e\leq C\d.
		\end{equation}
	 So, we obtain
		\begin{multline*}
			\left(\d\|e(u_\ed)\|_{L^2(\O^1_\e)}+\|e(u_\ed)\|_{L^2(\O^2_\e)}\right)^2
			\leq 2\left(\d^2\|e(u_\ed)\|^2_{L^2(\O^1_\e)}+\|e(u_\ed)\|^2_{L^2(\O^2_\e)}\right)\\
			\leq 2C\left(\|f\|_{L^2(\O)}+\|\theta_\e\|_{L^2(\O^1_\e)}\right)\left(\d\|e(u_\ed)\|_{L^2(\O^1_\e)}+\|e(u_\ed)\|_{L^2(\O^2_\e)}\right).
		\end{multline*}
%		which imply
%		$$\left(\d\|e(\wo{u}_\e)\|_{L^2(\O^1_\e)}+\|e(\wo{u}_\e)\|_{L^2(\O^2_\e)}\right)^2\leq 2C\left(\|f\|_{L^2(\O)}+\|\Theta_\e\|_{L^2(\O^1_\e)}\right)\left(\d\|e(\wo{u}_\e)\|_{L^2(\O^1_\e)}+\|e(\wo{u}_\e)\|_{L^2(\O^2_\e)}\right).$$
%		Finally, we get \eqref{E01}$_1$.
%		$$\d\|e(\wo{u}_\e)\|_{L^2(\O^1_\e)}+\|e(\wo{u}_\e)\|_{L^2(\O^2_\e)}\leq C\left(\|f\|_{L^2(\O)}+\|\Theta_\e\|_{L^2(\O^1_\e)}\right).$$
		The above inequality gives \eqref{E01}$_1$. Then, using the estimates \eqref{46}$_{3,4}$, we get
		$$\begin{aligned}
			 \d \|e(\Gw_\ed)\|_{L^2(\O^1_\e)} +\|e(\Gu_\ed)\|_{L^2(\O)}\leq C\left(\|f\|_{L^2(\O)}+\|\theta_\e\|_{L^2(\O^1_\e)}\right).
		\end{aligned}$$
		Now, since $\Gu_\ed$ and $\Gw_\ed$ satisfy the boundary conditions \eqref{46}$_{1,2}$, using Korn's inequality and the estimate \eqref{46}$_5$, we obtain
		$$\|\Gu_\ed\|_{H^1(\O)}+{\d \over \e}\|\Gw_\ed\|_{L^2(\O)}+\d\|\nabla \Gw_\ed\|_{L^2(\O)}\leq C\left(\|f\|_{L^2(\O)}+\|\theta_\e\|_{L^2(\O^1_\e)}\right).$$
		This completes the proof.
		\end{proof}
%%%%%%%%%%%%%%%%%%%%%%%%%%%%%%%%%%%%%%%%%%%%%%%%%%%%%%%%%%%%%%%%%%%%%%%%%%%%%%%%%%%%%%%%%%%%%%%%%%%%%%%%%%%%%%%%%%%%%%%%%%%%%%%%%%%%%%%%%%%%%%%%%%%%%%%%%%%%%%%%%%%%%%
\section{$\e,\d$-dependent optimal control problem}\label{S04}
In this section, we define the optimal control problem. Let use consider the cost functional 
$$\GJ_\e(u,\theta_\e)={1\over 2}\|u-u_d\|^2_{L^2(\O)}+{\gamma\over 2}\|\theta_\e\|^2_{L^2(\O^1_\e)},\quad \forall\,(u,\theta_\e)\in \GU\X \GV,$$
where $\GV=L^2(\O^1_\e)^N$. Let us define the reduced cost function as
$$\Gj_\ed(\theta_\e)=\GJ_\e(u_\ed(\theta_\e),\theta_\e),\quad \forall\,\theta_\e\in\GV.$$
The optimal control problem (OCP) reads as
\begin{equation}\label{OCPM01}
	\begin{aligned}
		&\Gi_\ed=\inf_{\theta_\e\in\GV}\big\{\Gj_\ed(\theta_\e)=\GJ_\e(u_\ed(\theta_\e),\theta_\e)\,|\,\text{subjected to $u_\ed(\theta_\e)$ satisfying \eqref{MainW01}}\big\},
	\end{aligned}
\end{equation}
where $u_\ed(\theta_\e)$ is the unique weak solution to the problem \eqref{MainW01}, corresponding to $\theta_\e$.
\begin{theorem}\label{Th2}
	For each $\e,\d>0$, the OCP \eqref{OCPM01} admits a solution.
\end{theorem}
\begin{proof}
As
		$$\Gi_\ed=\inf_{\theta_\e\in\GV}\GJ_\e(u_\ed(\theta_\e),\theta_\e)=\inf_{\theta_\e\in\GV}\Gj_\ed(\theta_\e),$$
		we have $\Gi_\ed\in [0,+\infty)$ since $\Gi_\ed\leq \Gj_\ed(0)<+\infty$. So, there exist a minimizing sequence $\{\theta_{\e,n}\}_{n}\subset\GV$ such that
		\begin{equation}\label{42}
				\Gi_\ed=\Gj_\ed(\theta_{\e,n})\leq \Gj_\ed(0),\quad \liminf_{n\to+\infty}\Gj_\ed(\theta_{\e,n})=\Gi_\ed.
			\end{equation}
		Since, using \eqref{NK01} and \eqref{E01}, we have
		$$\Gj_\ed(0)={1\over 2}\|u_\ed(0)-u_d\|^2_{L^2(\O)}\leq C,$$
		which imply using \eqref{l4}$_1$
		\begin{equation*}
				\|\theta_{\e,n}\|_{L^2(\O^1_\e)}\leq C ,\quad \|u_{\ed,n}(\theta_{\e,n})\|_{H^1(\O)}\leq C,
			\end{equation*}
		where $u_{\ed,n}(\theta_{\e,n})$ is the unique solution of \eqref{MainW01} corresponding to $\theta_{\e,n}$.
		The constant is independent of $n$. 
		
		So, we obtain there exist $\Theta_\e\in\GV$ and $u_\ed\in\GU$ such that
		\begin{equation*}
				\begin{aligned}
						\theta_{\e,n} &\rightharpoonup \Theta_\e,\quad &&\text{weakly in $L^2(\O^1_\e)^N$},\\
						u_{\ed,n} &\rightharpoonup u_\ed,\quad &&\text{weakly in $H^1(\O)^N$ and strongly in $L^2(\O)^N$}.
					\end{aligned} 
			\end{equation*}
		From the above convergences, we get for $w\in\GU$
		\begin{equation*}
				\lim_{n\to+\infty}\int_{\O^1_\e}\theta_{\e,n}\cdot w\,dx=\int_{\O^1_\e} \Theta_\e\cdot w,\quad \lim_{n\to+\infty}g(u_{\ed,n},w)=g(u_\ed,w),\quad \lim_{n\to+\infty}\Ga_\ed(u_{\ed,n},w)=a_\e(u_\e,w),
			\end{equation*}
		so, we get $u_\ed\in\GU$ is the unique solution corresponding to $\Theta_\e$, i.e. $u_\ed=u_\ed(\Theta_\e)$. Using the weakly lower semi continuity of $L^2$ and $H^1$-norm with \eqref{42}, we get 
		$$\Gj_\ed(\Theta_\e)\leq\liminf_{n\to+\infty}\Gj_\ed(\theta_{\e,n})=\Gi_\e.$$
		So, $\Theta_\e$ is an optimal control. This completes the proof.
\end{proof}
For $\a>0$, the uniqueness of the optimal control is not guaranteed in general, since the reduced problem is non-convex as the control to state map is nonlinear.
%%%%%%%%%%%%%%%%%%%%%%%%%%%%%%%%%%%%%%%%%%%%%%%%%%%%%%%%%%%%%%%%%%%%%%%%%%%%%%%%%%%%%%%%%%%%%%%%%%%%%%%%%%%%%%%%%%%%%%%%%%%%%%%%%%%%%%%%%%%%%%%%%%%%%%%%%%%%%%%%%%%%%%%%%%%%%%%%%%%%%%%%%%%%%%%%%%%%%%%%%%%%%%%%%%%%%%%%%%%%%%%%%%%
%\begin{proof}
%	Let us set
%	$$\Gm^\ast_\e=\inf_{\theta_\e\in\GV}\GJ_\e(u_\ed(\theta_\e),\theta_\e)=\inf_{\theta_\e\in\GV}\Gj_\e(\theta_\e).$$
%	Then, we have $\Gm^\ast_\e\in [0,+\infty)$ since $\Gm^\ast_\e\leq \GJ_\e(0)<+\infty$. So, there exist a minimizing sequence $\{\theta_{\e,n}\}_{n}\subset\GV$ such that
%	\begin{equation}\label{42}
%		\Gj_\e(\theta_{\e,n})\leq \Gj_\e(0),\quad \liminf_{n\to+\infty}\Gj_\e(\theta_{\e,n})=\Gm^\ast_\e.
%	\end{equation}
%	Since, using \eqref{NK01} and \eqref{E01}, we have
%	$$\Gj_\e(0)={1\over 2}\|u_\e(0)-u_d\|^2_{L^2(\O)}\leq C,$$
%	which imply 
%	\begin{equation*}
%		\|\theta_{\e,n}\|_{L^2(\O)}\leq C ,\quad \|u_{\ed,n}(\theta_{\e,n})\|_{H^1(\O)}\leq C,
%	\end{equation*}
%	where $u_{\ed,n}(\theta_{\e,n})$ is the unique solution of \eqref{MainW01} corresponding to $\theta_{\e,n}$.
%	The constant is independent of $n$. 
%	
%	So, we obtain there exist $\Theta_\e\in\GV$ and $u_\e\in\GU$ such that
%	\begin{equation*}
%		\begin{aligned}
%			\theta_{\e,n} &\rightharpoonup \Theta_\e,\quad &&\text{weakly in $L^2(\O)^N$},\\
%			u_{\e,n} &\rightharpoonup u_\e,\quad &&\text{weakly in $H^1(\O)^N$ and strongly in $L^2(\O)^N$}.
%		\end{aligned} 
%	\end{equation*}
%	From the above convergences, we get for $w\in\GU$
%	\begin{equation*}
%		\lim_{n\to+\infty}\int_\O\theta_{\e,n}w\,dx=\int_\O \Theta_\e\cdot u_\e,\quad \lim_{n\to+\infty}g_\e(u_{\e,n},w)=g_\e(u_\e,w),\quad \lim_{n\to+\infty}a_\e(u_{\e,n},w)=a_\e(u_\e,w),
%	\end{equation*}
%	so, we get $u_\e\in\GU$ is the unique solution corresponding to $\Theta_\e$. Using the weakly lower semi continuity of $L^2$ and $H^1$-norm with \eqref{42}, we get 
%	$$\Gj_\e(\Theta_\e)\leq\liminf_{n\to+\infty}\Gj_\e(\theta_{\e,n})=\Gm^\ast_\e.$$
%	So, $\Theta_\e$ is an optimal control. This completes the proof.
%\end{proof}	
%%%%%%%%%%%%%%%%%%%%%%%%%%%%%%%%%%%%%%%%%%%%%%%%%%%%%%%%%%%%%%%%%%%%%%%%%%%%%%%%%%%%%%%%%%%%%%%%%%%%%%%%%%%%%%%%%%%%%%%%%%%%%%%%%%%%%%%%%%%%%%%%%%%%%%%%%%%%%%%%%%%%%%%%%%%%%%%%%%%%%%%%%%%%%%%%%%%%%%%%%%%%%%%%%%%%%%%%%%%%%%%%%%%%%%%%%
 Let $\Theta_\e\in \GV$ be an  optimal control with $\wo u_\ed=u_\ed(\Theta_\e)\in\GU$ being the unique weak solution of \eqref{MainW01} corresponding to $\Theta_\e$.

Observe that $\Gj_{\ed}(\Theta_\e)\leq \Gj_{\ed}(0)$. Then, from the fact that $u_\ed(0)=\Gu_\ed(0)+\Gw_\ed(0)$ along with the estimates \eqref{E01}, we have
\begin{equation}\label{UBOC}
	\begin{aligned}
		\|\wo u_\ed-u_d\|^2_{L^2(\O)}+\gamma\|\Theta_\e\|^2_{L^2(\O^1_\e)}
		&\leq \|u_\ed(0)-u_d\|^2_{L^2(\O)}\\
		&\leq \|\Gu_\ed(0)\|^2_{L^2(\O)}+\|\Gw_\ed(0)\|^2_{L^2(\O)}+\|u_d\|^2_{L^2(\O)}\\
		&\leq {C\e^2\over \d^2}.
	\end{aligned}
\end{equation}
The constant is independent of $\e$ and $\d$.\\
Since $\kappa\in(0,+\infty]$, we have that there exists a $C>0$ independent of $\e$ and $\d$ such that (see also \eqref{R11})
\begin{equation}\label{UBOCC}
		\|\wo u_\ed-u_d\|_{L^2(\O)}+	\|\Theta_\e\|_{L^2(\O^1_\e)}\leq {C_1\e\over \d}\leq C.
\end{equation}		
% We denote by $\Gu_\e$ the unique solution of \eqref{MainW01} corresponding to $\Theta_\e$. We call $(\Gu_\e,\Theta_\e)$ the optimal pair. Observe that since $\Gj_\e(\Theta_\e)\leq \Gj_\e(0)$, we obtain using \eqref{34}
%\begin{equation}\label{43}
%	\|\Theta_\e\|_{L^2(\O)}\leq C.
%\end{equation}
%The constant is independent of $\e$.
We call $(\wo u_\ed,\Theta_\e)$ an optimal pair. 			Let us define the control to state map $\Sc_\ed:\GV\to\GU$ as
$$\Sc_\ed(\theta)=u_\ed(\theta),\quad \forall\,\theta\in \GV,$$
where $u_\ed(\theta)$ is the unique solution to \eqref{MainW01} corresponding to $\theta$.
\begin{theorem}\label{Th03}
	The mapping $\Sc_\ed$ is Lipschitz continuous from $\GV$ to $\GU$, i.e.~there exists a constant $L_\ed>0$ (dependent on $\e$ and $\d$) such that
	\begin{equation}\label{CS01}
		\|u^1_{\ed}-u^2_\ed\|_{H^1(\O)}\leq L_\ed\|\theta_\e^1-\theta_\e^2\|_{L^2(\O^1_\e)}.
	\end{equation}
	Moreover, there exist $L>0$ independent of $\e$ and $\d$ such that
	\begin{equation}\label{CS02}
		\|u^1_{\ed}-u^2_\ed\|_{L^2(\O)}\leq L\|\theta_\e^1-\theta_\e^2\|_{L^2(\O^1_\e)},
	\end{equation}
	where $\theta^i\in \GV$ and $\Sc_\ed(\theta^i)=u^i_\ed$ for $i=1,2$.
\end{theorem}
\begin{proof}
	Let $\theta_\e^1,\theta_\e^2\in \GV$ and denote by
	$u_\ed^i:=\Sc_\ed(\theta_\e^i)\in\GU$ $(i=1,2)$ the corresponding unique solutions of \eqref{MainW01}, i.e.
	\[
	\Ga_\ed(u_\ed^i,w)+\a\|u_\ed^i\|_{L^2(\O)}^{p}(u_\ed^i,w)
	=(f,w)+(\theta_\e^i,w)_{L^2(\O_\e^1)}
	\qquad\forall w\in\GU.
	\]
%	Set $\delta u:=u_\e^1-u_\e^2$ and $\delta\theta:=\theta_\e^1-\theta_\e^2$.
	Subtracting the two identities yields, for all $w\in\GU$,
	\begin{equation}\label{eq:diff_state}
		\Ga_\ed(u^1_\ed-u^2_\ed,w)
		+\a\big(\|u_\ed^1\|_{L^2(\O)}^{p}u_\ed^1-\|u_\ed^2\|_{L^2(\O)}^{p}u_\ed^2,\;w\big)
		=(\theta^1_\e-\theta^2_\e,w)_{L^2(\O_\e^1)}.
	\end{equation}
%	
%	\smallskip
%	\noindent\textbf{Monotonicity of the nonlinear operator.}
%	Define $\Phi:L^2(\O)^N\to\R$ by
%	\[
%	\Phi(u):=\frac{\a}{p+2}\,\|u\|_{L^2(\O)}^{p+2}.
%	\]
%	Then $\Phi$ is convex (for $p\ge 0$), and its G\^ateaux derivative satisfies
%	\[
%	\Phi'(u)[w]=\a\|u\|_{L^2(\O)}^{p}(u,w)
%	\qquad\forall u,w\in L^2(\O)^N.
%	\]
	The nonlinear operator $G(u)=\a\|u\|_{L^2(\O)}^{p}u$ for all $u\in\GV$ is monotone and satisfy
	\begin{equation}\label{eq:monotone_B}
		\big(G(u)-G(v),\,u-v\big)\ge 0
		\qquad\forall u,v\in \GV.
	\end{equation}
	Taking $w=u^1_\ed-u^2_\ed$ in \eqref{eq:diff_state}. Using \eqref{eq:monotone_B} we obtain
	\[
	\Ga_\ed(u^1_\ed-u^2_\ed,u^1_\ed-u^2_\ed)
	+\big(G(u_\ed^1)-G(u_\ed^2),\,u^1_\ed-u^2_\ed\big)
	=(\theta^1_\e-\theta^2_\e,u^1_\ed-u^2_\ed)_{L^2(\O_\e^1)},\]
	which imply
	\[\Ga_\ed(u^1_\ed-u^2_\ed,u^1_\ed-u^2_\ed)\le (\theta^1_\e-\theta^2_\e,u^1_\ed-u^2_\ed)_{L^2(\O_\e^1)}.\]
	Using the coercivity of $\Ga_\ed$ together with the inequality \eqref{NK01}, there exist a constant $C>0$ independent of $\e$ and $\d$, such that
	\begin{multline*}
		C\|u^1_\ed-u^2_\ed\|^2_{L^2(\O)}\leq \Ga_\ed(u^1_\ed-u^2_\ed,u^1_\ed-u^2_\ed)\le (\theta^1_\e-\theta^2_\e,u^1_\ed-u^2_\ed)_{L^2(\O_\e^1)}\\
		\leq C\|\theta^1_\e-\theta^2_\e\|_{L^2(\O^1_\e)}\|u^1_\ed-u^2_\ed\|_{L^2(\O)},
	\end{multline*}
	which give the estimate \eqref{CS02} with constant independent of $\e$ and $\d$. Similarly, using the coercivity of $\Ga_\ed$ together with Korn's inequality and continuous embedding (restriction estimate) $H^1(\O)^N\hookrightarrow L^2(\O)^N$, we get
	$$C(\e,\d)\|u^1_\ed-u^2_\ed\|^2_{H^1(\O)}\leq C\|\theta^1_\e-\theta^2_\e\|_{L^2(\O^1_\e)}\|u^1_\ed-u^2_\ed\|_{H^1(\O)},$$
	which imply \eqref{CS01} with constant dependent on $\e$ and $\d$.
%	
%	Assume the coercivity of $\Ga_\ed$ on $\GU$: there exists $c_\e>0$ such that
%	\[
%	\Ga_\e(v,v)\ge c_\e\|v\|_{\GU}^2\qquad\forall v\in\GU,
%	\]
%	and assume the continuous embedding (restriction estimate) $\GU\hookrightarrow L^2(\O_\e^1)^N$:
%	there exists $C_E>0$ such that
%	\[
%	\|v\|_{L^2(\O_\e^1)}\le C_E\|v\|_{\GU}\qquad\forall v\in\GU.
%	\]
%	Then, by Cauchy--Schwarz,
%	\[
%	c_\e\|\delta u\|_{\GU}^2
%	\le \Ga_\e(\delta u,\delta u)
%	\le \|\delta\theta\|_{L^2(\O_\e^1)}\,\|\delta u\|_{L^2(\O_\e^1)}
%	\le C_E\,\|\delta\theta\|_{L^2(\O_\e^1)}\,\|\delta u\|_{\GU}.
%	\]
%	If $\delta u\neq 0$, divide by $\|\delta u\|_{\GU}$; otherwise the estimate is trivial. Hence
%	\[
%	\|\delta u\|_{\GU}
%	\le \frac{C_E}{c_\e}\,\|\delta\theta\|_{L^2(\O_\e^1)}.
%	\]
	Therefore the control-to-state mapping $\Sc_\ed:\GV\to\GU$ is globally Lipschitz.
%	\[
%	\|\Sc_\e(\theta_\e^1)-\Sc_\e(\theta_\e^2)\|_{\GU}
%	\le \frac{C_E}{c_\e}\,\|\theta_\e^1-\theta_\e^2\|_{L^2(\O_\e^1)}.
%	\]
\end{proof}

\section{Homogenization of the high-contrast state system \eqref{MainW01}}\label{prototype}
\subsection{Unfolding operator}
The convergences are done via the periodic unfolding operator $\Te$ for homogenization in $\O$. 
Below we recall the definition of the periodic unfolding operator for functions defined in $\O$, respectively. For the complete properties of the operator, see \cite{CDG}, specifically see Proposition 1.12, Theorem 1.36, Corollary 1.37 and Proposition 1.39 in \cite{CDG}. \begin{definition}
	For every measurable function $\phi$ on $\O$ the unfolding operator $\Te$ is defined by 
	$$ \Te(\phi) (x,y) \doteq \left\{ \begin{aligned}
		&\phi \left(\e\left[{x\over \e}\right] + \e y\right) \qquad &&\hbox{for a.e. }\; (x,y)\in  \wh\O_\e\times Y,\\
		&0 \quad &&\hbox{for a.e. }\; (x,y)\in  (\O\setminus\wh\O_\e)\times Y,
	\end{aligned}\right.$$
\end{definition} 
where 
$$\wh \O_\e=\hbox{interior}\left\{\bigcup_{k\in\Kc_\e}(\e k+\e \wo Y)\right\}.$$
The unfolding operator is a continuous linear operator from $L^2(\O)$ into $L^2(\O\X Y)$ which satisfies
$$\|\Te(\phi)\|_{L^2(\O\X Y)}\leq C\|\phi\|_{L^2(\O)}\quad\text{for every}\quad \phi\in L^2(\O).$$
The constants $C$ do not depend on $\e$.
Moreover, for every $\psi\in H^1(\O)$ one has (see \cite[Proposition 1.35]{CDG})
$$ \nabla_y\Te(\psi)(x,y)=\e\Te(\nabla \psi)(x,y) \quad \hbox{a.e. in}\quad \O\X Y.$$
Since $\O$ is a bounded domain with Lipschitz boundary, we have
$$\left|\int_{\O}\psi\,dx-{1\over |Y|}\int_{\O\X Y}\Te(\psi)\,dxdy\right|\to 0,\quad\text{as $\e\to0$},$$
so, we have
\begin{equation}\label{uci}
	\int_{\O}\psi\,dx\simeq{1\over |Y|}\int_{\O\X Y}\Te(\psi)\,dxdy.
\end{equation}
%			The lemma below is a consequence of Proposition 1.12, Theorem 1.36, Corollary 1.37 and Proposition 1.39 in \cite{CDG}.
%			\begin{lemma}\label{A11}
	%				$ $
	%				\begin{enumerate}
		%					\item Let $\{v_\e\}_\e$ be a weakly convergent sequence in $L^2(\O)$ with weak limit $v\in L^2(\O)$ then, there exist a $\wh{v}\in L^2(\O\X Y)$ with $\Mc_{Y}(\wh{v})=0$ such that  upto a subsequence
		%					$$\Te(v_\e)\rightharpoonup v+\wh{v},\quad \text{weakly in $L^2(\O\X Y)$}.$$
		%					\item Let $\{v_\e\}_\e$ be a sequence in $H^1(\O)$ such that
		%					$$\|v_\e\|_{L^2(\O)}+\e\|\nabla v_\e\|_{L^2(\O)}\leq C.$$
		%					Then, we have
		%					$$\begin{aligned}
			%						v_\e &\rightharpoonup v,\quad &&\text{weakly in $L^2(\O)$},\\
			%						\e\nabla v_\e & \rightharpoonup 0,\quad&&\text{weakly in $L^2(\O)^N$},\\
			%						\Te( v_\e) & \rightharpoonup v+\wh v,\quad &&\text{weakly in $L^2(\O; H^1_{per,0}(Y))$},\\
			%						\Te(\e\nabla v_\e)=\nabla_y\Te(v_\e) & \rightharpoonup\nabla_y \wh v,\quad &&\text{weakly in $L^2(\O\X Y)^N$}.
			%					\end{aligned}$$
		%					\item Let $\{v_\e\}_\e$ be a weakly in convergent sequence in $H^1(\O)$ with weak limit $v\in H^1(\O)$. Then, there exist a $\wh{v}\in L^2(\O;H^1_{per,0}(Y))$ such that  upto a subsequence
		%					$$\begin{aligned}
			%						\Te(v_\e)&\to v,\quad &&\text{strongly in $L^2(\O\X Y)$},\\
			%						\Te(\nabla v_\e)&\rightharpoonup \nabla v+\nabla_y \wh{v}\quad&&\text{weakly in $L^2(\O\X Y)$}.
			%					\end{aligned}$$
		%				\end{enumerate}
	%			\end{lemma}	
\subsection{Two-scale limit system}
We define the following space
$$ \GH^1_{per}(Y)\doteq\left\{\phi\in H^1_{per}(Y)\,|\, \phi=0\text{ on $Y_2$}\right\}.$$
Let $f\in L^2(\O)^N$ and $(\wo u_\ed,\Theta_\e)$ be an optimal pair. Then, due to the inequality \eqref{UBOCC}, there exist a $C\in\R^+$ independent of $\e$ and $\d$ such that
\begin{equation}\label{OCC01}
	\|\Theta_\e\|_{L^2(\O^1_\e)}\leq C,\quad \text{and}\quad \quad \Te(\Theta_\e) \rightharpoonup \wh \Theta \text{ weakly in }L^2(\Omega\times Y_1)^N,  
\end{equation}
where $\Theta_\e\in L^2(\O^1_\e)^N$ and $\wh\Theta\in L^2(\Omega\times Y_1)^N$.  
Then, for the corresponding fields $(\wo\Gu_\ed,\wo\Gw_\ed)$, we have the following estimates using \eqref{E01}
\begin{equation}\label{EM01}
	\begin{aligned}
		\|\wo\Gu_\ed\|_{H^1(\O)}&\leq C,\qquad {\d\over \e}\|\wo\Gw_\ed\|_{L^2(\O^1_\e)}+\d\|\nabla \wo\Gw_\ed\|_{L^2(\O^1_\e)}\leq C.
	\end{aligned}
\end{equation}
The constant is independent of $\e$ and $\d$.

While analyzing the optimal control problem in Section \ref{S06}, observe that $ \Theta_\e\in L^2(\O_\e^1)^N$ is the optimal microscopic control.

As $(\e,\d)\to (0,0)$, with $\kappa\in(0,+\infty]$, we get\footnote{In all the lemmas and theorems below, we extract  a subsequence  of $\{\ed\}_{\e,\d}$ (still denoted by $\{\ed\}_{\e,\d}$)  in order to get the desired convergences.}
\begin{lemma}\label{L04}
	There exist $\Gu_0\in H^1(\O)^N$ and $\wh U\in L^2(\O; H^1_{per,0}(Y))^N$ with
	$$\Gu_0=0,\quad \text{on $\Gamma_0$},$$
	such that
	\begin{equation}\label{Con01}
		\begin{aligned}
			\wo\Gu_\ed &\rightharpoonup \Gu_0,\quad&&\text{weakly in $H^1(\O)^N$},\\
			\Te(\wo\Gu_\ed)&\to \Gu_0,\quad &&\text{strongly in $L^2(\O\X Y)^N$},\\
			\Te(\nabla \wo\Gu_\ed) &\rightharpoonup \nabla \Gu_0+\nabla_y\wh U,\quad &&\text{weakly in $L^2(\O\X Y)^{N\X N}$}.
		\end{aligned}
	\end{equation}
	Moreover,
	there exist $\wh W\in L^2(\O; \GH^1_{per}(Y))^N$ such that
	\begin{equation}\label{Con02}
		\begin{aligned}
%			{\d\over \e} \wo\Gw_\ed &\rightharpoonup \Gw_0,\quad&&\text{weakly in $L^2(\O)^N$},\\
%			\d\nabla \wo\Gw_\ed&\rightharpoonup 0,\quad&&\text{weakly in $L^2(\O)^{N\X N}$},\\
			{\d\over \e}\Te(\wo\Gw_\ed) &\rightharpoonup \wh W,\quad&&\text{weakly in $L^2(\O\X Y)^N$},\\
			\d\Te(\nabla \wo\Gw_\ed) &\rightharpoonup \nabla_y\wh W,\quad&&\text{weakly in $L^2(\O\X Y)^N$}.
		\end{aligned}
	\end{equation}
	Furthermore, we have
	\begin{equation}\label{Con03}
		\begin{aligned}
%			\wo\Gw_\ed &\rightharpoonup {1\over \kappa}\Gw_0,\quad &&\text{weakly in $L^2(\O)^N$},\\
			\Te(\wo u_\ed) &\rightharpoonup \Gu_0+{1\over \kappa}\wh W,\quad&&\text{weakly in $L^2(\O\X Y)^N$}.
		\end{aligned}
	\end{equation}
\end{lemma}
\begin{proof}
	The convergences \eqref{Con01} are a direct consequence of the estimates \eqref{EM01}$_1$ along with the property of unfolding operator (see Corollary 1.37 in \cite{CDG}). The convergences \eqref{Con02} are a consequence of the estimate \eqref{EM01}$_2$ and of the properties of the unfolding operator given in Proposition 1.12, and Theorem 1.36 of \cite{CDG}. The convergences \eqref{Con03} are a consequence of \eqref{Con01}--\eqref{Con02} along with the fact that $\ds{\e\over \d} \to \ds{1\over \kappa}$.    This completes the proof.
\end{proof}
As a consequence of the above lemma, we get the limit two-scale problems. Before we give the limit problems, we define the following limit spaces
\begin{equation*}
	\begin{aligned}
		&\D=\GU\X L^2(\O;\GH^1_{per}(Y))^N\X L^2(\O; H^1_{per,0}(Y))^N,\quad \D_0=\GU \X L^2(\O;\GH^1_{per}(Y))^N.
	\end{aligned}
\end{equation*}
The two-scale limit energy is given by: for all $(w_0,W_1,W_2)\in\D$ and $\kappa\in(0,+\infty]$
\begin{equation}\label{EXE01}
	\GT_\kappa(w_0,W_1,W_2)={1\over 2}\Ga_1(W_1)+{1\over 2}\Ga_2(w_0,W_2)+{1\over p+2}g_\kappa(w_2,W_1)-\ell_\kappa(w_0,W_1)
\end{equation}
with
$$\begin{aligned}
	\Ga_1(W_1)&=\int_{\O\X Y}A(y)e_y(W_1):e_y(W_1)\,dydx,\quad 	g_\kappa(w_0,W_1)=\alpha\|w_0+{1\over \kappa}W_1\|^{p+2}_{L^2(\O\X Y)},\\
	\Ga_2(w_0,W_2)&=\int_{\O\X Y_2} A(y)\big(e(w_0)+e_y(W_2)\big):\big(e(w_0)+e_y(W_2)\big)\,dy dx,\\
	\ell_\kappa(w_0,W_1)&=\int_{\O\X Y}(f+\1_{\O\X Y_1}\wh\Theta)\cdot \left(w_0+{1\over \kappa} W_1\right)\,dydx.
\end{aligned}$$
%where $\wt{\wh \Theta}$ is the zero extension of $\wh \Theta$, defined in $\O\X Y_1$ to the whole two-scale domain $\O\X Y$.
\begin{theorem}\label{Th3}
	There exist a unique minimizer to the problem
	$$\Gm_\kappa=\min_{(w_0,W_1,W_2)\in\D}\GT_\kappa(w_0,W_1,W_2).$$
\end{theorem}
\begin{proof}
	First, we have
	$$\Gm_\kappa\leq \GT_\kappa(0,0,0)=0 \implies \Gm_\kappa\in[-\infty,0].$$
	Let us consider $(w_0,W_1,W_2)\in \D$ such that $\GT_\kappa(w_0,W_1,W_2)\leq 0$. Then, using the coercivity of $A$ with the inequality \eqref{Ex01} and Korn's inequality, we get
	$$ \|w_0\|^2_{H^1(\O)}+\|W_2\|^2_{L^2(\O;H^1(Y))}\leq C\Ga_2(w_0,W_2),\quad \|W_1\|^2_{L^2(\O;H^1(Y))}\leq C\Ga_1(W_1).$$
	We also have
	$$| \ell_\kappa(w_0,W_1)|\leq C(\|f\|_{L^2(\O)}+\|\wh\Theta\|_{L^2(\O\X Y_1)})\left(\sqrt{\Ga_2(w_0,W_2)}+\sqrt{\Ga_1(W_1)}\right),$$
	which together with the fact that $\GT_\kappa(w_0,W_1,W_2)\leq 0$ and $g_\kappa(w_0,W_1)\geq0$, give
	$$\Ga_1(W_1)+\Ga_2(w_0,W_2)\leq \sqrt{2}C(\|f\|_{L^2(\O)}+\|\wh\Theta\|_{L^2(\O\X Y_1)})\sqrt{\Ga_2(w_0,W_2)+\Ga_1(W_1)}.$$
	So, we obtain
	\begin{equation*}\label{LES01}
		\begin{aligned}
			\|w_0\|^2_{H^1(\O)}+\|W_2\|^2_{L^2(\O;H^1(Y))}+ \|W_1\|^2_{L^2(\O;H^1(Y))}&\leq  C\Ga_2(w_0,W_2)+C\Ga_1(W_1)\\
			&\leq C(\|f\|_{L^2(\O)}+\|\wh\Theta\|_{L^2(\O\X Y_1)})^2.
		\end{aligned}
	\end{equation*}
	Then, we have $\Gm_\kappa\in(-\infty,0]$. Then, proceeding as in Theorem \ref{Th1}, we get the existence of unique minimizer in $\D$. This completes the proof.
%	Observe that proceeding as in Theorem \ref{Th1} combined with the estimates \eqref{Ex01}, we have that there exist a 
\end{proof}
Finally, we have the convergence of the total energy.
\begin{theorem}\label{ENH01}
	We have 
	\begin{equation}\label{ECM01}
		\lim_{(\e,\d)\to(0,0)}\Gm_\ed=\Gm_\kappa=\GT_\kappa(\Gu_0,\wh W,\wh U).
	\end{equation}
\end{theorem}
\begin{proof}
	The proof is presented in $3$ steps using a form of $\Gamma$-convergence.
	
	{\bf Step 1.} We show that
	\begin{equation}
		\Gm_\kappa\leq \liminf_{(\e,\d)\to(0,0)}\Gm_\ed.
	\end{equation}
	Let $\{\wo u_\ed\}_{\e,\d}$ be the sequence of unique solutions of the elasticity problem \eqref{MainW01} in $\GU$. Then, we extend $(\wo u_\ed)_{|\O^2_\e}$ to $\wo\Gu_\ed\in H^1(\O)^N$ using the extension result \eqref{ExRF} and we have the estimates \eqref{46} and \eqref{E01} for $\wo\Gu_\ed$ and $\wo\Gw_\ed$. Then, with the assumption on the forces, we have the estimates \eqref{EM01}. Then, we get the following convergences at least for a subsequence (as in Lemma \ref{L04})
	\begin{equation}\label{II02}
		\begin{aligned}
			\Te\big(e(\wo{u}_\ed)\big) &\rightharpoonup e(\Gu_0)+e_y(\wh U),\quad&&\text{weakly in $L^2(\O\X Y_2)^{N\X N}$},\\
			\d\Te\big(e(\wo u_\ed)\big)& \rightharpoonup e_y(\wh W),\quad&&\text{weakly in $L^2(\O\X Y_1)^{N\X N}$},\\
%			\Te\big(\wo{u}_\ed\big) &\rightharpoonup \Gu_0,\quad&&\text{strongly in $L^2(\O\X Y_2)^{N}$},\\
%			\wo u_\ed& \rightharpoonup \Gu_0+{1\over \kappa}\Gw_0,\quad&&\text{weakly in $L^2(\O)^{N}$},\\
			\Te\big(\wo u_\ed\big)& \rightharpoonup \Gu_0+{1\over \kappa}\wh W,\quad&&\text{weakly in $L^2(\O\X Y)^{N}$},
		\end{aligned}
	\end{equation}
	with $(\Gu_0,\wh W,\wh U)\in \D$.
	
	So, first we transorm the energy $\GT_\ed$ using the unfolding operator (see the property \eqref{uci}), then the above convergences combined with \eqref{OCC01} and the weak lower semicontinuity of the functional $\GT_\ed$ give
	$$\Gm_\kappa\leq \GT_\kappa(\Gu_0,\wh W,\wh U)\leq \liminf_{(\e,\d)\to(0,0)}\GT_\ed(\wo u_\ed)=\liminf_{(\e,\d)\to(0,0)}\Gm_\ed.$$
	This completes Step 1.
	
	{\bf Step 2.} 	We prove
	\begin{equation}\label{He01}
		\begin{aligned}
			\limsup_{(\e,\d)\to(0,0)}\Gm_\ed\leq \GT_\kappa(w_0,W_1,W_2),\quad&&&\forall\, (w_0,W_1,W_2)\in\D.
		\end{aligned}	
	\end{equation}
	For that, we construct $\{w_\ed\}_{\e,\d}\subset \GU$ from $(w_0,W_1,W_2)\in\D$ such that
	\begin{equation}\label{59}
		\GT_\ed(\wo u_\ed)\leq \GT_\ed(w_\ed),\quad \lim_{(\e,\d)\to(0,0)}\GT_\ed(w_\ed)=\GT_\kappa(w_0,W_1,W_2).
	\end{equation}
	Observe that by density argument, it is enough to prove the above statement for $(w_0,W_1,W_2)\in\D\cap \C^1(\wo\O)^N\X \C^1(\wo{\O}\X \wo{Y})^N\X\C^1(\wo{\O}\X \wo{Y})^N$. 
	
	Let $(w_0,W_1,W_2)\in \D \cap \left(\C^1(\O)^N\X \C^1(\wo{\O}\X \wo{Y})^N\X\C^1(\wo{\O}\X \wo{Y})^N\right)$.  Then we set
	\begin{equation}\label{Test01}
		w_\ed(x)=\left\{\begin{aligned}
			&w_0(x)+\e W_2\left(x,\left\{{x\over \e}\right\}\right),\quad&&\text{for $x\in \O^2_\e$},\\
			&w_0(x)+\e W_2\left(x,\left\{{x\over \e}\right\}\right)+{\e\over \d}W_1\left(x,\left\{{x\over \e}\right\}\right),\quad&&\text{for $x\in\O^1_\e$},
		\end{aligned}\right.
	\end{equation} 
	which give
	$$\nabla w_\ed=\left\{\begin{aligned}
		&\nabla w_0+\e \nabla_x W_2+\nabla_y W_1,\quad&&\text{in $\O^2_\e$},\\
		&\nabla w_0+\e \nabla_x W_2+\nabla_y W_2+{\e\over \d}\nabla_x W_1+{1\over \d}\nabla_y W_1,\quad&&\text{in $\O^1_\e$},
	\end{aligned}\right.$$ 
	which imply
	$$\begin{aligned}
		\Te\big(e(w_\ed)\big)&\to e(w_0)+e_y(W_2),\quad &&\text{strongly in $L^2(\O\X Y_2)^{N\X N}$},\\
		\Te\big(\d e(w_\ed)\big) &\to e_y(W_1),\quad&&\text{strongly in $L^2(\O\X Y_1)^{N\X N}$},\\
		\Te(w_\ed) &\to w_0+{1\over \kappa}W_1,\quad &&\text{strongly in $L^2(\O\X Y)^N$}.
	\end{aligned}$$
	So, with the above above convergences and proceeding as in Step 1, we obtain \eqref{59}. Since $(w_0,W_1,W_2)\in\D$ is arbitrary, we obtain \eqref{He01}.

Finally, combining Step 1 and 2, we get
$$\Gm_\kappa\leq  \GT_\kappa(\Gu_0,\wh W,\wh U)\leq \liminf_{(\e,\d)\to(0,0)}\GT_\ed(\wo u_\ed)\leq \limsup_{(\e,\d)\to(0,0)}\GT_\ed(\wo u_\ed)\leq \Gm_\kappa,$$
by taking the unique minimizer of $\Gm_\kappa$ in place of $(w_0,W_1,W_2)\in\D$. Hence, we get the energy convergence \eqref{ECM01} and the fact that $(\Gu_0,\wh W,\wh U)\in\D$ is the unique minimizer of $\GT_\kappa$ in $\D$. Due to the existence of unique minimizer, we have that the convergences \eqref{II02} holds for the whole sequence.
This completes the proof.
\end{proof}
As a consequence of the above energy convergence, 
\begin{corollary}\label{Cor1}
	We have the following strong convergence for the whole sequence
	\begin{equation}\label{SCM}
		\Te\big(\wo u_\ed\big) \to \Gu_0+{1\over \kappa}\wh W,\quad\text{strongly in $L^2(\O\X Y)^{N}$}.
	\end{equation}
\end{corollary}
\begin{proof}
	
	From the convergence of total energy \eqref{ECM01} and the weak convergence \eqref{II02}, we have
	$$\lim_{(\e,\d)\to(0,0)}\left[{1\over 2}\Ga_\ed(\wo u_\ed,\wo u_\ed)+{1\over p+2}g(\wo{u}_\ed,\wo u_\ed)\right]={1\over 2}\left[\Ga_1(\wh W)+\Ga_2(\Gu_0,\wh U)\right]+{1\over p+2}g_\kappa(\Gu_0,\wh W),$$
	which combined with the weak lower-semicontinuity, non-negativity of $\Ga_\ed(\cdot,\cdot)$, $\|\cdot\|^{p+2}_{L^2(\O)}$ and the convergences \eqref{II02} give
	$$\liminf_{(\e,\d)\to(0,0)}\Ga_\ed(\wo u_\ed,\wo u_\ed)=\Ga_1(\wh W)+\Ga_2(\Gu_0,\wh U),\quad \liminf_{(\e,\d)\to(0,0)}g(\wo{u}_\ed,\wo u_\ed)=g_\kappa(\Gu_0,\wh W).$$
	Then, we get
	\begin{multline*}
		\limsup_{(\e,\d)\to(0,0)}{1\over p+2}g(\wo u_\ed,\wo u_\ed)=\limsup_{(\e,\d)\to(0,0)}\left[{1\over 2}\Ga_\ed(\wo u_\ed,\wo u_\ed)+{1\over p+2}g(\wo{u}_\ed,\wo u_\ed)-{1\over 2}\Ga_\ed(\wo u_\ed,\wo u_\ed)\right]\\
		=\limsup_{(\e,\d)\to(0,0)}\left[{1\over 2}\Ga_\ed(\wo u_\ed,\wo u_\ed)+{1\over p+2}g(\wo{u}_\ed,\wo u_\ed)\right]-{1\over 2}\liminf_{(\e,\d)\to(0,0)}\Ga_\ed(\wo u_\ed,\wo u_\ed)\\
		={1\over p+2}g_\kappa(\Gu_0,\wh W),
	\end{multline*}
	which imply
	$$\lim_{(\e,\d)\to(0,0)}\a\|\Te(\wo u_\ed)\|^{p+2}_{L^2(\O\X Y)}=\a\|\Gu_0+{1\over\kappa}\wh W\|^{p+2}_{L^2(\O\X Y)}.$$
Moreover, by \eqref{II02} we have
$T_\e(\bar u_{\e\delta}) \rightharpoonup \Gu_0+\frac1\kappa \wh W$ weakly in $L^2(\Omega\times Y)^N$.
Since $t\mapsto t^{p+2}$ is strictly increasing on $[0,\infty)$, the previous limit implies
$\|T_\e(\bar u_{\e\delta})\|_{L^2(\Omega\times Y)}\to \|\Gu_0+\frac1\kappa \wh W\|_{L^2(\Omega\times Y)}$.
Therefore, by uniform convexity of $L^2(\Omega\times Y)^N$, we conclude the strong convergence \eqref{SCM}.

%	Since, due to the definition of the unfolding operator, we have
%	$$\lim_{(\e,\d)\to(0,0)}\left|\int_\O \wo u_\ed\,dx-\int_{\O\X Y}\Te(\wo u_\ed)\,dydx\right|^2=0.$$
	This completes the proof.
\end{proof}
%%%%%%%%%%%%%%%%%%%%%%%%%%%%%%%%%%%%%%%%%%%%%%%%%%%%%%%%%%%%%%%%%%%%%%%%%%%%%%%%%%%%%%%%%%%%%%%%%%%%%%%%%%%%%%%%%%%%%%%%%%%%%%%%%%%%%%%%%%%%%%%%%%%%%%%%%%%%%%%%%%%%%%%%%%%%%%%%%%%%%%%%%%%%%%%%%%%%%%%%%%%%%%%%%%%%%%%%%%%%%%%%%%%%
%	Now, we give the limit of the LHS part in the main problem \eqref{MainW01} when $\kappa\in (0,+\infty]$. 
%	%	Observe that in vectorial form, we have using the unfolding operator
%	%	%				\begin{equation*}
%		%		%					\int_{\O^1_\e} \d^2 A_\e(x) \big(e(\Gu_\ed)+e(\Gw_\ed)\big): e(w_\ed)\,dx+\int_{\O^2_\e} A_\e(x) e(\Gu_\ed): e(w_\ed)\,dx=\int_\O f\cdot w_\ed\,dx+\int_{\O^1_\e}\Theta_\e\cdot w_\ed\,dx,
%		%		%				\end{equation*}
%	%	%				or,
%	%	\begin{multline*}
%		%		\int_{\O\X Y_1}A(y)\Te\big(\d e(\Gu_\e)+\d e(\Gw_\e)\big):\Te\big(\d e(w_\ed)\big)\,dx dy+\int_{\O\X Y_2} A(y) \Te\big(e(\Gu_\e)\big):\Te(e(w_\ed))\,dx dy\\
%		%		\|\Gu_0\|^p_{L^2(\O\X Y_2)}\int_{\O\X Y_2}\Gu_0\cdot w_0\,dydx=\int_{\O\X Y_2}\Te(f)\cdot\Te(w_\ed)\,dxdy+\int_{\O\X Y_1}\Te(f+\Theta_\e)\cdot\Te(w_\ed)\,dxdy,
%		%	\end{multline*}
%	As $(\e,\d)\to(0,0)$ with $\kappa\in(0,+\infty]$, we have limit LHS side given by
%	\begin{multline*}
%		\int_{\O\X Y_1} A(y)e_y(\wh W): e_y(w_2)\,dx dy+\int_{\O\X Y_2} A(y) \big(e(\Gu_0)+e_y(\wh U)\big): \big(e(w_0)+e_y(w_1)\big)\\
%		+\|\Gu_0\|^p_{L^2(\O\X Y_2)}\int_{\O\X Y_2}\Gu_0\cdot w_0\,dydx.
%		%					&\text{when $\kappa=0$},\quad&&\int_{\O\X Y_2} A(y)\big(e(\Gu_0)+e_y(\wh U)\big):\big(e(w_0)+e_y(w_1)\big)\,dx dy.
%	\end{multline*}
%	Since, we have
%	$$\begin{aligned}
%		\Te\big(\wo{u}_\ed\big) &\to \Gu_0,\quad&&\text{strongly in $L^2(\O\X Y_2)^{N}$},\\
%		\d\Te\big(\wo u_\ed\big)& \to 0,\quad&&\text{strongly in $L^2(\O\X Y_1)^{N}$}.
%	\end{aligned}$$
%	%	Since,
%	%	\begin{equation*}
%		%		\begin{aligned}
%			%			%		\1_{\O^1_\e}\Theta_\e & \rightharpoonup \Theta_0,\quad &&\text{weakly in $L^2(\O)^N$},\\
%			%			\Te(\Theta_\e)&\rightharpoonup \wh\Theta,\quad&&\text{weakly in $L^2(\O\X Y_1)^{N}$}.
%			%		\end{aligned}
%		%	\end{equation*}
%	Using the convergence \eqref{OCC01}, we have limit RHS side by passing to the limit in RHS of \eqref{MainW01}, which is given by
%	$$\begin{aligned}
%		&\text{when $\kappa\in(0,+\infty]$},\quad  &&\int_{\O\X Y_2} f\cdot w_0\,dx dy+\int_{\O\X Y_1}\big(f+\wh\Theta\big)\cdot (w_0+{1\over \kappa} w_2)\,dx dy,
%		%					&\text{when $\kappa=0$},\quad  &&\int_{\O\X Y_2} f\cdot w_0\,dx dy+\int_{\O\X Y_1}\big(f+\wh\Theta\big)\cdot w_0\,dx dy,
%	\end{aligned}$$
%	which gives \eqref{L01}, since $\D\cap \C^1(\O)^N\X \C^1(\wo{\O}\X \wo{Y_1})^N\X\C^1(\wo{\O}\X \wo{Y})^N$ is dense in $\D$.
%	
%	
%	
%	The corresponding energies of the unfolded limit problem \eqref{L01} is given by
%	\begin{multline*}
%		\GT_\kappa(w_0,w_1,w_2)={1\over 2}\Ga_1(w_2)+{1\over 2}\Ga_2(w_0,w_1)+{1\over p+2}g(w_0)\\
%		-|Y_2|\int_{\O} f\cdot w_0\,dx-\int_{\O\X Y_1}\big(f+\wh\Theta\big)\cdot \left(w_0+{1\over \kappa} w_2\right)\,dxdy,
%	\end{multline*}
%	where
%	$$\begin{aligned}
%		\Ga_1(w_2)&=\int_{\O\X Y_1}A(y) e_y(w_2):e_y(w_2)\,dxdy,\quad 		g(w_0)=\|w_0\|^{p+2}_{L^2(\O\X Y_2)},\\
%		\Ga_2(w_0,w_1)&=\int_{\O\X Y_2} A(y)\big(e(w_0)+e_y(w_1)\big):\big(e(w_0)+e_y(w_1)\big)\,dx dy,
%	\end{aligned}\quad \text{for all $(w_0,w_1,w_2)\in \D$}.$$	
%	Then, from Lemma \ref{L5} and Subsection \ref{SS44}, we have
%	\begin{equation*}				\begin{aligned}
%			\GT_\kappa(\Gu_0,\wh U,\wh W)=\min_{(w_0,w_1,w_2)\in\D}\GT_\kappa(w_0,w_1,w_2)=\min_{w_0\in\GU}\GT^{hom}(w_0)=\Gm.
%			%						&\text{for $\kappa=0$},\quad&&\GT_0(\Gu_0,\wh U)=\min_{(w_0,w_1)\in\GU\X L^2(\O; H^1_{per,0}(Y))^N}\GT_0(w_0,w_1)=\min_{w_0\in\GU}\GT^{hom}(w_0)=\Gm.
%		\end{aligned}
%	\end{equation*}
%
%	Observe that by density argument, it is enough to prove the above statement for $(w_0,w_1,w_2)\in\D\cap \C^1(\O)^N\X \C^1(\wo{\O}\X \wo{Y_1})^N\X\C^1(\wo{\O}\X \wo{Y})^N$. As in Lemma \ref{L5}, we have a sequence $\{w_\ed\}_{\e,\d}$ in $\GU$ with $\Gm_\ed\leq \GT_\ed(w_\ed)$ (see \eqref{Test01}). Then, we have \eqref{He01}
%	for all $(w_0,w_1,w_2)\in\D\cap \C^1(\wo{\O})^N\X \C^1(\wo{\O}\X \wo{Y_1})^N\X\C^1(\wo{\O}\X \wo{Y})^N$. This completes Step 1.\\[1mm]
%	We prove 
%	\begin{equation*}
%		\GT_\kappa(w_0,w_1,w_2)\leq \liminf_{(\e,\d)\to(0,0)}\Gm_\ed.
%	\end{equation*}
%	Observe that using the estimates\eqref{E01}, we have there exist a $c_1>0$ (independent of $\e$ and $\d$) such that
%	$$-c_1\leq \GT_\ed(\wo u_\ed)\leq \GT_\ed(0)=0.$$
%	So, as in Lemma \ref{L5}, we have the convergences \eqref{II02}. Then, we have with non-negativity of $\sigma_\ed(\cdot):e(\cdot)$ and weak lower semi continuity of $\GT_\ed$
%	$$	\GT_\kappa(w_0,w_1,w_2)\leq \liminf_{(\e,\d)\to(0,0)}\GT_\ed(\wo u_\ed).$$
%	This completes the proof.
%%%%%%%%%%%%%%%%%%%%%%%%%%%%%%%%%%%%%%%%%%%%%%%%%%%%%%%%%%%%%%%%%%%%%%%%%%%%%%%%%%%%%%%%%%%%%%%%%%%%%%%%%%%%%%%%%%%%%%%%%%%%%%%%%%%%%%%%%%%%%%%%%%%%%%%%%%%%%%%%%%%%%%%%%%%%%%%%%%%%%%%%%%%%%%%%%%%%%%%%%%%%%%%%%%%%%%%%%%%%%%%%%%%%%%%%%%

The unique minimzer $(\Gu_0,\wh W,\wh U)\in\D$ satisfy the following two-scale system (in variational form)
\begin{multline}\label{L01}
	\int_{\O\X Y_2} A(y)\big(e(\Gu_0)+e_y(\wh U)\big):\big(e(w_0)+e_y(W_2)\big)\,dx dy\\
	+\int_{\O\X Y}A(y) e_y(\wh W):e_y(W_1)\,dxdy+\alpha\|\Gu_0+{1\over \kappa}\wh W\|^p_{L^2(\O\X Y)}\int_{\O\X Y}(\Gu_0+{1\over \kappa}\wh W)\cdot (w_0+{1\over \kappa}W_1)\,dydx\\
	=\int_{\O\X Y}\big(f+\1_{\O\X Y_1}\wh\Theta\big)\cdot \left(w_0+{1\over \kappa} W_1\right)\,dydx,
\end{multline}
for all $(w_0,W_1,W_2)\in\D$.

%\begin{lemma}\label{L5}
%	Let $\{\wo{u}_\ed\}_{\e,\d}$ be the sequence of unique solutions of the elasticity problem \eqref{MainW01} in $\GU$ corresponding to $\Theta_\e\in L^2(\O^1_\e)^N$. Then, for the whole sequence $\{(\e,\d)\}_{\e,\d}$ and for $\kappa\in (0,+\infty]$, the following convergence holds:	
%	and $(\Gu_0,\wh U,\wh W)\in\D$ is the solution of the unfolded limit problem given by (in variational form):
%
%	%		where $F(x)=f(x)+\theta_0(x)$.
%	%				$\bullet$ for $\kappa=0$
%	%					{ \begin{multline}\label{L02}
%			%							\int_{\O\X Y_2} A(y)\big(e(\Gu_0)+e_y(\wh U)\big):\big(e(w_0)+e_y(w_1)\big)\,dx dy\\
%			%							=|Y_2|\int_{\O} f\cdot w_0\,dx+\int_{\O\X Y_1}\big(f+\wh\Theta\big)\cdot w_0\,dxdy,
%			%					\end{multline}}
%%	$\bullet$ for $\kappa=+\infty$
%%	{ \begin{multline}\label{L03}
%%			\int_{\O\X Y_1}A(y) e_y(\wh W):e_y(w_2)\,dxdy+\int_{\O\X Y_2} A(y)\big(e(\Gu_0)+e_y(\wh U)\big):\big(e(w_0)+e_y(w_1)\big)\,dx dy\\
%%			=|Y_2|\int_{\O} f\cdot w_0\,dx+\int_{\O\X Y_1}\big(f+\wh\Theta\big)\cdot w_0\,dxdy,
%%	\end{multline}}
%	for all $(w_0,w_2,w_1)\in\D$ with $\wh\Theta\in L^2(\O\X Y_1)^N$.
%\end{lemma}
%\begin{proof}
%	Let $\{\wo u_\ed\}_{\e,\d}$ be the sequence of unique solutions of the elasticity problem \eqref{MainW01} in $\GU$. Then, we extend $(\wo u_\ed)_{|\O^2_\e}$ to $\wo\Gu_\ed\in H^1(\O)^N$ using the extension result \eqref{ExRF} and we have the estimates \eqref{46} and \eqref{E01} for $\wo\Gu_\ed$ and $\wo\Gw_\ed$. Then, with the assumption on the forces, we have the estimates \eqref{EM01}. Then, the weak convergences \eqref{II02} at least for a subsequence follow from Lemma \ref{L04}.
%	
%	Now, We derive \eqref{L01}.
%	
%	Firstly, we define the sequence of test functions.
%	
%
%	
%	Finally, using the inequality \eqref{Ex01}, we have the unfolded limit problems \eqref{L01} have unique solutions using the Lax-Milgram lemma. Hence all the above convergences hold for the whole sequence.
%	%    Now, we prove the strong convergences \eqref{II02}. We have the weak convergences for the whole sequence. Since, $\sigma_\ed:e(u_\e)$ is no-negative, we have
%	%    $$\int_{\O\X Y}\Te\big(\sigm$$
%	%
%	%Due to the properties satisfied by $A^\alpha$ along with the Lemma \ref{NEq1} and the Lax-Milgram theorem, we have that the problem(s) \eqref{L01}--\eqref{L03} have a unique solution. This uniqueness implies that it is enough to prove the convergence \eqref{II02} for a subsequence.\\
%\end{proof}
\subsection{Homogenization and cell problems}\label{SS44}
In this subsection, we derive the homogenized limit problems for $\kappa\in(0,+\infty]$.

%{\bf Case 1:} When $\kappa\in(0,+\infty)$.\\[1mm]
In \eqref{L01}, we replace $w_0=0$ and $W_1=0$ with localizing, we obtain\footnote{Observe that $\wh U\in H^1_{per,0}(Y)^N$, since $\nabla_y\wh U$ enters the variational form, we can choose $\wh U\in H^1_{per,0}(Y_2)^N$.}
\begin{equation}\label{514}
	\begin{aligned}
		&\text{Find $\wh U\in H^1_{per,0}(Y_2)^N$ such that}\\
		&\hskip 10mm \int_{Y_2}A(y)(e(\Gu_0)+e_y(\wh U)):e_y(W_2)\,dy=0,\quad\text{for all $W_2\in H^1_{per,0}(Y_2)^N$}.
	\end{aligned}
\end{equation}
This shows that $\wh U$ can be expressed in terms of the elements of the tensor $e(\Gu_0)$ and some correctors. First, we define the following $N\X N$ symmetric matrices by
$$\GM^{ij}_{kl}={1\over 2}\left(\d_{ki}\d_{lj}+\d_{kj}\d_{li}\right),\quad i,j,k,l\in\{1,\dots, N\},$$
where $\d_{ij}$ is the Kronecker symbol. The cell problems are given by
\begin{equation*}
	\begin{aligned}
		&\text{Find $\chi^2_{ij}\in H^1_{per,0}(Y_2)^N$ such that}\\
		&\hskip 10mm \int_{Y_2} A(y)(\GM^{ij}+e_y(\chi^2_{ij})):e_y(W_2)\,dy=0,\quad\text{for all $W_2\in H^1_{per,0}(Y_2)^N$}.
	\end{aligned}
\end{equation*}
The above equation implies
\begin{equation}\label{Warping}
	\wh U(x,y)=\sum_{i,j=1}^N e_{ij}(\Gu_0)(x)\chi^2_{ij}(y),\quad\text{for a.e. $(x,y)\in \O\X Y_2$}.
\end{equation}
Observe that using the expression \eqref{Warping}, we have
\begin{equation}\label{Hom01}
	A^{hom}_{ijkl}=\int_{Y_2}\left(A_{ijkl}(y)\big(\GM^{kl}+e_y(\chi^2_{kl})\big)\right)\,dy.
%	[A^{hom}e(\Gu_0)]_{ij}=\int_{Y_2}A_{ijkl}(y) \left(e(\Gu_0)+\sum_{k,l=1}^N e_{kl}(\Gu_0)e_y(\chi^2_{kl})\right)\,dy\\
%	=\sum_{k,l=1}^N\left(\int_{Y_2}\left(A_{ijkl}(y)\big(\GM^{kl}+e_y(\chi^2_{kl})\big)\right)\,dy\right) e_{kl}(\Gu_0)=\sum_{k,l=1}^N A^{hom}_{ijkl}e_{kl}(\Gu_0).
\end{equation}
So, our limit problem \eqref{L01} becomes using \eqref{514} and \eqref{5171}
\begin{multline*}
	\int_{\O}A^{hom}e(\Gu_0):e(w_0)\,dx+\a\|\Gu_0+{1\over \kappa}\wh W\|^p_{L^2(\O\X Y)}\int_{\O\X Y}(\Gu_0+{1\over \kappa}\wh W)\cdot w_0\,dydx\\
	=\int_{\O}f\cdot w_0\,dx+|Y_1|\int_\O\Mc_{Y_1}(\wh\Theta)\cdot w_0\,dx.
\end{multline*}
Similarly, as in Theorem \ref{ENH01}, the homogenized limit stiffness tensor satisfies the following as a direct consequence of the hypothesis on $A_{ijkl}$ and the inequality \eqref{Ex01} (also see  Corollary 10.13 in \cite{Donato} and Theorem II.1.1 in \cite{Olenik}).
\begin{lemma}\label{HT01}
	We have $A^{hom}_{ijkl}\in\R$ and
	\begin{itemize}
		\item Symmetry: $A^{hom}_{ijkl}=A^{hom}_{jikl}=A^{hom}_{klij}$.
		\item Ellipticity and uniform boundedness: There exists  $0<C_1<C_2$ such that
		\begin{equation}\label{424}
			\begin{aligned}
				\text{for all $\zeta\in \R^{N\X N}_\sy$},&\quad A_{ijkl}^{hom}\zeta_{kl}\zeta_{ij}\geq C_1|\zeta|^2_F,\\
				\text{for all $\zeta\in \R^{N\X N}$},&\quad A_{ijkl}^{hom}\zeta_{kl}\leq C_2|\zeta|_F. 
			\end{aligned}
		\end{equation}
	\end{itemize}
\end{lemma}
%%%%%%%%%%%%%%%%%%%%%%%%%%%%%%%%%%%%%%%%%%%%%%%%%%%%%%%%%%%%%%%%%%%%%%%%%%%%%%%%%%%%%%%%%%%%%%%%%%%%%%%%%%%%%%%%%%%%%%%%%%%%%%%%%%%%%%%%%%%%%%%%%%%%%%%%%%%%%%%%%%%%%%%%%%%%%%%%%%%%%%%%%%%%%%%%%%%%%%%%%%%%%%%%%%%%%%%%%%%%%%%%%%%%%%%%%%%%%%%%
%% ------------------------------------------------------------
%% From the global variational formulation to the (global) alpha–cell problem
%% and the corrector representation of \widehat W.
%% ------------------------------------------------------------
%
%\begin{multline}\label{eq:L01}
%	\int_{\Omega\times Y_2} A(y)\big(e(\Gu_0)+e_y(\widehat U)\big):\big(e(w_0)+e_y(w_2)\big)\,dx\,dy\\
%	+\int_{\Omega\times Y}A(y) e_y(\widehat W):e_y(w_1)\,dx\,dy
%	+\Big\|\Gu_0+\frac1\kappa\widehat W\Big\|^p_{L^2(\Omega\times Y)}
%	\int_{\Omega\times Y}\Big(\Gu_0+\frac1\kappa\widehat W\Big)\cdot \Big(w_0+\frac1\kappa w_1\Big)\,dy\,dx\\
%	=\int_{\Omega\times Y}\big(f+\widetilde{\widehat\Theta}\big)\cdot \left(w_0+\frac1\kappa w_1\right)\,dy\,dx .
%\end{multline}
%
%% ------------------------------------------------------------
%% Step 1: isolate the \widehat W-equation by choosing w_0 = 0 and w_2 = 0.
%% ------------------------------------------------------------
%Choose $w_0=0$ and $w_2=0$ in \eqref{eq:L01}. Then for all admissible
%$w_1$ (e.g. $w_1\in L^2(\Omega;H^1_{per,0}(Y))^N$) we obtain
%\begin{equation}\label{eq:R_global}
%	\int_{\Omega\times Y}A(y) e_y(\widehat W):e_y(w_1)\,dx\,dy
%	+\Big\|\Gu_0+\frac1\kappa\widehat W\Big\|^p_{L^2(\Omega\times Y)}
%	\int_{\Omega\times Y}\Big(\Gu_0+\frac1\kappa\widehat W\Big)\cdot \frac1\kappa w_1\,dy\,dx
%	=\frac1\kappa\int_{\Omega\times Y}\big(f+\widetilde{\widehat\Theta}\big)\cdot w_1\,dy\,dx.
%\end{equation}
%
%% ------------------------------------------------------------
%% Step 2: define the (global) nonlinear scalar alpha and rewrite the problem.
%% ------------------------------------------------------------
%Define the global scalar
%\begin{equation}\label{eq:def_alpha}
%	\alpha \;:=\;\Big\|\Gu_0+\frac1\kappa\widehat W\Big\|^p_{L^2(\Omega\times Y)} \;\ge 0.
%\end{equation}
%Since $\alpha$ does not depend on $(x,y)$, \eqref{eq:R_global} becomes
%\begin{equation}\label{eq:R_alpha}
%	\int_{\Omega\times Y}A(y) e_y(\widehat W):e_y(w_1)\,dx\,dy
%	+\frac{\alpha}{\kappa}\int_{\Omega\times Y}\Big(\Gu_0+\frac1\kappa\widehat W\Big)\cdot w_1\,dy\,dx
%	=\frac1\kappa\int_{\Omega\times Y}\big(f+\widetilde{\widehat\Theta}\big)\cdot w_1\,dy\,dx.
%\end{equation}
%Expanding the term $\Gu_0+\frac1\kappa\widehat W$ and moving the $\Gu_0$–part to the right yields
%\begin{equation}\label{eq:R_alpha_linear}
%	\int_{\Omega\times Y}A(y) e_y(\widehat W):e_y(w_1)\,dx\,dy
%	+\frac{\alpha}{\kappa^2}\int_{\Omega\times Y}\widehat W\cdot w_1\,dy\,dx
%	=\frac1\kappa\int_{\Omega\times Y}\big(f+\widetilde{\widehat\Theta}-\alpha\Gu_0\big)\cdot w_1\,dy\,dx.
%\end{equation}
%
%% ------------------------------------------------------------
%% Step 3: localization in x (alpha stays global).
%% ------------------------------------------------------------
%Take test functions of the form $w_1(x,y)=\varphi(x)\psi(y)$ with
%$\varphi\in C_c^\infty(\Omega)$ and $\psi\in H^1_{per,0}(Y)^N$.
%From \eqref{eq:R_alpha_linear} and the fundamental lemma, we obtain:
%for a.e. $x\in\Omega$, $\widehat W(x,\cdot)\in H^1_{per,0}(Y)^N$ satisfies
%\begin{equation}\label{eq:cell_alpha}
%	\int_Y A(y) e_y(\widehat W(x,\cdot)):e_y(\psi)\,dy
%	+\frac{\alpha}{\kappa^2}\int_Y \widehat W(x,y)\cdot\psi(y)\,dy
%	=\frac1\kappa\int_Y\big(f(x)+\widetilde{\widehat\Theta}(x,y)-\alpha\Gu_0(x)\big)\cdot\psi(y)\,dy,
%	\qquad \forall \psi\in H^1_{per,0}(Y)^N.
%\end{equation}
%
%% ------------------------------------------------------------
%% Step 4: define alpha-dependent correctors and represent \widehat W.
%% ------------------------------------------------------------
%For each $\alpha\ge 0$, define the bilinear form on $V:=H^1_{per,0}(Y)^N$:
%\[
%a_\alpha(u,v):=\int_Y A(y)e_y(u):e_y(v)\,dy+\frac{\alpha}{\kappa^2}\int_Y u\cdot v\,dy.
%\]
%
%% (i) Basis correctors for constant right-hand sides in y
%For $i=1,\dots,N$, let $\chi_\alpha^i\in V$ be the (unique) solution of
%\begin{equation}\label{eq:chi_alpha}
%	a_\alpha(\chi_\alpha^i,\psi)=\int_Y e_i\cdot\psi\,dy \qquad \forall \psi\in V,
%\end{equation}
%where $(e_i)_{i=1}^N$ is the standard basis of $\mathbb{R}^N$.
%
%% (ii) Theta-corrector for the y-dependent source \widetilde{\widehat\Theta}
%Let $\zeta_\alpha(x,\cdot)\in V$ solve
%\begin{equation}\label{eq:zeta_alpha}
%	a_\alpha(\zeta_\alpha(x,\cdot),\psi)=\int_Y \widetilde{\widehat\Theta}(x,y)\cdot\psi(y)\,dy
%	\qquad \forall \psi\in V,
%	\quad\text{for a.e. }x\in\Omega.
%\end{equation}
%
%Write $f(x)=\sum_{i=1}^N f_i(x)e_i$ and $\Gu_0(x)=\sum_{i=1}^N(\Gu_0)_i(x)e_i$.
%By linearity of \eqref{eq:cell_alpha} in the right-hand side (for fixed $\alpha$),
%the solution admits the corrector representation
%\begin{equation}\label{eq:W_representation}
%	\widehat W(x,y)
%	=
%	\frac1\kappa\,\zeta_\alpha(x,y)
%	+\frac1\kappa\sum_{i=1}^N\Big(f_i(x)-\alpha(\Gu_0)_i(x)\Big)\chi_\alpha^i(y),
%	\qquad\text{for a.e. }x\in\Omega,\ \text{a.e. }y\in Y.
%\end{equation}
%
%% ------------------------------------------------------------
%% Step 5: the global fixed-point equation for alpha.
%% ------------------------------------------------------------
%Finally, $\alpha$ is not free; it must satisfy \eqref{eq:def_alpha}, i.e.
%\begin{equation}\label{eq:alpha_fixed_point}
%	\alpha
%	=
%	\left(
%	\int_\Omega\int_Y
%	\left|\Gu_0(x)+\frac1\kappa \widehat W(x,y)\right|^2\,dy\,dx
%	\right)^{p/2}.
%\end{equation}
%Inserting \eqref{eq:W_representation} into \eqref{eq:alpha_fixed_point} gives the explicit fixed-point map
%\begin{equation}\label{eq:alpha_fixed_point_explicit}
%	\alpha
%	=
%	\left(
%	\int_\Omega\int_Y
%	\left|
%	\Gu_0(x)
%	+\frac{1}{\kappa^2}\zeta_\alpha(x,y)
%	+\frac{1}{\kappa^2}\sum_{i=1}^N\Big(f_i(x)-\alpha(\Gu_0)_i(x)\Big)\chi_\alpha^i(y)
%	\right|^2\,dy\,dx
%	\right)^{p/2}.
%\end{equation}
%
%% Summary:
%%  - For each fixed \alpha, solve \eqref{eq:chi_alpha} once (i=1..N) and \eqref{eq:zeta_alpha} for each x,
%%    then obtain \widehat W by \eqref{eq:W_representation}.
%%  - Determine \alpha by the global fixed-point equation \eqref{eq:alpha_fixed_point_explicit}.
%%%%%%%%%%%%%%%%%%%%%%%%%%%%%%%%%%%%%%%%%%%%%%%%%%%%%%%%%%%%%%%%%%%%%%%%%%%%%%%%%%%%%%%%%%%%%%%%%%%%%%%%%%%%%%%%%%%%%%%%%%%%%%%%%%%%%%%%%%%%%%%%%%%%%%%%%%%%%%%%%%%%%%%%%%%%%%%%%%%%%%%%%%%%%%%%%%%%%%%%%%%%%%%%%%%%%%%%%%%%%%%%%%%%%%%%%%%%
Similarly, choosing $w_0=0$ and $W_2=0$ in \eqref{L01}, with localizing, we obtain
\begin{equation}\label{5171}
	\begin{aligned}
		&\text{Find $\wh W\in \GH^1_{per}(Y)^N$ such that for all $W_1\in \GH^1_{per}(Y)^N$}\\
		&\hskip 10mm \int_{Y} A(y)e_y(\wh W): e_y(W_1)\,dy+\alpha\|\Gu_0+{1\over \kappa}\wh W\|^p_{L^2(\O\X Y)}\int_{Y}(\Gu_0+{1\over \kappa}\wh W)\cdot {1\over \kappa}W_1\,dy\\
		&\hskip 100mm={1\over \kappa}\int_{Y}(f+\1_{\O\X Y_1}{\wh \Theta})\cdot W_1\,dy.
	\end{aligned}
\end{equation}
%The cell problems in $Y_1$ for $i=1,\ldots,N$ are given by
%\begin{equation*}
%	\begin{aligned}
%		&\text{Find $(\chi^1_{0},\chi^1_i)\in [H^1_{per,0}(Y_1)^N]^{N+1}$ such that}\\
%		&\hskip 10mm \int_{Y_1} A(y) e_y(\chi^1_i):e_y(w_2)\,dy=\int_{Y_1} w_2\cdot\Ge_i\,dy,\quad\text{for all $w_2\in H^1_{per,0}(Y_1)^N$},\\
%		&\hskip 10mm \int_{Y_1} A(y) e_y(\chi^{1}_0):e_y(w_2)\,dy=\int_{Y_1}\wh\Theta\cdot w_2\,dy,\quad\text{for all $w_2\in H^1_{per,0}(Y_1)^N$}
%	\end{aligned}
%\end{equation*}
%Finally, we have the expression
%\begin{equation}\label{421}
%	\wh W(x,y)={1\over \kappa}\left(\sum_{i=1}^N\big((f\cdot\Ge_i)(x)\chi^1_i(y)\big)+\chi^1_0(x,y)\right),\quad\text{for a.e. $(x,y)\in \O\X Y_1$}.
%\end{equation}
%Observe that $\chi^1_0\in L^2(\O; H^1_{per,0}(Y_1))^N$.\\
% where $|Y_2|f=F$ and $|Y_1|\theta_0=\theta$,\\
%with $\Mc_{Y_1}(\wh\Theta)=0$.\\
When $\kappa=+\infty$. Choosing $w_0=0$ and $W_2=0$ in \eqref{L01}, with localizing, we obtain 
\begin{equation}\label{5172}
	\begin{aligned}
		&\text{Find $\wh W\in \GH^1_{per}(Y)^N$ such that}\\
		&\hskip 10mm \int_{Y} A(y)e_y(\wh W): e_y(W_1)\,dy=0,\quad\text{for $W_1\in \GH^1_{per}(Y)^N$},
	\end{aligned}
\end{equation}
which has a unique solution using Lax-Milgram lemma and is given by $\wh W=0$. Similarly in \eqref{L01} we replace $w_0=0$ and $W_1=0$ with localizing, we obtain \eqref{514}.
%\begin{equation*}
%	\begin{aligned}
%		&\text{Find $\wh U\in H^1_{per,0}(Y_2)^N$ such that}\\
%		&\hskip 10mm \int_{Y_2}A(y)(e(\Gu_0)+e_y(\wh U)):e_y(W_2)\,dy=0,\quad\text{for all $W_2\in H^1_{per,0}(Y_2)^N$}.
%	\end{aligned}
%\end{equation*}
%So, our limit problem \eqref{L01} becomes using \eqref{514} and \eqref{5172}
%\begin{multline*}
%	\int_{\O\X Y_2}A(y)\big(e(\Gu_0)+e_y(\wh U)\big):e(w_0)\,dxdy+\alpha\|\Gu_0\|^p_{L^2(\O\X Y)}\int_{\O\X Y}\Gu_0\cdot w_0\,dydx\\
%	=\int_{\O}f\cdot w_0\,dx+|Y_1|\int_\O\Mc_{Y_1}(\wh\Theta)\cdot w_0\,dx.
%\end{multline*}
%So, 
% we have the limit homogenized problem given as in case 1,
%   \begin{equation}
	% \begin{aligned}
		% &\text{Find $\Gu_0\in \GH^1(\O)^N$ such that}\\
		% &\hskip 30mm        
		% \int_\O A^{hom}e(\Gu_0):e(w_0)\,dx=\int_{\O}\big(|Y|f+|Y_1|\theta_0+|Y_1|\Mc_{Y_1}(\wh\theta)\big)\cdot w_0\,dx,\quad \text{for all $w_0\in \GH^1(\O)^N$},
		% \end{aligned}
	% \end{equation}
% with $A^{hom}$ given by \eqref{Hom01}.\\
%			\noindent {\bf Case 3:} When $\kappa=0$.\\
%			Similarly in \eqref{L03} we replace $w_0=0$ with localizing, we obtain
%			\begin{equation*}
	%				\begin{aligned}
		%					&\text{Find $\wh U\in H^1_{per,0}(Y_2)^N$ such that}\\
		%					&\hskip 10mm \int_{Y_2}A(y)(e(\Gu_0)+e_y(\wh U)):e_y(w_1)\,dy=0,\quad\text{for all $w_1\in H^1_{per,0}(Y_2)^N$}.
		%				\end{aligned}
	%			\end{equation*}
\begin{remark}\label{ReM02}
	Note that the cell problems defined on $Y_2$ for any $\kappa\in(0,+\infty]$ are all the same, and this common problem is given by~\eqref{514}. Moreover, the two-scale limit problem~\eqref{L01} for $\kappa=+\infty$ can be obtained by first passing to the limit $(\varepsilon,\delta)\to(0,0)$ with $\kappa\in(0,+\infty)$ to derive~\eqref{L01}, and then passing to the limit $\kappa\to+\infty$.
\end{remark}
%Hence, using the above calculations along with Theorem \ref{ENH01} and Lemma \ref{HT01}, we have the limit homogenized problem (in variational form) for $\kappa\in(0,+\infty]$ and for all $(w_0,W_1)\in \D_0$
%\begin{multline*}
%	\int_\O A^{hom}e(\Gu_0):e(w_0)\,dx+\int_{\O\X Y}A(y) e_y(\wh W):e_y(W_1)\,dxdy\\
%	+\a\|\Gu_0+{1\over \kappa}\wh W\|^p_{L^2(\O\X Y)}\int_{\O\X Y}(\Gu_0+{1\over \kappa}\wh W)\cdot (w_0+{1\over \kappa}W_1)\,dydx\\
%	=\int_{\O\X Y}(f+1_{Y_1}\wh \Theta)\cdot (w_0+{1\over \kappa}W_1)\,dydx.
%\end{multline*}
% For $\kappa=+\infty$, the homogenized system (in variational form) is given by for all $w_0\in \GU$
%\begin{equation*}
%	\int_\O A^{hom}e(\Gu_0):e(w_0)\,dx+\a^{hom}\|\Gu_0\|^p_{L^2(\O)}\int_{\O}\Gu_0\cdot w_0\,dx
%	=\int_{\O}f\cdot w_0\,dx+|Y_1|\int_\O\Mc_{Y_1}(\wh\Theta)\cdot w_0\,dx,
%\end{equation*}
%where $\alpha^{hom}=\a|Y_2|^{{p\over2}+1}$. 
%
%\medskip
Finally, using the two-scale energy \eqref{EXE01} and the cell problems \eqref{514}, \eqref{5171}, \eqref{5172} with the expression \eqref{Hom01}, Theorem \ref{ENH01} and Lemma \ref{HT01}, we have the following
%\begin{theorem}
%	The limit displacement $\Gu_0\in \GU$ is the unique minimizer of the homogenized energy, i.e.
%	\begin{equation}\label{425}
%		\Gm_\kappa=\GT^{hom}_\kappa(\Gu_0),\quad \text{for $\kappa\in(0,+\infty]$},
%	\end{equation}
%	where
%	\begin{multline*}
%		\GT_\kappa^{hom}(\Gu_0)=\int_\O A^{hom}e(\Gu_0):e(\Gu_0)\,dx+{\a\over p+2}\|\Gu_0+{1\over\kappa}\wh W\|^{p+2}_{L^2(\O\X Y)}
%		-\int_{\O\X Y}\big(f+\wt{\wh\Theta}\big)\cdot \left(\Gu_0+{1\over \kappa} \wh W\right)\,dydx
%	\end{multline*}
%	with $A^{hom}$ is given by \eqref{Hom01} and $\wh W$ is the unique solution of \eqref{5171}.
%%	
%%	for all $w_0\in \GU$ and for $\kappa\in(0,+\infty]$, where $A^{hom}$ is given by \eqref{Hom01}, $\a=|Y_2|^{{p\over 2}+1}$ and
%%	$$\Te(\Theta_\e) \rightharpoonup \wh \Theta,\quad \text{weakly in $L^2(\O\X Y_1)^N$}.$$
%\end{theorem}
\begin{theorem}[Homogenized system]
	\label{56}
	For $\kappa\in(0,+\infty]$ and let $A^{hom}$ be given by \eqref{Hom01}. 
	Then the limit displacement $(\Gu_0,\wh W)\in \D_0$ is the unique minimizer of the homogenized energy functional
	$\GT^{hom}_\kappa:\D_0\to\R$, and
	\begin{equation}\label{425}
		\Gm_\kappa=\GT^{hom}_\kappa(\Gu_0,\wh W),
	\end{equation}
	where $\GT^{hom}_\kappa$ is given by
	\begin{multline}\label{Thm:HomEnergy}
		\GT^{hom}_\kappa(\Gu_0,\wh W)
		=
		\frac12\int_\O A^{hom}e(\Gu_0):e(\Gu_0)\,dx+{1\over 2}\int_{\O\X Y}A(y) e_y(\wh W):e_y(\wh W)\,dxdy\\
		+\frac{\a}{p+2}\Big\|\Gu_0+\frac1\kappa\wh W\Big\|_{L^2(\O\times Y)}^{p+2}
		-\int_{\O\times Y}\big(f+\1_{\O\X Y_1}\wh\Theta\big)\cdot\Big(\Gu_0+\frac1\kappa\wh W\Big)\,dx\,dy.
	\end{multline}
	Moreover, $(\Gu_0,\wh W)\in\D_0$ is the unique solution of 
		\begin{multline*}
		\int_{\O}A^{hom}e(\Gu_0):e(w_0)\,dx+\int_{\O\X Y}A(y) e_y(\wh W):e_y(W_1)\,dxdy\\
		+\a\|\Gu_0+{1\over \kappa}\wh W\|^p_{L^2(\O\X Y)}\int_{\O\X Y}(\Gu_0+{1\over \kappa}\wh W)\cdot (w_0+{1\over \kappa}W_1)\,dydx\\
		=\int_{\O\X Y}(f+\1_{\O\X Y_1}\wh \theta)\cdot (w_0+{1\over \kappa}W_1)\,dydx,\quad\forall\,(w_0,W_1)\in\D_0.
	\end{multline*}
%	where $\wh W\in L^2(\O;\GH^1_{per,0}(Y)^N)$ is the unique solution of the following cell problem \eqref{5171}.
%	\begin{equation}\label{5171}
%		\begin{aligned}
%			&\text{Find $\wh W\in L^2(\O;\GH^1_{per,0}(Y)^N)$ such that for all $w_1\in L^2(\O;\GH^1_{per,0}(Y)^N)$}\\
%			&\hskip 10mm \int_{\O\times Y} A(y)e_y(\wh W): e_y(w_1)\,dx\,dy
%			+\a\Big\|\Gu_0+\frac1\kappa\wh W\Big\|^p_{L^2(\O\times Y)}
%			\int_{\O\times Y}\Big(\Gu_0+\frac1\kappa\wh W\Big)\cdot \frac1\kappa w_1\,dx\,dy\\
%			&\hskip 100mm=\frac1\kappa\int_{\O\times Y}\big(f+\wt{\wh \Theta}\big)\cdot w_1\,dx\,dy.
%		\end{aligned}
%	\end{equation}
\end{theorem}

%The expression of $\GT_\kappa$ is direct consequence of \eqref{EX01}, 
%So, proceeding as in Theorem \ref{Th1}, we also have $\Gu_0\in\GU$ is also the minimizer for the homogenized energy given by
%$$\Gm=\GT^{hom}(\Gu_0)=\min_{w_0\in\GU}\GT^{hom}(w_0),$$
%with
%$$ \GT^{hom}(w_0)={1\over 2}\int_\O A^{hom}e(w_0):e(w_0)\,dx+{\a\over p+2}\|w_0\|^{p+2}_{L^2(\O)}-\int_\O \left(f+|Y_1|\Mc_{Y_1}(\wh \Theta)\right)\cdot w_0\,dx,\quad\forall\,w_0\in\GU.$$
%Below, we present 
%%%%%%%%%%%%%%%%%%%%%%%%%%%%%%%%%%%%%%%%%%%%%%%%%%%%%%%%%%%%%%%%%%%%%%%%%%%%%%%%%%%%%%%%%%%%%%%%%%%%%%%%%%%%%%%%%%%%%%%%%%%%%%%%%%%%%%%%%%%%%%%%%%%%%%%%%%%%%%%%%%%%%%
	% Part done by Abu.
	%\section{Unfolded limit optimal control problem}
\section{Asymptotic behavior of the $\e,\d$-OCP}\label{S06}
%	The expression of the $L^2$-cost functional is given by
%	\begin{align*}
%		\GJ_{1,\ed}(\theta_\e)=\frac{1}{2}\int_{\Omega}|{u}_\varepsilon(\theta_\varepsilon
%		)-u_0|^2\,dx+{\gamma\over 2}\int_{\O^1_\e}|\theta_\e|^2\,dx,\quad \forall\,\theta_\e\in L^2(\O^1_\e)^N,
%	\end{align*}
%	and $\Theta_\e\in L^2(\O^1_\e)^N$ is the optimal control with $\wo u_\e\in\GU$ is unique weak solution of \eqref{MainW01} corresponding to $\Theta_\e$.
%	
%	Observe that $\GJ_{1,\ed}(\Theta_\e)\leq \GJ_{1,\ed}(0)$. Then, from the fact that $u_\e(0)=\Gu_\e(0)+\Gw_\e(0)$ along with the estimates \eqref{E01}, we have
%	\begin{equation}\label{UBOC}
%		\begin{aligned}
%			\|\wo u_\e-u_0\|^2_{L^2(\O)}+\gamma\|\Theta_\e\|^2_{L^2(\O^1_\e)}
%			&\leq \|u_\e(0)-u_0\|^2_{L^2(\O)}\\
%			&\leq \|\Gu_\e(0)\|^2_{L^2(\O)}+\|\Gw_\e(0)\|^2_{L^2(\O)}+\|u_0\|^2_{L^2(\O)}\\
%			&\leq {C\e^2\over \d^2}.
%		\end{aligned}
%	\end{equation}
%	The constant is independent of $\e$ and $\d$.\\
%	Since $\kappa\in(0,+\infty]$, we have that there exists a $C>0$ independent of $\e$ and $\d$ such that (see also \eqref{R11})
%	\begin{equation}\label{UBOCC}
%		\|\Theta_\e\|_{L^2(\O^1_\e)}\leq {C_1\e\over \d}\leq C.
%	\end{equation}
	The main aim is to analyze the asymptotic behavior of the $\varepsilon$-$\d$ dependent optimal control problem as $(\e,\d)\to(0,0)$ satisfying \eqref{R1}, for which first we need a proper characterization of the optimal control.
	
	From, Theorem \ref{Th2}, we have there exist an optimal control $\Theta_\e$ such that
	$$\Gj_\ed(\Theta_\e)=\inf_{\theta_\e\in\GV}\Gj_\ed(\theta_\e),\quad \|\Theta_\e\|_{L^2(\O)}\leq C.$$
	The constant is independent of $\e$ and $\d$.
	
	The adjoint system w.r.t the unique solution $\wo u_\ed$ corresponding to $\Theta_\e$ is given by
	\begin{equation}\label{AdjW01}
		\begin{aligned}
			&\text{Find $\wo v_\ed\in \GU$ such that for all $w\in\GU$}\\
			&\hskip 10mm \Ga_\ed(w,\wo v_\ed)+\a\|\wo u_\ed\|^p_{L^2(\O)}(w,\wo v_\ed)+\a p\|\wo u_\ed\|^{p-2}_{L^2(\O)}(\wo u_\ed,w)(\wo u_\ed,\wo v_\ed)=(\wo u_\ed-u_d,w),	
		\end{aligned}
	\end{equation}
	where $(\cdot,\cdot)$ denotes the $L^2(\O)$ inner product. In the energy form the adjoint system is given by: find $\wo v_\ed\in\GU$ such that
	\[
	\Gm^\ast_\ed=\Jc_\ed(\wo v_\ed)=\min_{v\in\GU}\mathcal J_\ed(v),
	\]
	where the quadratic functional $\mathcal J_\ed:\GU\to\R$ by
	\begin{equation}\label{AdjEnergy}
		\mathcal J_\ed(v)
		=
		\frac12\,\Ga_\ed(v,v)
		+\frac{\a}{2}\,\|\wo u_\ed\|_{L^2(\O)}^{p}\,\|v\|_{L^2(\O)}^{2}
		+\frac{\a p}{2}\,\|\wo u_\ed\|_{L^2(\O)}^{p-2}\,(\wo u_\ed,v)^2
		-\;(\wo u_\ed-u_d,\,v),
		\quad v\in\GU.
	\end{equation}
	The Euler--Lagrange condition $\frac{d}{dt}\mathcal J_\ed(\wo v_\ed+t w)\big|_{t=0}=0$
	for all $w\in\GU$ is exactly \eqref{AdjW01}.
	\begin{theorem}
		There exist a unique minimizer of the adjoint energy \eqref{AdjEnergy}. The unique minimizer $\wo v_\ed\in\GU$ satisfy the following:
		\begin{equation}\label{EsA01}
			\d\|e(\wo v_\ed)\|_{L^2(\O^1_\e)}+\|e(\wo v_\ed)\|_{L^2(\O^2_\e)}\leq C\|\wo u_\ed-u_d\|_{L^2(\O)}.
		\end{equation}
		The constant is independent of $\e$ and $\d$.
	\end{theorem}
	\begin{proof}
		Fix $\wo u_\ed\in \GU$ and since $p\ge 2$. 
		Define the bilinear form $B_\ed:\GU\times\GU\to\R$ and the linear functional $F_\ed:\GU\to\R$ by
		\begin{align*}
			B_\ed(w,v)
			&=\Ga_\ed(w,v)
			+\a\|\wo u_\ed\|_{L^2(\O)}^{p}\,(w,v)
			+\a p\,\|\wo u_\ed\|_{L^2(\O)}^{p-2}\,(\wo u_\ed,w)\,(\wo u_\ed,v),\quad
			F_\ed(w)
			=(\wo u_\ed-u_d,w).
		\end{align*}
		Then the adjoint problem \eqref{AdjW01} can be written as
		\begin{equation}\label{AdjW01_LM}
			\text{Find $\wo v_\ed\in\GU$ such that } \quad B_\ed(w,\wo v_\ed)=F_\ed(w)\quad\forall w\in\GU.
		\end{equation}
		By the continuity of $\Ga_\ed$ on $\GU\times\GU$ and the Cauchy--Schwarz inequality in $L^2(\O)$,
		there exists $C>0$ (independent of $w,v$) such that
		\begin{multline*}
			|B_\ed(w,v)|
			\le C\|w\|_{H^1(\O)}\|v\|_{H^1(\O)}\\
			+\a\|\wo u_\ed\|_{L^2(\O)}^{p}\,\|w\|_{L^2(\O)}\|v\|_{L^2(\O)}
			+\a p\,\|\wo u_\ed\|_{L^2(\O)}^{p-2}\,\|\wo u_\ed\|_{L^2(\O)}^2\,\|w\|_{L^2(\O)}\|v\|_{L^2(\O)}.
		\end{multline*}
		Using the continuous embedding $H^1(\O)^N\hookrightarrow L^2(\O)^N$, there exists $C_0>0$
		such that $\|z\|_{L^2(\O)}\le C_0\|z\|_{H^1(\O)}$ for all $z\in\GU$. Hence
		\[
		|B_\ed(w,v)|\le C_1\|w\|_{H^1(\O)}\|v\|_{H^1(\O)}\qquad\forall w,v\in\GU,
		\]
		for some constant $C_1>0$. Similarly,
		\[
		|F_\ed(w)|=|(\wo u_\ed-u_d,w)|\le \|\wo u_\ed-u_d\|_{L^2(\O)}\,\|w\|_{L^2(\O)}
		\le C_0\|\wo u_\ed-u_d\|_{L^2(\O)}\,\|w\|_{H^1(\O)},
		\]
		so $F_\ed$ is continuous on $\GU$.

		Since $p\ge 2$, $\|\wo u_\ed\|_{L^2(\O)}^{p}\ge 0$ and $\|\wo u_\ed\|_{L^2(\O)}^{p-2}\ge 0$. Thus,
		for any $v\in\GU$,
		\[
		B_\ed(v,v)
		=\Ga_\ed(v,v)+\a\|\wo u_\ed\|_{L^2(\O)}^{p}\,\|v\|_{L^2(\O)}^{2}
		+\a p\,\|\wo u_\ed\|_{L^2(\O)}^{p-2}\,(\wo u_\ed,v)^2
		\ge \Ga_\ed(v,v).
		\]
		By the coercivity assumption on $\Ga_\ed$, there exists $c>0$ such that
		\[
		\Ga_\ed(v,v)\ge c\,\|v\|_{H^1(\O)}^2\qquad\forall v\in\GU\implies B_\ed(v,v)\ge c\,\|v\|_{H^1(\O)}^2\qquad\forall v\in\GU.
		\]
		Therefore, $B_\ed$ is coercive on $\GU$.

		By the Lax--Milgram theorem, there exists a unique $\wo v_\ed\in\GU$ such that \eqref{AdjW01_LM} holds, i.e.
		$\wo v_\ed$ is the unique solution of the adjoint variational problem \eqref{AdjW01}.
%		
%		\smallskip
%		\noindent\textbf{Energy characterization.}
%		Consider the functional $\mathcal J_\ed:\GU\to\R$ defined by
%		\[
%		\mathcal J_\ed(v):=\frac12\,B_\ed(v,v)-F_\ed(v).
%		\]
%		Expanding $B_\ed$ and $F_\ed$ yields exactly \eqref{AdjEnergy}. 
%		For any $v,w\in\GU$,
%		\[
%		\frac{d}{dt}\mathcal J_\ed(v+t w)\Big|_{t=0}
%		= B_\ed(w,v)-F_\ed(w).
%		\]
%		Hence $v\in\GU$ is a critical point of $\mathcal J_\ed$ if and only if it satisfies
%		$B_\ed(w,v)=F_\ed(w)$ for all $w\in\GU$, i.e. if and only if $v$ solves \eqref{AdjW01}.
		Since $B_\ed$ is coercive, $\mathcal J_\ed$ is strictly convex and therefore has at most one minimizer.
		Consequently, the unique solution $\wo v_\ed$ given by Lax--Milgram is the unique minimizer of $\mathcal J_\ed$.
		
		Finally, taking $w=\wo v_\ed$ in \eqref{AdjW01}, we get
		$$C\d^2\|e(\wo v_\ed)\|^2_{L^2(\O^1_\e)}+\|e(\wo v_\ed)\|^2_{L^2(\O^2_\e)}\leq C\|\wo u_\ed-u_d\|_{L^2(\O)}\|\wo v_\ed\|_{L^2(\O)}.$$
		Then, with the inequality \eqref{NK01} and proceeding as in Lemma \ref{L3}, we get \eqref{EsA01}.
		This completes the proof.
	\end{proof}
%	The above system is linear in $v_\ed$ for fixed $\wo u_\ed$. So using Korn's inequality together with the facts that $\Ga_\ed$ is coercive on $\GU$, the added $L^2$-mass term is nonnegative and the rank-one term is bounded. Then, using Lax-Milagram lemma, there exist a unique solution $\wo v_\ed\in\GU$ to the problem \eqref{AdjW01}.
	
	%------------------------------------------------------------
	% Energy functional for the adjoint system (all inner products in L^2(\Omega))
	%------------------------------------------------------------
	
	% Given: fixed state u_\ed^0 (denoted \wo u_\ed) and bilinear form \Ga_\ed(\cdot,\cdot)
%	
%	\paragraph{Adjoint variational system.}
%	Let $\wo u_\ed\in L^2(\O)^N$ be fixed. The adjoint state $\wo v_\ed\in\GU$ is defined by
%	\begin{equation}\label{AdjW01}
%		\begin{aligned}
%			&\text{Find $\wo v_\ed\in \GU$ such that for all $w\in\GU$}\\
%			&\hskip 10mm \Ga_\ed(w,\wo v_\ed)
%			+\a\|\wo u_\ed\|^p_{L^2(\O)}(w,\wo v_\ed)
%			+\a p\|\wo u_\ed\|^{p-2}_{L^2(\O)}(\wo u_\ed,w)(\wo u_\ed,\wo v_\ed)
%			=(\wo u_\ed-u_d,w),
%		\end{aligned}
%	\end{equation}
%	where $(\cdot,\cdot)$ denotes the $L^2(\O)$ inner product.
	
%	\paragraph{Equivalent energy formulation.}
%	Define the quadratic functional $\mathcal J_\ed:\GU\to\R$ by
%	\begin{equation}\label{AdjEnergy}
%			\mathcal J_\ed(v)
%			=
%			\frac12\,\Ga_\ed(v,v)
%			+\frac{\a}{2}\,\|\wo u_\ed\|_{L^2(\O)}^{p}\,\|v\|_{L^2(\O)}^{2}
%			+\frac{\a p}{2}\,\|\wo u_\ed\|_{L^2(\O)}^{p-2}\,(\wo u_\ed,v)^2
%			-\;(\wo u_\ed-u_d,\,v),
%			\quad v\in\GU.
%	\end{equation}
%	Then $\wo v_\ed$ is the unique minimizer of $\mathcal J_\ed$ over $\GU$, i.e.
%	\[
%	\Jc_\ed(\wo v_\ed)=\min_{v\in\GU}\mathcal J_\ed(v),
%	\]
%	and the Euler--Lagrange condition $\frac{d}{dt}\mathcal J_\ed(\wo v_\ed+t w)\big|_{t=0}=0$
%	for all $w\in\GU$ is exactly \eqref{AdjW01}.
	
%	{\clm Existence of minimizer proof is needed.}
	
	
	Below, we give the necessary condition and converse statement of the first order optimality.		
		\begin{theorem}[First-order necessary conditions]
			\label{CC01}
			Let $\wo u_\ed\in\GU$ be the (unique) state associated with an optimal control $\Theta_\e\in\GV$ for \eqref{MainW01}.
			Let $\wo v_\ed\in\GU$ solve the adjoint problem \eqref{AdjW01}, for the fixed $\wo u_\ed$. 
			Then the optimal control satisfies
			\begin{equation}\label{OCF01}
				\Theta_\e=-\frac{1}{\gamma}\,\wo v_\ed \qquad \text{in }\Omega^1_\e.
			\end{equation}
			Conversely, if $(\wo u_\ed,\wo v_\ed,\Theta_\e)\in\GU\times\GU\times\GV$ satisfies, for all $w\in\GU$,
			\begin{equation}\label{OS01}
				\begin{aligned}
					&\Ga_\ed(\wo u_\ed,w)+\alpha\|\wo u_\ed\|_{L^2(\Omega)}^{p}\,(\wo u_\ed,w)=\ell_\e(w),\\
					&\Ga_\ed(w,\wo v_\ed)
					+\alpha\|\wo u_\ed\|_{L^2(\Omega)}^{p}\,(w,\wo v_\ed)
					+\alpha p\,\|\wo u_\ed\|_{L^2(\Omega)}^{p-2}\,(\wo u_\ed,w)\,(\wo u_\ed,\wo v_\ed)
					= \,(\wo u_\ed-u_d,w),\\
					&\Theta_\e=-\frac{1}{\gamma}\,\wo v_\ed,
				\end{aligned}
			\end{equation}
			then $\Theta_\e$ is a stationary point of the OCP.
		\end{theorem}
		\begin{proof}			
			\textbf{Step 1:} Linearization of the state map.
			
			Fix $\theta_\e\in\GV$ and denote $u_\ed=u_\ed(\theta_\e)\in \GU$ the unique state solving \eqref{MainW01}. For any $h\in\GV$, the Fréchet derivative $\dot{u}_\ed=\Sc_\ed'(\theta_\e)h\in \GU$ exists and solves the
			linearized state problem
			\begin{equation}\label{lin-state}
				\Ga_\ed(\dot{u}_\ed,w)+\alpha\|u_\ed\|_{L^2(\O)}^{p}\,(\dot{u}_\ed,w)
				+\alpha p\,\|u_\ed\|_{L^2(\O)}^{p-2}\,(u_\ed,\dot{u}_\ed)(u_\ed,w)
				=(h,w)\quad\forall w\in \GU.
			\end{equation}
			Coercivity of $\Ga_\ed$ (see \eqref{l2}) on $\GU$ using Korn's inequality and uniform ellipticity together with  boundedness of the lower-order terms yield existence and uniqueness by Lax--Milgram.
			
			\textbf{Step 2:} Adjoint and duality identity.
			
			For the fixed state $u_\ed$, let $v_\ed\in \GU$ solve the adjoint problem \eqref{AdjW01} for the fixed state $u_\ed$ and control $\theta_\e\in\GV$.
			Choose $w=\dot{u}_\ed$ in \eqref{AdjW01} and  $w=v_\ed$ in \eqref{lin-state}; since the left-hand sides coincide (due to the fact that $\Ga_\ed$ is symmetric), we obtain
			\[
			(u_\ed-u_d,\dot{u}_\ed)=(h,v_\ed).
			\]
			
			\textbf{Step 3:} Derivative of the reduced cost and optimality condition.
			
			By the chain rule,
			\[
			\Gj_\ed'(\theta_\e)[h]=(u_\ed-u_d,\dot{u}_\ed)+\gamma(\theta_\e,h)
			=(h,v_\ed)+\gamma(\theta_\e,h)
			=(\gamma\theta_\e+v_\ed,h)\quad\forall h\in\GV.
			\]
			Since $\Theta_\e$ is an unconstrained optimal control, the stationary condition $\Gj_\ed'(\Theta_\e)[h]=0$ for all $h\in\GV$ implies
			\[
			(\gamma\Theta_\e+\wo v_\ed,h)=0\quad\forall h\in\GV
			\quad\Rightarrow\quad
			\gamma\Theta_\e+\wo v_\ed=0\ \text{ in $\Omega^1_\e$}
			\]
			i.e.~\eqref{OCF01} holds.
			
			{\bf Step 4:} Stationary condition.
			
			Conversely, if $(\wo u_\ed,\wo v_\ed,\Theta_\e)$ satisfies the state equation, the adjoint equation, and
			$\Theta_\e=-\wo v_\ed/\gamma$, then for every $h\in\GV$ the computation above gives
			\[
			\Gj_\ed'(\Theta_\e)[h]=(\gamma\Theta_\e+\wo v_\ed,h)=0,
			\]
			so $\Theta_\e$ is a stationary point of the OCP. This completes the proof.
		\end{proof}
		\begin{remark}[On sufficiency]
			In general with the nonlinear term $\alpha\|u\|_{L^2(\O)}^p u$, the system \eqref{OS01} gives
			first-order necessary conditions. Since, for $\alpha>0$ the control-to-state map $\Sc_\ed:\GV\to\GU$ is nonlinear (indeed, non-affine). 
			A quick scaling argument shows this: for simplicity we take $f\equiv 0$. If $u=\Sc_\ed(\theta)$ and $t>0$, then 
			$u_t:=\Sc_\ed(t\theta)$ solves
			\[
			\Ga_\ed(u_t,w)+\alpha\|u_t\|_{L^2(\Omega)}^{p}(u_t,w)=(t\theta,w)\quad\forall w\in\GU.
			\]
			If $\Sc_\ed$ were linear, we would have $u_t=tu$, hence
			\[
			\Ga_\ed(tu,w)+\alpha\|tu\|_{L^2(\O)}^{p}(tu,w)
			=t\,\Ga_\ed(u,w)+t^{p+1}\alpha\|u\|_{L^2(\O)}^{p}(u,w),
			\]
			which equals $t\big(\Ga_\ed(u,w)+\alpha\|u\|_{L^2(\O)}^{p}(u,w)\big)$ only when $\alpha=0$ or $u\equiv 0$. 
			Thus $\Sc_\ed$ is not linear for $\alpha>0$.
			As a consequence, the reduced functional 
			\[
			\Gj_\ed(\theta)=\tfrac12\|\Sc_\ed(\theta)-u_d\|_{L^2(\Omega)}^2+\tfrac{\gamma}{2}\|\theta\|_{L^2(\Omega)}^2
			\]
			is, in general, non-convex for $\alpha>0$. Indeed, differentiating twice in a direction $h\in\GV$ yields
			\[
			\Gj_\ed''(\theta)[h,h]
			=\|\Sc_\ed'(\theta)h\|_{L^2(\Omega)}^2
			+\big(\Sc_\ed(\theta)-u_d,\;\Sc_\ed''(\theta)[h,h]\big)_{L^2(\Omega)}
			+\gamma\|h\|_{L^2(\Omega)}^2,
			\]
			and the middle term has no fixed sign in general. In the linear case $\alpha=0$ we have $\Sc_\ed''\equiv 0$, so
			\[
			\Gj_\ed''(\theta)[h,h]=\|\Sc_\ed'(\theta)h\|_{L^2(\O)}^2+\gamma\|h\|_{L^2(\O)}^2\ge \gamma\|h\|_{L^2(\O)}^2,
			\]
			hence $\Gj_\ed$ is (strictly) convex and the optimal control is unique. For $\alpha>0$, nonconvexity may lead to
			multiple stationary points and non-uniqueness of global minimizers.
		\end{remark}
%%%%%%%%%%%%%%%%%%%%%%%%%%%%%%%%%%%%%%%%%%%%%%%%%%%%%%%%%%%%%%%%%%%%%%%%%%%%%%%%%%%%%%%%%%%%%%%%%%%%%%%%%%%%%%%%%%%%%%%%%%%%%%%%%%%%%%%%%%%%%%%%%%%%%%%%%
%	\subsection{Characterization of $L^2$-Energy Optimal Control}
%	The following theorem gives the characterization of the optimal control.
%	Before stating the theorem, we introduce  the adjoint state corresponding to the optimal control $\Theta_\e$, denoted by $\wo{v}_\varepsilon$, which satisfies
%	\begin{equation}\label{opt-adj-e}
%		\int_\O \sigma_\ed(\wo{v}_\e):e(w)\,dx=\int_\O\big(\wo{u}_\e-u_d\big)\cdot w\,dx,\quad\forall\,w\in\GU.
%	\end{equation}
%	Below, for $\e,\d>0$, we present the first-order optimality condition for the OCP.
%	\begin{theorem}\label{CC01}
%		Let $(\bar{u}_\varepsilon,\Theta_\varepsilon)$ be the optimal solution to the problem \eqref{ocp_ep} and $\wo v_\e$ is the corresponding adjoint state. Then 
%		\begin{equation}\label{CE01}
%			\Theta_{\varepsilon}=-\frac{1}{\gamma}\wo{v}_{\varepsilon}\quad\text{     in }\Omega^1_\e .
%		\end{equation} 
%		Conversely, if a pair  $(\wo{u}_\varepsilon,\wo{v}_\varepsilon)\in \GU \times \GU $ satisfies the following system
%		\begin{equation}\label{op-sys}
%			\left.\begin{aligned}
%				\int_\O \sigma_\ed(\wo{u}_\e):e(w)\,dx&=\int_{\O} f\cdot w\,dx-\frac{1}{\gamma}\int_{\O^1_\e}\wo{v}_\e\cdot w\,dx,\\
%				\int_\O \sigma_\ed(\wo{v}_\e):e(w)\,dx&=\int_\O\big(\wo{u}_\e-u_d\big)\cdot w\,dx,
%			\end{aligned}\right\}\quad \forall\,w\in\GU
%			%	\left\{\begin{aligned}
%				% -\nabla \cdot \sigma_\ed(\wo{u}_\e (x))&=f(x)-\frac{1}{\gamma}\wo{v}_\varepsilon \G1_\varepsilon,\quad &&\text{$x\in\O$},\\
%				%		-\nabla \cdot \sigma_\ed(\wo{v}_\varepsilon(x))&=\wo{u}_\varepsilon(x)-u_d(x) \quad &&\text{$x\in\O$},\\
%				%	\wo{u}_\e(x)=\wo{v}_\e(x)&=0,\quad &&\text{$x\in\Gamma_0$},\\
%				%\sigma_\ed(\wo{u}_\varepsilon(x))\Gn(x)&=0,\quad&&\text{$x\in \Gamma_1$},\\
%				%\sigma_\ed(\wo{v}_\varepsilon(x))\Gn(x)&=0,\quad&&\text{$x\in \Gamma_1$}.
%				%	\end{aligned}\right.
%		\end{equation}
%		Then $\left(\wo{u}_\varepsilon,-\frac{1}{\gamma}\wo{v}_\varepsilon \right)\in \GU\times L^2(\Omega^1_\e)^N$ is the optimal solution to the problem \eqref{ocp_ep}.
%	\end{theorem}
%	\begin{proof}
%		Let $(\bar{u}_\varepsilon,\Theta_\varepsilon)$ be the optimal solution to the problem \eqref{ocp_ep}. Since the constrained PDE \eqref{sta-e} is linear,  we have 
%		%$u_{\varepsilon}(\bar{\theta}_\varepsilon+\lambda\theta)$ as 
%		$$\wo{u}_{\e,\lambda}(\Theta_\e)=\wo{u}_{\varepsilon}(\Theta_\varepsilon+\lambda\theta_\e)=\bar{u}_{\varepsilon}({\Theta}_\varepsilon)+\lambda w_{\theta_\e},\quad\text{for all $\lambda>0$ and $\theta_\e\in L^2(\O^1_\e)^N$},$$
%		where $w_{\theta_\e}$ satisfies
%		\begin{equation}\label{wtheta}
%			\int_\O \sigma_\ed\big(w_{\theta_\e}\big):e(w)\,dx=\int_{\O^1_\e}\theta_\e\cdot w\,dx,\quad \forall\,w\in \GU.
%			%	\left\{\begin{aligned}\label{wtheta}
%				%		-\nabla \cdot \sigma_\ed(w_{\theta_\e} (x))&= \theta_\e\G1_\e,\quad &&\text{$x\in\O$},\\
%				%		w_{\theta_\e}(x)&=0,\quad &&\text{$x\in \Gamma_0$},\\
%				%	\sigma_\ed(w_{\theta_\e}(x))\Gn(x)&=0,\quad&&\text{$x\in \Gamma_1$},
%				%	\end{aligned}\right.
%		\end{equation}
%		Now, we evaluate the directional derivative of the cost functional $\GJ_{1,\ed}$ in the direction of $\theta_\e$ at ${\Theta}_\varepsilon$. 
%		Evaluating the limit of $$\frac{1}{\lambda}(\GJ_{1,\ed}({\Theta_\varepsilon}+\lambda \theta_\varepsilon)-\GJ_{1,\ed}({\Theta}_{\varepsilon}))$$ as $\lambda \to 0$ and denoting the limit by $\GJ_{1,\ed}^\prime({\Theta}_{\varepsilon})\theta_\e$, we get
%		\begin{align*}
%			\GJ'_{1,\ed}({\Theta}_\varepsilon) \theta_\e=\int_{\Omega}(\bar{u}_\varepsilon-u_d)\cdot w_{\theta_\e}\, dx+\gamma \int_{\Omega^1_\varepsilon} {\Theta}_{\varepsilon}\cdot \theta_\varepsilon\, dx.
%		\end{align*}
%		As $(\bar{u}_{\varepsilon}, {\Theta}_{\varepsilon})$ is the optimal solution, we have
%		$\GJ'_{1,\ed}({\Theta}_{\varepsilon})\theta_\varepsilon=0,~~\text{for~all} ~~\theta_\varepsilon \in L^{2}(\Omega^1_\varepsilon)^N$, which give
%		\begin{align}\label{beptheta}
%			\int_{\Omega} (\bar{u}_\varepsilon-u_d)\cdot w_{\theta_\varepsilon}\, dx=-\gamma \int_{\Omega_\varepsilon^1} {\Theta}_{\varepsilon}\cdot  \theta_{\varepsilon}\, dx.
%		\end{align}
%		Using $w_{\theta_\varepsilon}$ and $\wo{v}_\e$ as  test functions in \eqref{opt-adj-e} and \eqref{wtheta} respectively, we obtain
%		\begin{align}\label{bepv}
%			\int_{\Omega} (\bar{u}_\varepsilon-u_d)\cdot w_{\theta_\varepsilon}\, dx=\int_{\Omega_\varepsilon^1} \bar{v}_{\varepsilon} \cdot \theta_{\varepsilon}\, dx.
%		\end{align} 
%		Hence from \eqref{beptheta} and \eqref{bepv}, we have 
%		\begin{align*}
%			\int_{\Omega_\varepsilon^1}\bar{v}_\varepsilon\cdot \theta_\varepsilon=-\gamma\int_{\Omega^1_\varepsilon}{\Theta}_\varepsilon\cdot \theta_\varepsilon,\quad\forall\,\theta_\e\in L^2(\O^1_\e)^N,
%		\end{align*}
%		which imply  ${\Theta}_\varepsilon=-\frac{1}{\gamma}\bar{v}_\epsilon,$ a.e. in $\O^1_\e$.
%		
%		To see the converse part, suppose $(\wo{u}_\varepsilon,\bar{v}_\varepsilon)$ satisfy the system \eqref{op-sys}. Then for any $\theta_\varepsilon\in L^2(\Omega_\varepsilon^1)^N$, a direct computation give $$\GJ_{1,\ed}\left(-\frac{1}{\gamma}\bar{v}_\varepsilon+\theta_\e\right)-\GJ_{1,\ed}\left(-\frac{1}{\gamma}\bar{v}_\varepsilon\right)=\frac{1}{2}\int_{\Omega}|w_{\theta_\varepsilon}|^2\,dx +\frac{\gamma}{2}\int_{\Omega_\varepsilon^1}|\theta_\e|^2\,dx\geq 0.$$
%		This shows the optimality of the pair $(\wo{u}_\varepsilon,-\frac{1}{\gamma}\wo{v}_\varepsilon)$ for the problem \eqref{ocp_ep}. 
%		
%		This completes the proof.
%	\end{proof}
%%%%%%%%%%%%%%%%%%%%%%%%%%%%%%%%%%%%%%%%%%%%%%%%%%%%%%%%%%%%%%%%%%%%%%%%%%%%%%%%%%%%%%%%%%%%%%%%%%%%%%%%%%%%%%%%%%%%%%%%%%%%%%%%%%%%%%%%%%%%%%%%%%%%%
	From the above theorem, it is sufficient to analyze the asymptotic behavior of the adjoint state $\wo v_\ed$ to get the strong convergence of the optimal control. So, below, we first present the asymptotic behavior of the adjoint state,  followed by the derivation of two-scale adjoint problem, finally we end with the limit optimal control problem. 
	
	Let $\wo v_\ed$ be the unique adjoint state of the problem \eqref{AdjW01}. Then, using Lemma \ref{L2}, there exist $\wo\Gv_\ed\in \GU$ satisfying  \eqref{46}$_{1,3}$ and $\wo\Gz_\ed\in \GU$ satsifying \eqref{46}$_{2,4,5}$ such that
	$$\wo v_\ed=\wo\Gv_\ed+\wo\Gz_\ed,\quad \text{in $\O$}.$$
%	 we obtain $\Gv_\e={\bm\Pc}_\e\left({\wo v_\e}_{|\O^2_\e}\right)$ and $\Gz_\e=\wo v_\e-\Gv_\e$. 
	 Then, using the estimates \eqref{46}--\eqref{E01}, \eqref{UBOC}--\eqref{UBOCC} and \eqref{EsA01}, we obtain
	\begin{equation}\label{61}
		\|\wo \Gv_\ed\|_{H^1(\O)}+{\d\over \e}\|\wo\Gz_\ed\|_{L^2(\O^1_\e)}+\d\|\nabla\wo\Gz_\ed\|_{L^2(\O^1_\e)}\leq C\|\wo u_\ed-u_d\|_{L^2(\O)}\leq C.
	\end{equation}
	The constant is independent of $\e$ and $\d$.
	\begin{lemma}[Two-scale adjoint system]\label{L08}
		Let $(\bar{u}_\ed,\Theta_\e)\in\GU\X \GV$ be the optimal pair to the OCP \eqref{OCPM01} and $\wo v_\ed\in\GU$ be the adjoint state. Then, there exist $(\Gv_0,\wh Z,\wh V)\in\D$ such that
		\begin{equation}\label{Con011}
			\begin{aligned}
%				\wo\Gv_\ed &\rightharpoonup \Gv_0&&\text{weakly in $H^1(\O)^N$},\\
				\Te(\wo\Gv_\ed)&\to \Gv_0, &&\text{strongly in $L^2(\O\X Y)^N$},\\
				\Te(\nabla \wo\Gv_\ed) &\rightharpoonup \nabla \Gv_0+\nabla_y\wh V, &&\text{weakly in $L^2(\O\X Y)^{N\X N}$},\\
				 \Te(\wo\Gz_\ed) &\rightharpoonup{1\over \kappa}\wh Z,&& \text{weakly in $L^2(\O\X Y)^N$},\\
				 \d\Te(\nabla \wo\Gz_\ed) &\rightharpoonup \nabla_y\wh Z,&&\text{weakly in $L^2(\O\X Y)^{N\X N}$}.
			\end{aligned}
		\end{equation}
		Moreover, we have
		\begin{equation}\label{ECA01}
			\lim_{(\e,\d)\to(0,0)}\Jc_\ed(\wo v_\ed)=\Jc_\kappa(\Gv_0,\wh Z,\wh V)=\min_{(w_0, W_1,W_2)\in\D}\Jc(w_0,W_1,W_2)=\Gm^\ast_\kappa,
		\end{equation}
		where the two-scale adjoint energy is given by
		\begin{multline*}
			\Jc(\Gv_0,\wh Z,\wh V)={1\over 2}\left[\Ga_1(\wh Z)+\Ga_2(\Gv_0,\wh V)\right]+{\a\over 2}\left[\|\Gu_\kappa\|^p_{L^2(\O\X Y)}\|\Gv_\kappa\|^2_{L^2(\O\X Y)}+p\|\Gu_\kappa\|^{p-2}_{L^2(\O\X Y)}(\Gu_\kappa,\Gv_\kappa)^2\right]\\
			-(\Gu_\kappa-u_d,\Gv_\kappa),
		\end{multline*}
		with
		$$\Gu_\kappa=\Gu_0+{1\over \kappa}\wh W,\quad \Gv_\kappa=\Gv_0+{1\over\kappa}\wh Z,\quad\text{in $\O\X Y$}.$$
%		$\bullet$ When $\kappa=+\infty$, the duplet $(\Gu_0,\Gv_0)\in \GU\X \GU$ satisfies the following optimality system
%		\begin{align}\label{lim-ops-case 3}
%			\begin{aligned}
%				&\int_{\O} A^{hom}e(\Gu_0):e(w_0)\,dx
%				=\int_{\O} f\cdot w_0\,dx-\frac{|Y_1|}{\gamma}\int_{\O}\Gv_0\cdot w_0\,dx,\\ &
%				\int_{\O} A^{hom}e(\Gv_0):e(w_0)\,dx
%				=\int_{\O} (\Gu_0-u_0)\cdot  w_0\,dx
%			\end{aligned},
%		\end{align}
%		$\forall(w_0,w_2)\in \D_0$.\\
		Furthermore, we have the strong convergence for the microscopic optimal control
		\begin{equation}\label{OCC03}
				\Te(\Theta_\e) \to -{1\over \gamma}\left(\Gv_0+{1\over \kappa}\wh Z\right)=\wh \Theta,\quad\text{strongly in $L^2(\O\X Y_1)^N$}.
		\end{equation}
	\end{lemma}
	%{\clm \Large Give the proof that $\|\theta_\e\|_{L^2(\O^1_\e)}\leq K$.}
	\begin{proof}
		The convergences \eqref{Con011} are obtained using the estimate \eqref{61} as in Lemma \ref{L04}.
		
		Proceeding as in Theorem \ref{Th3}, we get the existence of unique minimizer of $\Jc_\kappa$ in $\D$, i.e.
		$$\min_{(w_0, W_1,W_2)\in\D}\Jc(w_0,W_1,W_2)=\Gm^\ast_\kappa\in(-\infty,0],$$
		then proceeding as in Theorem \ref{ENH01}, we obtain
		$$\lim_{(\e,\d)\to(0,0)}\Jc_\ed(\wo v_\ed)=\Jc_\kappa(\Gv_0,\wh Z,\wh V)=\Gm^\ast_\kappa,$$
		using the strong convergence \eqref{SCM}, weak convergences \eqref{Con011} and the recovery sequence $\{w_\ed\}_{\e,\d}\subset \GU$ from Step 2 of Theorem \ref{ENH01}. This gives \eqref{ECA01}.
		
		Finally, proceeding as in Corollary \ref{Cor1} and the total energy convergence \eqref{ECA01} combined with the convergences \eqref{SCM}, \eqref{Con011} give
		$$
		\lim_{(\e,\d)\to(0,0)}\left[\|\Te(\wo u_\ed)\|^p_{L^2(\O\X Y)}\|\Te(\wo v_\ed)\|^2_{L^2(\O\X Y)}\right]=\|\Gu_\kappa\|^p_{L^2(\O\X Y)}\|\Gv_\kappa\|^2_{L^2(\O\X Y)},$$
		which combined with strong convergence \eqref{SCM} and weak convergences \eqref{Con011}$_{1,3}$ imply the strong convergence
		\begin{equation}\label{SCM+}
			\Te\big(\wo v_\ed\big) \to \Gv_\kappa=\Gv_0+{1\over \kappa}\wh Z,\quad\text{strongly in $L^2(\O\X Y)^{N}$}.
		\end{equation}
		Now, we give the unfolded limit of the optimal control. Using Theorem \ref{CC01}, we have 
		$$\Theta_\e=\begin{aligned}
			%    &-{1\over \gamma}\Gv_\e,\quad&&\text{a.e. in $\O^2_\e$},\\
			&-{1\over \gamma}\wo v_\ed\quad &&\text{a.e. in $\O^1_\e$}.
		\end{aligned}
		$$
		So, using the convergences \eqref{OCC01} and \eqref{SCM+}, we get the strong convergence \eqref{OCC03} for $\kappa\in(0,+\infty]$.
%		\begin{equation}\label{OCC02}
%			\begin{aligned}
%				%   \Te(\theta_\e) &\rightharpoonup -{1\over \gamma}\Gv_0,\quad &&\text{weakly in $L^2(\O\X Y_2)^N$},\\
%				\Te(\Theta_\e) &\to -{1\over \gamma}\left(\Gv_0+{1\over \kappa}\wh Z\right)=\wh \Theta,\quad&&\text{strongly in $L^2(\O\X Y_1)^N$}.
%			\end{aligned}
%		\end{equation}
%		 In addition, when $\kappa=+\infty$, using cell problems in Subsection \ref{SS44}, we get $\wh Z=0$ a.e. in $Y_1$. 
		 This completes the proof.
%		 
%		Observe that $\bar{u}_\ed=\wo\Gu_\ed+\wo\Gw_\ed$, which implies that, using the convergences \eqref{Con01}, we get for $\kappa\in (0,+\infty]$
%		$$\begin{aligned}
%			\Te(\wo u_\e) &\rightharpoonup \Gu_0,\quad &&\text{weakly in $L^2(\O; H^1(Y_2))^N$},\\
%			\Te(\wo u_\e) &\rightharpoonup \Gu_0+{1\over \kappa}\wh W,\quad&&\text{weakly in $L^2(\O; H^1(Y_1))^N$}.
%		\end{aligned}$$
%		Hence, using the above convergences and following as in Subsection \ref{SS44} with Lemma \ref{L5}, we get \eqref{lim-ops-case 1}.
	\end{proof}
	The unique minimizer $(\Gv_0,\wh Z,\wh V)\in\D$ of $\Jc_\kappa$ satisfy the following two-scale system (in variational form)
	\begin{multline*}
		\int_{\O\X Y_2} A(y)\big(e(\Gv_0)+e_y(\wh V)\big):\big(e(w_0)+e_y(W_2)\big)\,dx dy\\
		+\int_{\O\X Y}A(y) e_y(\wh Z):e_y(W_1)\,dxdy+\alpha\|\Gu_\kappa\|^p_{L^2(\O\X Y)}\int_{L^2(\O\X Y)}\Gv_\kappa\cdot\Gw_\kappa\,dydx\\
		+\a p\|\Gu_\kappa\|^{p-2}_{L^2(\O\X Y)}\int_{\O\X Y}\Gu_\kappa\cdot\Gw_\kappa\,dydx\int_{\O\X Y}\Gu_\kappa\cdot\Gv_\kappa\,dydx\\
		=\int_{\O\X Y}(\Gu_\kappa-u_d)\cdot\Gw_\kappa\,dydx,
	\end{multline*}
	for all $(w_0,W_1,W_2)\in\D$ with $\Gw_\kappa=w_0+{1\over \kappa}W_1$.
	
	Then, proceeding as in Subsection \ref{SS44}, we can express $\wh V$ in terms of $\Gv_0$ and correctors $\chi^2_{ij}$ to get the two-scale limit adjoint system for $(\Gv_0,\wh Z)\in\D_0$
	\begin{equation}\label{LAS01}
		\begin{aligned}
		\int_{\O} A^{hom}e(\Gv_0)&:e(w_0)\,dx
		+\int_{\O\X Y}A(y) e_y(\wh Z):e_y(W_1)\,dxdy\\
		&+\a p\|\Gu_\kappa\|^{p-2}_{L^2(\O\X Y)}\int_{\O\X Y}\Gu_\kappa\cdot\Gw_\kappa\,dydx\int_{\O\X Y}\Gu_\kappa\cdot\Gv_\kappa\,dydx\\
		&\hskip 10mm +\alpha\|\Gu_\kappa\|^p_{L^2(\O\X Y)}\int_{L^2(\O\X Y)}\Gv_\kappa\cdot\Gw_\kappa\,dydx=\int_{\O\X Y}(\Gu_\kappa-u_d)\cdot\Gw_\kappa\,dydx.
		\end{aligned}
	\end{equation}
	Now, we are in position to give the limit OCP.
	\subsection{Limit optimal control problem}
	 The two-scale limit OCP is given by
	\begin{equation}\label{LOCP}
		\Gi_\kappa=\inf_{\wh \theta\in L^2(\O\X Y_1)^N}\left\{\Gj_\kappa(\wh\theta)=\GJ_\kappa(v_0,V_1,\wh\theta)\,|\,\text{subjected to $(v_0,V_1)\in\D_0$ satisfying \eqref{TS01}}\right\},
	\end{equation}
	where
	$$\GJ_\kappa(v_0,V_1,\wh\theta)= \left\|v_0+{1\over \kappa}V_1-u_d\right\|_{L^2(\O\X Y)}^2+\frac{\gamma}{2}\left\|\wh\theta\right\|^2_{L^2(\O\X Y_1)} ,$$
	and $(v_0,V_1)\in\D_0$ is the unique solution of
	\begin{multline}\label{TS01}
		\int_{\O}A^{hom}e(v_0):e(w_0)\,dx+\int_{\O\X Y}A(y) e_y(V_1):e_y(W_1)\,dxdy\\
		+\a\|v_0+{1\over \kappa}V_1\|^p_{L^2(\O\X Y)}\int_{\O\X Y}(v_0+{1\over \kappa}V_1)\cdot (w_0+{1\over \kappa}W_1)\,dydx\\
		=\int_{\O\X Y}(f+\1_{\O\X Y_1}\wh \theta)\cdot (w_0+{1\over \kappa}W_1)\,dydx,\quad\forall\,(w_0,W_1)\in\D_0,
	\end{multline}
	 corresponding to $\wh\theta\in L^2(\O\X Y_1)^N$.
	 

	\begin{theorem}
		For each $\kappa\in(0,+\infty]$, the OCP \eqref{LOCP} admits a solution. The limit control to state map $\Sc_\kappa:L^2(\O\X Y_1)^N\to \D_0$ is Lipschitz continuous.
	\end{theorem}
	\begin{proof}
		Observe that the unique solution $(v_0,V_1)\in\D_0$ of \eqref{TS01} corresponding to $\wh \theta$ satisfy
		$$\|v_0\|_{H^1(\O)}+\|V_1\|_{L^2(\O;H^1(Y))}\leq C(\|f\|_{L^2(\O)}+\|\wh \theta\|_{L^2(\O\X Y_1)}),$$
		due to the inequality \eqref{LES01}. 
		
		As
		$$\Gi_\kappa=\inf_{\wh \theta\in L^2(\O\X Y_1)^N}\Gj_\kappa(\wh \theta),$$
		we have $\Gi_\kappa\in[0,+\infty)$ since $\Gi_\kappa\leq \Gj_\kappa(0)<+\infty$. So, proceeding as in the proof of Theorem \ref{Th2}, we get the existence of optimal control for the limit OCP \eqref{LOCP}.
		
		Finally, using the properties of the homogenized stiffness tensor $A^{hom}$ from Lemma \ref{HT01} and proceeding as in Theorem \ref{Th03}, we get the limit control to state mapping is Lipschitz continuous.
	\end{proof}	
	
	Observe that the adjoint system w.r.t the unique solution $(\Gu_0,\wh W)\in\D_0$ of \eqref{TS01} corresponding to $\wh\Theta$ is given by \eqref{LAS01}. The adjoint system \eqref{LAS01} has unique solution $(\Gv_0,\wh Z)\in\D_0$. Now, we prove that $\wh \Theta\in L^2(\O\X Y_1)^N$ is an optimal control of the limit OCP \eqref{LOCP}.
	\begin{theorem}
		The limit of the microscopic optimal control $\wh\Theta\in L^2(\O\X Y_1)^N$ is an optimal control of the limit OCP \eqref{LOCP}, i.e.
		$$\Gj_\kappa(\wh\Theta)=\Gi_\kappa.$$
	\end{theorem} 
	\begin{proof}
		To prove the main result of this theorem, we use a form of $\Gamma$-convergence, to derive the following energy convergence
		\begin{equation}\label{LM}
			\Gj_\kappa(\wh \Theta)=\lim_{(\e,\d)\to(0,0)}\Gj_\ed(\Theta_\e)=\lim_{(\e,\d)\to(0,0)}\Gi_\ed=\Gi_\kappa.
		\end{equation}
		First, due to the strong convergences \eqref{SCM} and \eqref{SCM+}, we get
		\begin{equation}\label{LI}
			\Gi_\kappa\leq \Gj_\kappa(\wh\Theta)=\liminf_{(\e,\d)\to(0,0)}\Gj_\ed(\Theta_\e)=\liminf_{(\e,\d)\to(0,0)}\Gi_\ed.
		\end{equation}
		Now, we prove
		\begin{equation}\label{LS}
			\lim_{(\e,\d)\to(0,0)}\Gi_\ed\leq \Gj_\kappa(\wh \theta),\quad \forall\,\wh \theta\in L^2(\O\X Y_1)^N.
		\end{equation}
		Let $\wh \theta\in L^2(\O\X Y_1)^N$, then there exist a unique state $(v_0,V_1)\in\D_0$, which is the unique solution of \eqref{TS01}, corresponding to $\wh\theta$. 
		
		Let us set
		\begin{equation*}
			\begin{aligned}
				\theta_\e(x)&=\wh\theta\left(x,\left\{x\over \e\right\}\right) 
%				v_\ed(x)=v_0(x)+{\e\over \d}V_1\left(x,\left\{x\over \e\right\}\right)
			\end{aligned},\quad x\in \O^1_\e.
		\end{equation*}
		Observe that $\theta_\e\in\GV$ and using the properties of the unfolding operator, we have
		\begin{equation*}
			\begin{aligned}
				\Te(\theta_\ed)&\to \wh\theta,\quad &&\text{strongly in $L^2(\O\X Y_1)^N$}.
%				\Te(v_\ed) &\to v_0+{1\over \kappa}V_1,\quad &&\text{strongly in $L^2(\O\X Y)^N$}.
			\end{aligned}
		\end{equation*} 
		Observe that, there exist a unique solution $v_\ed\in\GU$ of \eqref{MainW01} corresponding to $\theta_\ed\in\GV$. Proceeding as Theorem \ref{ENH01} and Corollary \ref{Cor1}, we get: there exist
		$(v_0',V_1')\in\D_0$ such that 
		\begin{equation}\label{strong-to-some-limit}
			T_\e(v_{\e\d})\to v_0'+\frac1\kappa V_1'
			\quad\text{strongly in }L^2(\Omega\times Y)^N.
		\end{equation}
				Using $w_\ed$ from Step 2 of Theorem \ref{ENH01} as test function in the state equation \eqref{MainW01} satisfied by $v_{\e\d}$, corresponding to $\theta_\e$, give
		\begin{equation}\label{weak-eps}
			\Ga_{\e\d}(v_{\e\d},w_\ed)
			+\alpha\|v_{\e\d}\|_{L^2(\Omega)}^{p}\int_\Omega v_{\e\d}\cdot w_\ed\,dx
			=\int_\Omega f\cdot w_\ed\,dx+\int_{\Omega_\e^1}\theta_\e\cdot w_\ed\,dx.
		\end{equation}
				We now pass to the limit in each term, using the unfolding operator and the strong convergences \eqref{Test01} and \eqref{strong-to-some-limit}, to get
			\begin{multline*}	\int_{\O}A^{hom}e(v'_0):e(w_0)\,dx+\int_{\O\X Y}A(y) e_y(V'_1):e_y(W_1)\,dxdy\\
			+\a\|v'_0+{1\over \kappa}V'_1\|^p_{L^2(\O\X Y)}\int_{\O\X Y}(v'_0+{1\over \kappa}V'_1)\cdot (w_0+{1\over \kappa}W_1)\,dydx\\
			=\int_{\O\X Y}(f+\1_{\O\X Y_1}\wh \theta)\cdot (w_0+{1\over \kappa}W_1)\,dydx.
		\end{multline*}
		To obtained the above, we have proceeded as in Section \eqref{SS44}, specifically Theorem \ref{56}. So, we got the exact limit variational formulation \eqref{TS01}
		with the pair $(v_0',V_1')$ and the control $\widehat\theta$. Therefore $(v_0',V_1')$ is a solution of \eqref{TS01}
		associated with $\widehat\theta$. By uniqueness of the limit state (for the given $\widehat\theta$),
		we conclude
		\[
		(v_0',V_1')=(v_0,V_1).
		\]
% We prove that 
%		\begin{equation*}
%			\begin{aligned}
%%				\Te(\theta_\ed)&\to \wh\theta,\quad &&\text{strongly in $L^2(\O\X Y_1)^N$},\\
%				\Te(v_\ed) &\to v_0+{1\over \kappa}V_1,\quad &&\text{strongly in $L^2(\O\X Y)^N$}.
%			\end{aligned}
%		\end{equation*}
		The above convergence \eqref{strong-to-some-limit} together with the fact that $\Gi_\ed\leq \Gj_\ed(\theta_\e)=\GJ_\e(v_\ed,\theta_\e)$ gives
		$$\limsup_{(\e,\d)\to(0,0)}\Gi_\ed\leq \limsup_{(\e,\d)\to(0,0)}\Gj_\e(\theta_\e)=\limsup_{(\e,\d)\to(0,0)}\GJ_\e(v_\ed,\theta_\e)=\GJ_\kappa(v_0,V_1,\wh\theta)=\Gj_\kappa(\wh \theta).$$
		Since, $\wh \theta\in L^2(\O\X Y_1)^N$ is arbitary, we get \eqref{LS}. Now, taking an optimal control of the OCP \eqref{LOCP} in place of $\wh \theta$ in \eqref{LS}, gives (together with \eqref{LI})
		$$\Gi_\kappa\leq \Gj_\kappa(\wh\Theta)=\liminf_{(\e,\d)\to(0,0)}\Gi_\ed=\limsup_{(\e,\d)\to(0,0)}\Gi_\ed\leq \Gi_\kappa.$$
		So, we obtain \eqref{LM}. This completes the proof.
	\end{proof}
%%%%%%%%%%%%%%%%%%%%%%%%%%%%%%%%%%%%%%%%%%%%%%%%%%%%%%%%%%%%%%%%%%%%%%%%%%%%%%%%%%%%%%%%%%%%%%%%%%%%%%%%%%%%%%%%%%%%%%%%%%%%%%%%%%%%%%%%%%%%%%%%%%%%%%%%
%	\begin{proposition}[Strong unfolded convergence and identification of the limit state]\label{Prop:strong-state}
%		Let $\widehat\theta\in L^2(\Omega\times Y_1)^N$ and define the realized control
%		\[
%		\theta_\e(x):=\widehat\theta\!\left(x,\left\{\frac{x}{\e}\right\}\right)\mathbf 1_{\Omega_\e^1}(x),
%		\qquad x\in\Omega,
%		\]
%		(with $\theta_\e$ extended by $0$ on $\Omega\setminus\Omega_\e^1$). Let $v_{\e\d}\in\mathcal U$ be the unique weak solution of \eqref{MainW01}
%		corresponding to $\theta_\e$.
%		Then, as $(\e,\d)\to(0,0)$ with $\d/\e\to\kappa\in(0,+\infty]$, we have
%		\begin{equation}\label{strong-state}
%			T_\e(\overline v_{\e\d})\to v_0+\frac1\kappa V_1
%			\quad\text{strongly in }L^2(\Omega\times Y)^N,
%		\end{equation}
%		where $(v_0,V_1)\in\mathcal D_0$ is the (unique) solution of the limit state problem \eqref{TS01}
%		associated with the control $\widehat\theta$ (with the convention $1/\kappa=0$ when $\kappa=+\infty$).
%	\end{proposition}
%	
%	\begin{proof}
%		\textbf{Step 1: compactness and strong convergence to some limit.}
%		By the a priori estimate for the state equation (uniform in the index of the sequence), the family
%		$\{v_{\e\d}\}$ is bounded in $H^1(\Omega)^N$, hence (up to a subsequence) there exist
%		$(v_0',V_1')\in\mathcal D_0$ such that
%		\[
%		T_\e(\overline v_{\e\d})\rightharpoonup v_0'+\frac1\kappa V_1'
%		\quad\text{weakly in }L^2(\Omega\times Y)^N.
%		\]
%		Moreover, proceeding exactly as in Corollary~\ref{Cor1} (using the corresponding energy convergence
%		for the sequence $v_{\e\d}$), we obtain the convergence of norms
%		\[
%		\|T_\e(\overline v_{\e\d})\|_{L^2(\Omega\times Y)} \to 
%		\left\|v_0'+\frac1\kappa V_1'\right\|_{L^2(\Omega\times Y)}.
%		\]
%		Since $L^2(\Omega\times Y)^N$ is uniformly convex, the combination of weak convergence and convergence
%		of norms yields
%		\begin{equation}\label{strong-to-some-limit}
%			T_\e(\overline v_{\e\d})\to v_0'+\frac1\kappa V_1'
%			\quad\text{strongly in }L^2(\Omega\times Y)^N.
%		\end{equation}
%		
%		\textbf{Step 2: identification of the limit by passing to the limit equation.}
%		Fix an arbitrary test pair $(\varphi,\Phi)\in\mathcal D_0$. Let $w_\e\in\mathcal U$ be the usual oscillating
%		test function associated with $(\varphi,\Phi)$ (the same choice as in the derivation of the limit system),
%		so that, as $(\e,\d)\to(0,0)$,
%		\[
%		T_\e(\overline w_\e)\to \varphi+\frac1\kappa\Phi
%		\quad\text{strongly in }L^2(\Omega\times Y)^N,
%		\]
%		and $w_\e\to\varphi$ weakly in $H^1(\Omega)^N$, with the corresponding convergences for the strains
%		used in your limit passage (cf.\ \eqref{Con011}).
%		
%		Using $w_\e$ as test function in the state equation \eqref{MainW01} satisfied by $v_{\e\d}$, we have
%		\begin{equation}\label{weak-eps}
%			\mathcal A_{\e\d}(v_{\e\d},w_\e)
%			+\alpha\|v_{\e\d}\|_{L^2(\Omega)}^{p}\int_\Omega v_{\e\d}\cdot w_\e\,dx
%			=\int_\Omega f\cdot w_\e\,dx+\int_{\Omega_\e^1}\theta_\e\cdot w_\e\,dx.
%		\end{equation}
%		
%		We now pass to the limit in each term:
%		
%		\smallskip
%		\noindent\emph{(i) Control term.}
%		By the definition of $\theta_\e$ and the standard property of periodic unfolding,
%		\[
%		T_\e(\theta_\e)\to \widehat\theta
%		\quad\text{strongly in }L^2(\Omega\times Y_1)^N.
%		\]
%		Together with $T_\e(\overline w_\e)\to \varphi+\frac1\kappa\Phi$ strongly in $L^2$,
%		we obtain
%		\[
%		\int_{\Omega_\e^1}\theta_\e\cdot w_\e\,dx
%		=\int_{\Omega\times Y_1}T_\e(\theta_\e)\cdot T_\e(\overline w_\e)\,dx\,dy + o(1)
%		\to \int_{\Omega\times Y_1}\widehat\theta\cdot\Big(\varphi+\frac1\kappa\Phi\Big)\,dx\,dy.
%		\]
%		
%		\smallskip
%		\noindent\emph{(ii) Nonlinear (nonlocal) term.}
%		Using the unfolding norm identity (up to a boundary remainder) and \eqref{strong-to-some-limit}, we have
%		\[
%		\|v_{\e\d}\|_{L^2(\Omega)}=\|T_\e(\overline v_{\e\d})\|_{L^2(\Omega\times Y)}+o(1)
%		\to \left\|v_0'+\frac1\kappa V_1'\right\|_{L^2(\Omega\times Y)}.
%		\]
%		Also,
%		\[
%		\int_\Omega v_{\e\d}\cdot w_\e\,dx
%		=\int_{\Omega\times Y}T_\e(\overline v_{\e\d})\cdot T_\e(\overline w_\e)\,dx\,dy + o(1)
%		\to \int_{\Omega\times Y}\Big(v_0'+\frac1\kappa V_1'\Big)\cdot\Big(\varphi+\frac1\kappa\Phi\Big)\,dx\,dy,
%		\]
%		hence
%		\[
%		\alpha\|v_{\e\d}\|_{L^2(\Omega)}^{p}\int_\Omega v_{\e\d}\cdot w_\e\,dx
%		\to
%		\alpha\left\|v_0'+\frac1\kappa V_1'\right\|_{L^2(\Omega\times Y)}^{p}
%		\int_{\Omega\times Y}\Big(v_0'+\frac1\kappa V_1'\Big)\cdot\Big(\varphi+\frac1\kappa\Phi\Big)\,dx\,dy.
%		\]
%		
%		\smallskip
%		\noindent\emph{(iii) Elastic term.}
%		By the convergences used in \eqref{Con011} and the standard unfolding limits for the bilinear form,
%		we have
%		\[
%		\mathcal A_{\e\d}(v_{\e\d},w_\e)\to \mathcal A_2(v_0',\varphi)+\mathcal A_1(V_1',\Phi)
%		\quad\text{(in the notation of your limit problem)}.
%		\]
%		
%		\smallskip
%		\noindent\emph{(iv) Forcing term.}
%		Since $w_\e\rightharpoonup\varphi$ in $H^1(\Omega)^N$, in particular $w_\e\to\varphi$ strongly in $L^2(\Omega)^N$,
%		hence
%		\[
%		\int_\Omega f\cdot w_\e\,dx\to\int_\Omega f\cdot\varphi\,dx.
%		\]
%		
%		Passing to the limit in \eqref{weak-eps} yields exactly the limit variational formulation \eqref{TS01}
%		with the pair $(v_0',V_1')$ and the control $\widehat\theta$. Therefore $(v_0',V_1')$ is a solution of \eqref{TS01}
%		associated with $\widehat\theta$. By uniqueness of the limit state (for the given $\widehat\theta$),
%		we conclude
%		\[
%		(v_0',V_1')=(v_0,V_1).
%		\]
%		
%		\textbf{Step 3: conclusion.}
%		Combining this identification with \eqref{strong-to-some-limit} gives \eqref{strong-state}.
%		Since the limit is unique, the convergence holds for the whole sequence (not only for a subsequence).
%	\end{proof}
%%%%%%%%%%%%%%%%%%%%%%%%%%%%%%%%%%%%%%%%%%%%%%%%%%%%%%%%%%%%%%%%%%%%%%%%%%%%%%%%%%%%%%%%%%%%%%%%%%%%%%%%%%%%%%%%%%%%%%%%%%%%%%%%%%%%%%%%%%%%%%%%%%%%%%%%
	
	Finally, we end this section with the first order optimality condition for the limit OCP \eqref{LOCP}.
	\begin{theorem}
		Let $(\Gu_0,\wh W)\in\D_0$ be the unique state associated with an optimal control $\wh \Theta\in L^2(\O\X Y_1)^N$ for \eqref{TS01}. Let $(\Gv_0,\wh Z)\in\D_0$ solve the adjoint system \eqref{LAS01}, for fixed $(\Gu_0,\wh W)$. Then, the optimal control satisfies
		\begin{equation}
			\wh \Theta=-{1\over \gamma}(\Gv_0+{1\over \kappa}\wh Z),\quad \text{in $\O\X Y_1$}.
		\end{equation}
	\end{theorem}
	\begin{proof}
		The result is direct consequence of Lemma \ref{L08}.
	\end{proof}
		The limit optimal control problem for the microscopic OCP \eqref{OCPM01} is given by \eqref{LOCP}. For $\kappa=+\infty$, we have a complete scale separated state equation, adjoint and limit OCP. 

%	Let us denote by $\wt \Theta\in L^2(\O\X Y_1)^N$ be an optimal control of the OCP. Moreover, let $(\wt\Gu_0,\wh W)\in\D_0$ be the unique state aasociated to $\wt \Theta$ and $(\wt \Gv_0,\wt Z)\in\D_0$ be the respective adjoint state.
%		
%	Below, we present the main result of this paper
%	\begin{theorem}[Limit $L^2$-OCP]
%		The limit optimal control problem for the microscopic OCP \eqref{OCPM01} is given by \eqref{LOCP}.
%	\end{theorem}
%	\begin{proof}
%		Due to the strong convergences \eqref{SCM} and \eqref{SCM+}, we get
%		$$\lim_{(\e,\d)\to(0,0)}\Gj_\ed(\Theta_\e)=\Gj_\kappa(\wh\Theta).$$
%		The main result is a consequence of the above.
%	\end{proof}
%
%	\begin{remark}
%					Observe that for $\kappa\in(0,+\infty]$, we have
%					$$\Theta=\Mc_{Y_1}(\wh\Theta)=-{1\over \gamma}\Gv_0,\quad\text{a.e. in $\O$}.$$
%				\end{remark}

				% where $|Y_2|f=F$ and $|Y_1|\theta_0=\theta$,\\
				%with $\Mc_{Y_1}(\wh\W)=0$.\\										
				
				
				%\begin{remark}
				%	If $A=B$ then we have the following version of Theorem \ref{TH08}
				%	$$\GJ_{2,\kappa}
				%	
				%\end{remark}
				
				%\begin{lemma}
				%There exists a subsequence (still denoted by $\e$),   $\Gu_0,\Gv_0\in H^1(\O)^N$ and $\wh U, \wh V\in L^2(\O; H^1_{per,0}(Y))^N$ with
				%	$$\Gv_0=\Gu_0=0,\quad \text{on $\Gamma_0$},$$
				%	such that
				%\end{lemma}
%				\section{Conclusions}
%					XXXXXXXXXXXXXXX
				
				\section*{Acknowledgment}
				The authors would like to thank Prof.~Dr.~Georges Griso and Dr.~Julia Orlik for their suggestions and corrections. Abu Sufian acknowledges support from the Agencia Nacional de Investigaci\'on y Desarrollo, Chile, under the FONDECYT postdoctoral fellowship,  project no ($N^\circ ~3240018$).
%				\section*{Statements and Declarations}
%				The authors declare no conflict of interest.
				
				%%%%%%%%%%%%%%%%%%%%%%%%%%%%%%%%%%%%%%%%%%%%%%%%%%%%%%%%%%%%%%%%%%%%%%%%%%%%%%%%%%%%%%%%%%%%%%%%%%%%%%%%%%%%%%%%%%%%%%%%%%%%%%%%%%%%%%%%%%%%%%%%%%%%%%%%%%%%%%%%%%%%%%
				
				
				
				%%%%%%%%%%%%%%%%%%%%%%%%%%%%%%%%%%%%%%%%%%%%%%%%%%%%%%%%%%%%%%%%%%%%%%%%%%%%%%%%%%%%%%%%%%%%%%%%%%%%%%%%%%%%%%%%%%%%%%%%%%%%%%%%%%%%%%%%%%%%%%%%%%%%%%%%%%%%%%%%%%%%% 
				
				%%%%%%%%%%%%%%%%%%%%%%%%%%%%%%%%%%%%%%%%%%%%%%%%%%%%%%%%%%%%%%%%%%%%%%%%%%%%%%%%%%%%%%%%%%%%%%%%%%%%%%%%%%%%%%%%%%%%%%%%%%%%%%%%%%%%%%%%%%%%%%%%%%%%%%%%%%%%%%%%%%%
				% Refernces
				{\small\bibliographystyle{plain}
				\bibliography{reference}} % Your .bib file without the .bib extension
				
				%\printbibliography
				%										\bibliographystyle{ios1}
				%										\bibliography{reference}
				%									\bibliographystyle{abbrv}
				%									\bibliography{reference}	
				%%%%%%%%%%%%%%%%%%%%%%%%%%%%%%%%%%%%%%%%%%%%%%%%%%%%%%%%%%%%%%%%%%%%%%%%%%%%%%%%%%%%%%%%%%%%%%%%%%%%%%%%%%%%%%%%%%%%%%%%%%%%%%%%%%%%%%%%%%%%%%%%%%%%%%%%%%%%%%%%%%%%%%
				\begin{appendices}
					\section{Appendix}
					In this section, we prove some coercivity results for $N\geq 2$.
					\begin{theorem}
						For $(\Gu_0,\wh U)\in \GU\X L^2(\O; H^1_{per,0}(Y_2))^N$, we have that there exists a $C_0>0$ such that
						\begin{equation}\label{Ex01}
							\|\Gu_0\|^2_{H^1(\O)}+\|\wh U\|_{L^2(\O;H^1(Y_2))}\leq C_0\int_{\O\X Y_2}\big(e(\Gu_0)+e_y(\wh U)\big):\big(e(\Gu_0)+e_y(\wh U)\big)\,dxdy.
						\end{equation}
					\end{theorem}
					The above inequality is a direct consequence of Korn's inequality and the fact that $\int_{Y_2}\nabla_y\wh U\,dy=0$ since $\wh U\in H^1_{per,0}(Y_2)^N$. Using the above inequality we prove using Lax-Milgram lemma that the unfolded limit problems have a unique solution.
%					
%					We have the following  norm equivalence result when $N\geq 2$. 
%					\begin{lemma}\label{NEq1}
%						Consider the space $\S=\R^{N\X N}_\sy\X H^1_{per,0}(Y)^N$ with the semi-norm 
%						$$	\|(F,w)\|^2_\S = \sum_{i,j=1}^N  \|[F,w]_{ij}\|^2_{L^2(Y)},
%						$$ 
%						where for every $(F,w) \in \S$, we denote $[F,w]$ the symmetric matrix by 
%						$$\begin{aligned}
%							%			\text{for $N=2$},\quad 	&{\cal E}(\eta,w)_{ij}&&=
%							%			\begin{pmatrix}
%								%					\eta_1+ e_{11,y}( w)  & \eta_3 + e_{12,y}( w) \\
%								%					* & \eta_2+ e_{22,y}( w) 
%								%				\end{pmatrix},\\
%							\text{for $N\geq 2$},\quad 	&[F,w]&&= F+e_y(w).
%							%			\begin{pmatrix}
%								%					\eta_1+ e_{11,y}( w)  & \eta_4 + e_{12,y}( w) & \eta_5+e_{13,y}(w) \\
%								%					* & \eta_2+ e_{22,y}( w) &\eta_6+e_{23,y}(w) \\
%								%					* & *&  \eta_3+  e_{33,y}( w) 
%								%				\end{pmatrix}.
%						\end{aligned}
%						$$ 
%						Then, $\|\cdot\|_{\GS}$ is a norm in $\GS$ and there exists constants $c,C>0$ such that for all $(F,w) \in \S$ we have
%						\begin{equation}\label{ENE1}
%							c\left(\|F\|_2^2  + \|w\|^2_{H^1_{per,0}(Y)}\right) \leq  \|(F,w)\|^2_\S
%							\leq C \left(\|F\|_2^2 + \|w\|^2_{H^1_{per,0}(Y)}\right).
%						\end{equation} 
%					\end{lemma}
%					\begin{proof}
%						{\bf Step 1:} We prove that	for $(F,w)\in \GS$ satisfy $\|(F,w)\|_\GS=0$ implies $(F,w)=0$.\\
%						For that we define the map
%						$$\tau(y)=Fy.
%						%    \begin{pmatrix}
%							%			y_1\eta_1+y_2\eta_4+y_3\eta_5\\
%							%			y_1\eta_4+y_2\eta_2+y_3\eta_6\\
%							%			y_1\eta_5+y_2\eta_6+y_3\eta_3
%							%		\end{pmatrix}.
%						$$
%						Then, we have
%						$$[F,w]=e_y(\tau+w)\quad\text{and}\quad [0,w]=e_y(w).$$
%						So, $\|(F,w)\|_\GS=0$ implies $e_y(\tau+w)=0$. Hence, $\tau(y)+w(y)=a+B y$ is a rigid motion, with $a\in\R^N$ and $B$ is an anti-symmetric matrix in $\R^{N\X N}$.
%						Since, $w$ is periodic in $\Ge_i$, for $i\in\{1,\ldots,N\}$ directions with mean zero, i.e.
%						\begin{equation}\label{II01}
%							\begin{aligned}
%								w(y+n\Ge_i)&=w(y),\quad\text{a.e. on $Y$, $\forall n\in\Z$, $i\in \{1,\ldots,N\}$},\\
%								\Mc_{Y}(w)&=\int_{Y}w(y)\,dy=0,
%							\end{aligned}
%						\end{equation}
%						and we also have
%						$$w(y)=a+(B-F)y.
%						%    =\begin{pmatrix}
%							%			y_1\eta_1+y_2\eta_4+y_3\eta_5-a_1-b_2y_3+b_3y_2\\
%							%			y_1\eta_4+y_2\eta_2+y_3\eta_6-a_2-b_3y_1+b_1y_3\\
%							%			y_1\eta_5+y_2\eta_6+y_3\eta_3-a_3-b_1y_2+b_2y_1
%							%		\end{pmatrix}.
%						$$
%						So, applying \eqref{II01}$_1$ on $w$ give
%						$$(B-F)\Ge_i=0,\quad\forall\, i\in \{1,\ldots,N\}.$$
%						Using the fact that $\{\Ge_i\}$ are an orthonormal basis of $\R^N$ and $F$ is symmetric with $B$ being anti-symmetric give $B=F=0$. So, we have $w(y)=a$. Finally, applying \eqref{II01}$_2$ to $w$ gives $a=0$, which in-turn, we get $w=0$.\\[1mm]
%						{\bf Step 2:} We prove \eqref{ENE1}$_1$ using contradiction method.\\
%						If \eqref{II01}$_1$ is not true then for all $n\in\N^*$, there exists a $(F_n,w_n)\in \GS$ such that
%						$$\|(F_n,w_n)\|_\GS<{1\over n}\quad\text{and}\quad\|F_n\|_2+\|w\|_{H^1( Y)}=1.$$
%						Then, using the above estimate and the result from Step 1, along with the fact that $\|e_y(\cdot)\|_{L^2(Y)}$ is a norm in $H^1_{per,0}(Y)^N$, we obtain
%						$$\lim_{n\to0}\big(\|F_n\|_2+\|w_n\|_{H^1(\Yc)}\big)=0.$$
%						A contradiction.\\
%						Observe that \eqref{ENE1}$_2$ is trivial.
%						This completes the proof. 
%						%     Now, applying \eqref{II01}$_2$ give $a=0$. Hence, we obtain $(F,w)=(0,0)$.\\
%						%	$$\eta_1=\eta_2=\eta_3=0,\quad \eta_4=-\eta_4=b_3,\quad \eta_5=-\eta_5=b_2\quad\text{and}\quad\eta_6=-\eta_6=b_1,$$
%						%	which imply $\eta=0$ and $b=0$. 
%						%	For $N=2$, we have
%						%	$$B=-B^T=\begin{pmatrix}
%							%			b_1 & b_2\\
%							%			-b_2 & B_3
%							%		\end{pmatrix}.$$ Then, similarly as above we    
%					\end{proof}
				\end{appendices}
\end{document}

					
%					\newpage
%					{\clm -----------------------------------------------------------------------------------------------------------------------------------------------------}
%					\subsection{Energy convergences: Commutativity of $\d\to0$ and $\e\to0$}\label{RR2}
%					{\clm As a consequence of above energy convergence, 
%						\begin{corollary}
%							We have
%							\begin{multline}
%								\lim_{(\e,\d)\to(0,0)}\Ga_\ed(\wo u_\ed,\wo u_\ed)\\
%								=\int_{\O\X Y_1}A(y) e_y(\wh W):e_y(\wh W)\,dxdy+\int_{\O\X Y_2} A(y)\big(e(\Gu_0)+e_y(\wh U)\big):\big(e(\Gu_0)+e_y(\wh U)\big)\,dx dy.
%							\end{multline}
%					\end{corollary}}
\section{Supplementary energy convergences}
					\begin{theorem}
						We have
						$$\lim_{\d\to0}\lim_{\e\to0}\Gm_\ed=\lim_{\d\to0}\Gm_\d=\Gm,$$
						where
						$$\Gm_\d=\min_{(w_0,w_1,w_2)\in \D}\GT_\d(w_0,w_1,w_2),$$
						with
						\begin{multline*}
							\GT_\d(w_0,w_1,w_2)={1\over 2}
							\int_{\O\X Y_1}A(y)\d\big(e(w_0)+e_y(w_1)+e_y(w_2)\big):\d\big(e(w_0)+e_y(w_1)+e_y(w_2)\big)\,dxdy
							\\+{1\over 2}\int_{\O\X Y_2}A(y)\big(e(w_0)+e_y(w_1)\big):\big(e(w_0)+e_y(w_1)\big)\,dxdy-\int_{\O\X Y}(f+\1_{\O\X Y_1}\wh\Theta)\cdot w_0\,dxdy,
						\end{multline*}
						for all $(w_0,w_1,w_2)\in \D$.
					\end{theorem}
					\begin{proof}
						Observe that the unfolded limit energy $\GJ_\d$ has a unique minimizer $(\Gu_{0,\d},\wh U_\d, \wh W_\d)$ such that 
						$$\Gm_\d=\GT_\d(\Gu_{0,\d},\wh U_\d, \wh W_\d)=\min_{(w_0,w_1,w_2)\in \D}\GT_\d(w_0,w_1,w_2).$$
						The proof is divided into three steps.
						
						{\bf Step 1:} We proceed as in Step 1 of the previous theorem using the recovery sequence
						$$  w_\e(x)=\left\{\begin{aligned}
							&w_0(x)+\e w_1\left(x,\left\{{x\over \e}\right\}\right),\quad&&\text{for $x\in \O^2_\e$},\\
							&w_0(x)+\e w_1\left(x,\left\{{x\over \e}\right\}\right)+{\e}w_2\left(x,\left\{{x\over \e}\right\}\right),\quad&&\text{for $x\in\O^1_\e$},
						\end{aligned}\right.$$
						for $(w_0,w_1,w_2)\in \D \cap \left(\C^1(\O)^N\X \C^1(\wo{\O}\X \wo{Y_1})^N\X\C^1(\wo{\O}\X \wo{Y})^N\right)$.
						Then, using Lemma \ref{L5}, we get
						$$	\limsup_{\e\to0}\Gm_\ed\leq \GT_\d(w_0,w_1,w_2)\,\quad\forall\,(w_0,w_1,w_2)\in\D.$$
						{\bf Step 2:}
						Using the Lemmas \ref{L2}--\ref{L3}, we have $\Gu_\e\in \GU$ and $\Gw_\e\in H^1(\O)^N$ with $\Gw_\e=0$ in $\O^2_\e\cup \Lambda_\e$ such that $\wo u_\e=\Gw_\e+\Gu_\e$ and 	
						$$\d\|\Gw_\e\|_{L^2(\O)}+\d\e\|\nabla \Gw_\e\|_{L^2(\O)}\leq C\e,\quad\|\Gu_\e\|_{H^1(\O)}\leq C.$$
						Then, $\e\to 0$ with $\d$ being fixed give there exist $\Gu_{0,\d}\in\GU$, $\wh U_\d\in L^2(\O; H^1_{per,0}(Y))^N$ and $\wh W_\d\in L^2(\O; \GH^1_{per,0}(Y))^N$, such that upto a subsequence
						$$\begin{aligned}
							\Te(e(\Gu_\e))& \rightharpoonup e(\Gu_{0,\d})+e_y(\wh U_\d),\quad &&\text{weakly in $L^2(\O\X Y_2)^{N\X N}$},\\
							\Te(e(\Gw_\e))& \rightharpoonup e_y(\wh W_\d),\quad &&\text{weakly in $L^2(\O\X Y_1)^{N\X N}$}.
						\end{aligned}$$
						So, using the above convergences  along with non-negativity of $\sigma_\ed(\cdot):e(\cdot)$ and weak lower semi continuity of $\GT_\d$ give
						$$\GT_\d(\Gu_{0,\d},\wh U_\d, \wh W_\d)\leq \liminf_{\e\to0}\GT_\ed(\wo u_\e)\leq \limsup_{\e\to0}\GT_\ed(\wo u_\e)\leq \GT_\d(\Gu_{0,\d},\wh U_\d, \wh W_\d),$$
						which imply
						$$\lim_{\e\to0}\Gm_\ed=\Gm_\d.$$
						{\bf Step 3:} Similarly, as in Step 3 of previous with recovery sequence $(w_0,w_1,{1\over \d}w_2)$ give
						$$\limsup_{\d\to0}\Gm_\d\leq \GT^{hom}(w_0),\quad\forall\,w_0\in\GU.$$
						From Step 2, we get
						$$\|\wh W_\d\|_{L^2(\O; H^1(Y_1))}\leq {C\over \d},\quad \|\Gu_{0,\d}\|_{H^1(\O)}+\|\wh U_\d\|_{L^2(\O; H^1(Y))}\leq C,$$
						which give as in Step 3 of Theorem \ref{v03}
						$$\GT^{hom}(\Gu_0)\leq \liminf_{\d\to0}\Gm_\d.$$
						Finally, we obtain
						$$\lim_{\d\to0}\Gm_\d=\Gm.$$
						%								Since, $\d\|\wh W_\d\|_{L^2(\O; H^1(Y_1))}\leq C$, there exist a $\wh W\in L^2(\O; \GH^1_{per,0}(Y_1))^N$ such that upto a subsequence 
						%								$$\d\wh W_\d \rightharpoonup \wh W,\quad \text{weakly in $L^2(\O; H^1(Y_1))^N$}.$$
						%								Then, passing $\d\to0$ in \eqref{ER02} with test functions $(w_0,w_1,{1\over \d}w_2)$ give \eqref{L03}. Hence, similarly like above, we have for the whole sequence
						%								$$\lim_{\d\to0}\Gm_\d=\Gm.$$
						Due to the existence of unique solutions in each step, we have the convergences for the whole sequence. This completes the proof.
					\end{proof}
%					\subsubsection*{\bf Case 2: First $\d\to0$ with $\e$ being fixed then $\e\to0$}
%					\begin{theorem}\label{v03}
%						We have 
%						$$\lim_{\e\to0}\limsup_{\d\to0}\Gm_\ed\leq\lim_{\e\to0}\Gm_\e= \Gm,$$
%						where
%						$$\Gm_\e=\min_{w\in \GU}\GT_\e(w),$$
%						with
%						$$\GT_\e(w)={1\over 2}\int_{\O^2_\e}A_\e(x)e(w):e(w)\,dx-\int_\O f\cdot w\,dx-\int_{\O^1_\e}\Theta_\e\cdot w\,dx,\quad\forall\, w\in\GU.$$	
%					\end{theorem}
%					\begin{proof}			
%						Observe that, using Korn's inequality and Lax-Milgram lemma, we have there exist a unique $\Gu_{0,\e}\in \GU$ such that 
%						$$\Gm_\e=\GT_\e(\Gu_{0,\e})=\min_{w\in \GU}\GT_\e(w).$$
%						Observe that there exist $c_1>0$ such that
%						$$-c_1\leq \GT_\e(\Gu_{0,\e})=\Gm_\e\leq \GT_\e(0)=0.$$
%						The proof is divided into two steps.
%						
%						{\bf Step 1:} We prove for $\e$ being fixed
%						\begin{equation}\label{v01}
%							\limsup_{\d\to0}\Gm_\ed\leq \GT_\e(w),\quad\forall\,w\in \GU,
%						\end{equation}
%						Observe that due to density argument it is enough to prove \eqref{v01} for every $w\in \GU\cap \C^1(\wo{\O})^N$. Let $w\in \GU\cap \C^1(\wo{\O})^N$, we define the recovery sequence as $w=w_{\d}\in \GU$. Then ,we have
%						$$\begin{aligned}
%							e(w_\d)\to e(w),\quad\text{strongly in $L^2(\O)^{N\X N}$},
%						\end{aligned}$$
%						and $\Gm_\ed\leq \GT_\ed(w_{\d})$. So we obtain
%						$$\lim_{\d\to0}\GT_\ed(w_{\d})=\int_{\O^2_\e}A_\e(x)e(w):e(w)\,dx-\int_\O f_\e\cdot w\,dx=\GJ_\e(w),$$
%						where $f_\e=f+\wt\Theta_\e$ ($\wt \Theta_\e$ is extension of $\Theta_\e$ by zero in $\O^2_\e$), which imply
%						$$\limsup_{\d\to0}\Gm_\ed\leq \limsup_{\d\to0}\GT_\ed(w_\d)= \GT_\e(w).$$
%						So, we have \eqref{v01}.
%						
%						{\bf Step 2:} Now, proceeding as in Step 1, we pass to $\e\to0$ in $\Gm_\e$ for that we use the recovery sequence defined as
%						$$w_\e(x)=w_0+\e w_1\left(x,\left\{{x\over \e}\right\}\right),\quad\text{for a.e. $\O$},$$
%						which, together with Lemma \ref{L5}, gives
%						$$\limsup_{\e\to0}\Gm_\e\leq \GT^{hom}(w),\quad\forall\,w\in\GU.$$
%						Now, using the non-negativity of $A_\e e(\cdot):e(\cdot)$ and weak lower semi-continuity of convex functional $\GT^{hom}$, we obtain (as in Lemma \ref{L5})
%						$$\GT^{hom}(\Gu_0)\leq \liminf_{\e\to0}\GT_\e(\Gu_{0,\e})\leq \limsup_{\e\to0}\GT_\e(\Gu_{0,\e})\leq \GT^{hom}(\Gu_0).$$
%						Hence, we obtain
%						$$\lim_{\e\to0}\Gm_\e=\Gm.$$
%						Due to the existence of unique solutions in each step, we have the convergences for the whole sequence. This completes the proof.
%					\end{proof}
					%%%%%%%%%%%%%%%%%%%%%%%%%%%%%%%%%%%%%%%%%%%%%%%%%%%%%%%%%%%%%%%%%%%%%%%%%%%%%%%%%%%%%%%%%%%%%%%%%%%%%%%%%%%%%%%%%%%%%%%%%%%%%%%%%%%%%%%%%%%%%%%%%%%%%%%%%%%%%%%%%%%%%%%%%%%%%%%%%%%%%%%%%%%%%%%%%%%%%%%%%%%%%%%%%%%%%%%%%%%%%%%%%%%%%%%%%
					%			\begin{theorem}\label{v03}
						%				We have
						%				$$\lim_{\e\to0}\limsup_{\d\to0}\Gm_\ed\leq\lim_{\e\to0}\Gm_\e= \Gm,$$
						%				where
						%				$$\Gm_\e=\min_{(w_2,w_1)\in \GU\X \GH^1(\O)}\GT_\e(w_1,w_2),$$
						%				with
						%				\begin{multline*}
							%					\GT_\e(w_1,w_2)={1\over 2}\int_{\O^1_\e}A_\e(x)e(w_1):e(w_1)\,dx+{1\over 2}\int_{\O^2_\e}A_\e(x)e(w_2):e(w_2)\,dx\\
							%					-\int_\O (f+ \Theta_\e)\cdot (w_2+{1\over \kappa}w_1)\,dx,\quad\forall\, (w_1,w_2)\in\GU\X \GH^1(\O),
							%				\end{multline*}
						%				where
						%				$$\GH^1(\O)=\left\{w\in H^1(\O)^N\,|\, w=0,\quad\text{a.e. in $\O^2_\e\cup \partial\O^1_\e$}\right\}.$$
						%			\end{theorem}
					%			\begin{proof}			
						%				Observe that, using Korn's inequality and Lax-Milgram lemma, we have there exist a unique $(\Gu_{2,\e},\Gu_{1,\e})\in \GU\X \GH^1(\O)$ such that 
						%				$$\Gm_\e=\GT_\e(\Gu_{1,\e},\Gu_{2,\e})=\min_{(w_2,w_1)\in \GU\X \GH^1(\O)}\GT_\e(w_1,w_2).$$
						%				Observe that there exist $c_1>0$ such that
						%				$$-c_1\leq \GT_\e(\Gu_{1,\e},\Gu_{2,\e})=\Gm_\e\leq \GT_\e(0,0)=0.$$
						%				$\bullet$ {\bf Step 1:} We prove for $\e$ being fixed
						%				\begin{equation}\label{v01}
							%					\limsup_{\d\to0}\Gm_\ed\leq \GT_\e(w_1,w_2),\quad\forall\,(w_2,w_1)\in \GU\X \GH^1(\O),
							%				\end{equation}
						%				Observe that due to density argument it is enough to prove \eqref{v01} for every $w\in \GU\cap \C^1(\wo{\O})^N$. Let $w_2\in \GU\cap \C^1(\wo{\O})^N$ and $w_1\in \GH^1(\O)\cap \C^1(\wo{\O})^N$, we define the recovery sequence as $w_{\d}=\d w_1+ w_2\in \GU$. Then ,we have
						%				$$\begin{aligned}
							%					e(w_\d)\to e(w_2),\quad\text{strongly in $L^2(\O^2_\e)^{N\X N}$},\\
							%					{1\over \d}e(w_\d)\to e(w_1),\quad\text{strongly in $L^2(\O^1_\e)^{N\X N}$},
							%				\end{aligned}$$
						%				and $\Gm_\ed\leq \GT_\ed(w_{\d})$. So we obtain
						%				$$\lim_{\d\to0}\GT_\ed(w_{\d})=\int_{\O^2_\e}A_\e(x)e(w):e(w)\,dx-\int_\O f_\e\cdot w\,dx=\GJ_\e(w),$$
						%				where $f_\e=f+\wt\Theta_\e$ ($\wt \Theta_\e$ is extension of $\Theta_\e$ by zero in $\O^2_\e$), which imply
						%				$$\limsup_{\d\to0}\Gm_\ed\leq \limsup_{\d\to0}\GT_\ed(w_\d)= \GT_\e(w).$$
						%				So, we have \eqref{v01}.\\[1mm]
						%				$\bullet$ {\bf Step 2:}
						%Using the Lemmas \ref{L2}--\ref{L3}, we have $\Gu_\e\in \GU$ and $\Gw_\e\in H^1(\O)^N$ with $\Gw_\e=0$ in $\O^2_\e\cup \partial \O^1_\e$ such that $\wo u_\e=\Gw_\e+\Gu_\e$ and 	
						%$$\d\|\Gw_\e\|_{L^2(\O)}+\d\e\|\nabla \Gw_\e\|_{L^2(\O)}\leq C\e,\quad\|\Gu_\e\|_{H^1(\O)}\leq C.$$
						%Then, $\d\to 0$ with $\e$ being fixed give there exist $\Gu_{0,\e}\in\GU$ and $\Gw_{0,\e}\in H^1(\O)^N$ with $\Gw_{0,\e}=0$ in $\O^2_\e\cup \partial \O^1_\e$, such that upto a subsequence
						%$$\begin{aligned}
							%	e(\Gu_\e)& \rightharpoonup e(\Gu_{0,\e}),\quad &&\text{weakly in $L^2(\O)^{N\X N}$},\\
							%	{1\over \d}e(\Gw_\e)& \rightharpoonup e(\Gw_{0,\e}),\quad &&\text{weakly in $L^2(\O)^{N\X N}$}.
							%\end{aligned}$$
							%So, using the above convergences  along with non-negativity of $\sigma_\ed(\cdot):e(\cdot)$ and weak lower semi continuity of $\GT_\e$ give
							%$$\GT_\e(\Gu_{0,\e})\leq \liminf_{\d\to0}\GT_\ed(\wo u_\e)\leq \limsup_{\d\to0}\GT_\ed(\wo u_\e)\leq \GT_\e(\Gu_{0,\e},\wh U_\e, \wh W_\e),$$
							%which imply
							%$$\lim_{\e\to0}\Gm_\ed=\Gm_\d.$$
							%													
							%		$$\GJ_\e(\Gu_\e)\leq \liminf_{\d\to0}\GJ_\ed(\Gu_\ed) \leq \limsup_{\d\to0}\GJ_\ed(\Gu_\ed)\leq \GJ_\e(w),\quad\forall\,w\in\GU$$
							%	which imply
							%		$$\lim_{\d\to0}\Gm_\ed=\Gm_\e=\GJ_\e(\Gu_\e).$$
							%		
							%				$\bullet$ {\bf Step 2:} Now, proceeding as in Step 1, we pass to $\e\to0$ in $\Gm_\e$ for that we use the recovery sequence defined as
							%				$$w_\e(x)=w_0+\e w_1\left(x,\left\{{x\over \e}\right\}\right),\quad\text{for a.e. $\O$},$$
							%				which give with Lemma \ref{L5}
							%				$$\limsup_{\e\to0}\Gm_\e\leq \GT^{hom}(w),\quad\forall\,w\in\GU.$$
							%				Now, using the non-negativity of $A_\e e(\cdot):e(\cdot)$ and weak lower semi-continuity of convex functional $\GT^{hom}$ give as in Lemma \ref{L5}
							%				$$\GT^{hom}(\Gu_0)\leq \liminf_{\e\to0}\GT_\e(u_\e)\leq \limsup_{\e\to0}\GT_\e(\Gu_\e)\leq \GT^{hom}(\Gu_0).$$
							%				Hence, we obtain
							%				$$\lim_{\e\to0}\Gm_\e=\Gm.$$
							%				Due to the existence of unique solutions in each step, we have the convergences for the whole sequence. This completes the proof.
							%				
							%				%								 \eqref{L02}. Then, proceeding as in the previous section, we get \eqref{425} as the homogenized limit, which has a unique solution satisfying
							%				%								$$\Gm=\GJ^{hom}(\Gu_0)=\min_{w_0\in\GU}\GJ^{hom}(w_0),$$
							%				%								where
							%				%								$$ \GJ^{hom}(w_0)={1\over |Y|}\int_\O A^{hom}e(w_0):e(w_0)\,dx-\int_\O \left(f+{|Y_1|\over |Y|}\Mc_{Y_1}(\wh \Theta)\right)\cdot w_0\,dx,\quad\forall\,w_0\in\GU.$$
							%				%								Similarly, like before we have for the whole sequence and $\kappa\in[0,+\infty]$
							%				%								\begin{equation}
								%					%									\Gm=\lim_{\e\to0}\Gm_\e.
								%					%								\end{equation}
							%			\end{proof}
						%%%%%%%%%%%%%%%%%%%%%%%%%%%%%%%%%%%%%%%%%%%%%%%%%%%%%%%%%%%%%%%%%%%%%%%%%%%%%%%%%%%%%%%%%%%%%%%%%%%%%%%%%%%%%%%%%%%%%%%%%%%%%%%%%%%%%%%%%%%%%%%%%%%%%%%%%%%%%%%%%%%%%%%%%%%%%%%%%%%%%%%%%%%%%%%%%%%%%%%%%%%%%%%%%%%%%%%%%%%%%%%%%%%%%%%%%%%%%%%%%%%%%%%%%%%%%%%%%%%%%%
						Finally, we have the following interchange of limits
						\begin{equation}\label{CUM01}
							\boxed{
								\lim_{\e\to0}\limsup_{\d\to0}\Gm_\ed\leq \Gm=	\lim_{\d\to0}\lim_{\e\to0}\Gm_\ed=\lim_{(\e,\d)\to(0,0)}\Gm_\ed.}
						\end{equation}
						{\clm -----------------------------------------------------------------------------------------------------------------------------------------------------}
						
					
					
					%\newpage		
					%{\color{magenta}From the variational form,  $\widetilde{\widetilde{\chi_{ij}^2}}=-\chi_{ij}^2,$  right? }
					%\section{Supplementary}
					%											Let us set
					%											$$\GU=\left\{v\in \GU| v=0,\quad \text{in $\O^2_\e\cup \partial\O^1_\e$}\right\},$$
					%											and the total energies
					%											$$\left.\begin{aligned}
						%												\GJ^H_\ed(w)&=\int_{\O^2_\e}\GA_\ed e(w):e(w)\,dx-\int_\O f_\e\cdot w\,dx,\\
						%												\GJ^S_\ed(w)&=\int_{\O^1_\e}\GA_\ed e(w):e(w)\,dx-\int_{\O^1_\e}f_\e\cdot w\,dx,												\end{aligned}\right\}\quad\forall\,w\in\GU$$
					%											then, we have
					%											$$\GJ_\ed(u_\ed)=\GJ_\ed^H(\Gu_\ed)+\GJ^S_\ed(\Gw_\ed),\quad \GJ_\ed(\Gu_\ed)=\GJ_\ed^H(\Gu_\ed).$$
					%											There exist $(U_\ed,V_\ed)\in \GU\X \GU$ such that	
					%											$$\GJ_\ed^H(U_\ed)+\GJ^S_\ed(V_\ed)=\min_{(U,V)\in \GU\X \GU}(\GJ^H_\ed(U)+\GJ^S_\ed(V)).$$	
					%											{\clm 								$\bullet$ {\bf Step 2:} We prove 
						%												\begin{equation}
							%													\GJ_\e(\Gu_\e)\leq \liminf_{\d\to0}\Gm_\ed.
							%												\end{equation}
						%												Using the Lemmas \ref{L2}--\ref{L3}, we have $\Gu_\ed\in \GU$ and $\Gw_\ed\in H^1(\O)^N$ with $\Gw_\ed=0$ in $\O^2_\e\cup \partial \O^1_\e$ such that $u_\ed=\Gw_\ed+\Gu_\ed$ and 	
						%												$$\d\|\Gw_\ed\|_{L^2(\O)}+\d\e\|\nabla \Gw_\ed\|_{L^2(\O)}\leq C\e,\quad\|\Gu_\ed\|_{H^1(\O)}\leq C.$$
						%												Then,  as $\d\to0$ with $\e$ being fixed give that there exist $\Gw_\e,\Gu_\e\in \GU$ such that (upto a subsequence)
						%												\begin{equation}\label{v02}
							%													\begin{aligned}
								%														\d\Gw_\ed &\rightharpoonup \Gw_\e,\quad&&\text{weakly in $H^1(\O)^N$},\\
								%														\Gu_\ed &\rightharpoonup \Gu_\e,\quad&&\text{weakly in $H^1(\O)^N$}.
								%													\end{aligned}
							%												\end{equation}
						%												Also using the estimate \eqref{E01}$_1$, we have that there exist $c_2>0$ such that
						%												\begin{equation}
							%													-c_2\leq \GJ_\ed(u_\ed)=\Gm_\ed\leq \GJ_\ed(0)=0.
							%												\end{equation}
						%												{\bf Step 2.1:} Let us set
						%												$$\GJ_\ed^*(u_\ed)=\GJ_\ed(u_\ed)+\int_{\O^1_\e} f_\e\cdot \Gw_\ed.$$
						%												Then, we have using the convergences \eqref{v02}
						%												$$\liminf_{\d\to0}\GJ_\ed^*(u_\ed)=\liminf_{\d\to0}\left(\GJ_\ed(u_\ed)+\int_{\O^1_\e}f_\e\cdot \Gw_\ed\,dx\right)=\GJ_\e(\Gu_\e)+\int_{\O^1_\e}A_\e e(\Gw_\e):e(w_\e)\,dx.$$					
						%												We also have from Step 1,
						%												$$\liminf_{\d\to0}\GJ_\ed(u_\ed)\leq \limsup_{\d\to0}\GJ_\ed(u_\ed)\leq \GJ_\e(\Gu_\e).$$
						%												So, we get
						%												$$\limsup_{\d\to0}\left(\GJ_\ed(u_\ed)+\int_{\O^1_\e}f_\e\cdot \Gw_\ed\,dx\right)\leq \GJ_\e(\Gu_\e)+\limsup_{\d\to0}\int_{\O^1_\e} f\cdot\Gw_\ed\,dx.$$
						%												Comparing the above inequalities, we obtain
						%												$$\GJ_\e(\Gu_\e)+\int_{\O^1_\e}A_\e e(\Gw_\e):e(w_\e)\,dx\leq \GJ_\e(\Gu_\e)+\limsup_{\d\to0}\int_{\O^1_\e} f\cdot\Gw_\ed\,dx,$$
						%												which give $X_\e\leq 0$ 	where
						%												$$X_\e=\limsup_{\d\to0}X_\ed=\limsup_{\d\to0}\left(\int_{\O^1_\e}\GA_\ed e(\Gw_\ed):e(\Gw_\ed)\,dx-\int_{\O^1_\e}f_\e\cdot \Gw_\ed\,dx\right).$$
						%												{\bf Step 2.2:} 
						%												%									Using Korn's inequality we have
						%												%									$$\d^2\|e(\Gw_\ed)\|_{L^2(\O^1_\e)}\leq C\|f_\e|\|e(\Gw_\ed)\|_{L^2(\O^1_\e)},$$
						%												%									 as $\GX^1_\ed(\Gw_\ed)\leq \GX^1_\ed(0)=0$
						%												Observe that we also have
						%												$$\GJ_\e(\Gu_\e)+X_\e=\liminf_{\d\to0}\GJ_\ed(u_\ed),$$
						%												which imply
						%												$$\GJ_\e(\Gu_\e)= \liminf_{\d\to0}\GJ_\ed(u_\ed)+(-X_\e).$$
						%												Since $-X_\e\geq 0$, we get
						%												$$\GJ_\e(\Gu_\e)\leq \liminf_{\d\to0}\GJ_\ed(u_\ed).$$
						%												Together with Step 1, we get
						%												$$\GJ_\e(\Gu_\e)\leq \liminf_{\d\to0}\GJ_\ed(u_\ed)\leq \limsup_{\d\to0}\GJ_\ed(u_\ed)\leq \GJ_\e(\Gu_\e).$$
						%												This give
						%												$$\lim_{\d\to0}\Gm_\ed=\Gm_\e.$$}									

			
			{\clm 
				\subsection{Special behavior: First $(\e,\d)\to(0,0)$ with $\kappa\in(0,+\infty)$ then $\kappa\to0$}
				{\clb I think, passing to the limit $\kappa\to 0$ may not be possible. Sinec, for each $\kappa>0$, corresponding solution $\wh{W}_k$ of \eqref{5171} is not a bounded sequence in $H^1_{0,per}(Y)$. We do not have even uniform bound for $\|\wh{W}_k\|_{L^2(\O)}$ with respect to $\kappa$, which is the minimum requirement.} {\clm Yes.}\\
				In the cell problem \eqref{5171}, if we pass $\kappa\to0$, we obtain
				\begin{equation}\label{5173}
					\int_{\O\X Y_1}\big(f+\wh\Theta\big)\cdot w_2\,dxdy=0,\quad\forall \,w_2\in L^2(\O;H^1_{per,0}(Y_1))^N,
				\end{equation}
				which imply
				$$f+\wh\Theta=0,\quad\text{a.e. in $\O\X Y_1$}.$$
				finally, the homogenized problem is given by
				\begin{equation}\label{426}
					{1\over |Y|}\int_\O A^{hom}e(\Gu_0):e(w_0)\,dx=\left(1-{|Y_1|\over |Y|}\right)\int_{\O}f\cdot w_0\,dx,
				\end{equation}
				for all $w_0\in \GU$, where $A^{hom}$ is given by \eqref{Hom01}, which is the same if we have started with only the perforated domain (no soft part). Hence, we only analyze the asymptotic behavior of the OCPs when $\kappa\in(0,+\infty]$.}
			
			
			
			
			
			
			
			
			
			
			Let $v_\e$ be a sequence in $H^1(\O)$ such that
			$${1\over \e}\|v_\e\|_{L^2(\O)}+\|\nabla v_\e\|_{L^2(\O)}\leq C.$$
			Then, we have
			$$\begin{aligned}
				{1\over \e}v_\e &\rightharpoonup v,\quad &&\text{weakly in $L^2(\O)$},\\
				\nabla v_\e & \rightharpoonup 0,\quad&&\text{weakly in $L^2(\O)^N$},\\
				\Te({1\over \e} v_\e) & \rightharpoonup v+\wh v,\quad &&\text{weakly in $L^2(\O; H^1_{per,0}(Y))$},\\
				\Te(\nabla v_\e)={1\over \e}\nabla_y\Te(v_\e) & \rightharpoonup\nabla_y \wh v,\quad &&\text{weakly in $L^2(\O\X Y)^N$}.
			\end{aligned}$$
			
			
			% and the weak convergences \eqref{II02} are actually strong. Finally, using the inequality \eqref{Ex01} and the Lemma \ref{NEq1},  we have \eqref{424}.
			%	\section{Properties of Classical unfolding operator}
			
			%%%%%%%%%%%%%%%%%%%%%%%%%%%%%%%%%%%%%%%%%%%%%%%%%%%%%%%%%%%%%%%%%%%%%%%%%%%%%%%%%%%%%%%%%%%%%%%%%%%%%%%%%%%%%%%%%%%%%%%%%%%%%%%%%%%%%%%%%%%%%%%%%%%%%%%%%%%%%%%%%%%%%%
			
			
			{\clm {\Large The proof is simple, We will provide it or not will depend on the number of total pages.}
				\begin{corollary}
					Let $\{u_\e\}_\e$ be sequence of unique weak solution of the problem \eqref{MainW01} in $\GU$ corresponding to $\Theta_\e\in L^2(\O^1_\e)^N$. Then, we have following strong convergences for $\kappa\in[0,+\infty]$ 
					\begin{equation}
						\begin{aligned}
							\Te\big(e(u_\e)\big) &\to e(\Gu_0)+e_y(\wh U),\quad&&\text{strongly in $L^2(\O\X Y_2)^{N\X N}$},\\
							\d\Te\big(e(u_\e)\big)& \to e_y(\wh W),\quad&&\text{strongly in $L^2(\O\X Y_1)^{N\X N}$},
						\end{aligned}
					\end{equation}
					where $(\Gu_0,\wh U,\wh W)\in\D$ is the unique solution of the limit unfolded problems \eqref{L01}--\eqref{L03}.
				\end{corollary}
			}
			
			
			
			
			
			
			
			
			\begin{itemize}
				\item When $\kappa=+\infty$, the duplet $(\Gu_0,\Gv_0)\in [\GU]^2$ satisfies the following optimality system
				\begin{align}\label{lim-ops-case 3+}
					\left.\begin{aligned}
						&\int_{\O} A^{hom}e(\Gu_0):e(w_0)\,dx
						=|Y|\int_{\O} f\cdot w_0\,dx-\frac{|Y_1|}{\gamma}\int_{\O}\Gv_0\cdot w_0\,dx,\\ &
						\int_{\O} A^{hom}e(\Gv_0):e(w_0)\,dx
						=|Y|\int_{\O} (\Gu_0-u_0)\cdot  w_0\,dx
					\end{aligned}\right\}\forall w_0\in \GU.
				\end{align}
			\end{itemize}  
			
			\newpage
			%%%%%%%%%%%%%%%%%%%%%%%%%%%%%%%%%%%%%%%%%%%%%%%%%%%%%%%%%%%%%%%%%%%%%%%%%%%%%%%%%%%%%%%%%%%%%%%%%%%%%%%%%%%%%%%%%%%%%%%%%%%%%%%%%%%%%%%%%%%%%%%%%%%%%%%%%%%%%%%%%%%%%%
			%Supplementary part for future.
			\section{Unfolding for small holes}
			Now, we define the unfolding operator for function defined on $\O^1_\e$\\
			{\clr This operator is more general then what is given in \cite{CDG} for small holes. So, details has to given when $\lim_{(\e,r)\to(0,0)}{r\over \e}=0$.}
			{\clm \begin{definition}[Small unfolding operator]
					For every measurable function $\phi$ on $\O^1_\e$ the small unfolding operator $\Te^r$ is defined by 
					$$ \Te^r(\phi) (x,y) \doteq \begin{aligned}
						&\psi \left(\e\left[{x\over r}\right] + r y\right) \qquad &&\hbox{for a.e. }\; (x,y)\in  \O\times Y_1.
						%		&0\qquad&&\hbox{for a.e. }\; (x',y')\in  \Lambda_\e\times Y.
					\end{aligned}$$
				\end{definition} 
				The small unfolding operator is also a continuous linear operator from $L^2(\O^1_\e)$ into $L^2(\O\X Y_1)$ which satisfies
				$$\|\Te^r(\phi)\|_{L^2(\O\X Y_1)}\leq C\|\phi\|_{L^2(\O^1_\e)}\quad\text{for every}\quad \phi\in L^2(\O^1_\e).$$
				The constants $C$ do not depend on $\e$ and $r$.\\
				Moreover, for every $\psi\in H^1(\O^1_\e)$ one has
				$$ \nabla_y\Te^r(\psi)(x,y)=r\Te^r(\nabla \psi)(x,y) \quad \hbox{a.e. in}\quad \O\X Y_1,$$
				and for $\phi\in H^1(\O)$, one has
				$$\Te\big(\phi_{|\O^1_\e}\big)(x,y)=\Te^r(\phi)(x,y),\quad\text{for $(x,y)\in \O\X Y_1$}.$$
				Using the above operators, we have
				\begin{equation}
					\begin{aligned}
						\|\Te(\Gu_\e)\|_{L^2(\O\X Y)}+\|\Te(\nabla \Gu_\e)\|_{L^2(\O\X Y)}&\leq C,\\
						\|\Te^r(\Gv_\e)\|_{L^2(\O\X Y_1)}+\|\nabla_y\Te^r(\Gv_\e)\|_{L^2(\O\X Y_1)}&\leq {Cr\over \d}.
					\end{aligned}
			\end{equation}}
			\section{Boundary control problem in elasticity}
			%Let us denote stress tensor in vectorial form
			%$$\sigma(u)=\GA e(u),\quad \text{for all $u\in H^1(\O)^N$}.$$
			%or
			%$$\sigma(u)=\lambda(tr(e(u)))\GI_N+2\mu e(u),\quad \text{for all $u\in H^1(\O)^N$ with Lame's constant $\lambda,\mu>0$}.$$
			%
			%where the elasticity tensor $\GA_{ijkl}(\cdot)\in L^\infty(\O)$ and the external force $f=(f_i)$ are given functions of $x=(x_i)$ in a bounded domain $\O\subset \R^N$ with lipchitz boundary $\partial\O$.\\
			%. It is usually assumed that
			%\begin{itemize}
			%	\item The elasticity tensor $\GA_{ijkl}$ satisfy the condition:
			%	\begin{equation}
				%		\GA_{ijkl}=\GA_{jikl}=\GA_{ijlk}=\GA_{klij}.
				%	\end{equation}
			%	\item There exist two positive constants $0<k_0$ such that
			%	\begin{equation}
				%		k_0S_{ij}S_{ij}\leq \GA_{ijkl}S_{ij}S_{kl},\quad\text{for $S=[S_{ij}]\in \R_\sy^{N\X N}$}.
				%	\end{equation}
			%\end{itemize}
			{\clr We consider the following boundary value problem of Dirichlet type in linear elasticity:
				\begin{equation}\label{Main01}
					\left\{\begin{aligned}
						-\nabla \cdot \sigma_\ed(x)&=f(x),\quad &&\text{$x\in\O$},\\
						u_\e(x)&=0,\quad &&\text{$x\in\partial \Gamma_0$},\\
						\sigma_\ed(x)\Gn(x)&=g_\e(x),\quad&&\text{$x\in \Gamma_1$},
					\end{aligned}\right.
				\end{equation}
				where $f=(f_i)$ is the body force and $g_\e=(g_{\e,i})$ is the surface force}.\\
			The problem is to determine the displacement field $u_\e=(u_{\e,i})$\footnote{For simplicity of notation, we index fields only by $\e$} satisfying the system \eqref{Main01}. By virtue of the assumptions on $\GA_\ed$, we define the bilinear form $\Ga_\ed$ by
			\begin{equation}
				\Ga_\ed(u,v)=\int_{\O}\sigma_{\ed,ij}(u)e_{ij}(v)\,dx,\quad \forall\,u,v\in H^1_0(\O)^N.
			\end{equation}
			So, the weak form of the problem \eqref{Main01} is given by: For a given function $f\in L^2(\O)^N$ and $g_\e\in L^2(\Gamma_1)^N$, find the solution $u_\e\in \GU$ such that
			\begin{equation}
				\Ga_\ed(u_\e,v)=\int_\O f\cdot v\,dx+\int_{\Gamma_1} g_\e\cdot v\,ds,\quad \forall\,v\in \GU_\e,
			\end{equation}
			where the set of admissible displacement is given by
			$$\GU=\left\{u\in H^1(\O)^N\,|\; u=0\;\text{on $\Gamma_0$}\right\}.$$
			The using the assumptions and Korn's inequality, we have the bilinear form $\Ga_\ed$ is continuous and coercive in $\GU$. Therefore, according to Lax-Milgram lemma, we have for any $f\in L^2(\O)^N$ and $g_\e\in L^2(\Gamma_1)^N$, there exist a unique solution $u_\e\in \GU_\e\cap H^2(\O)^N$ satisfying
			$$\|u_\e\|_{H^1(\O)}\leq C\{\|f\|_{L^2(\O)}+\|g_\e\|_{L^2(\Gamma_1)}\}.$$
			The unique solution corresponding to the given $g_\e$ will be denoted by $u_\e(g_\e)$.\\
			
			Now let us turn to the formulation of an optimal control problem to obtain a prescribed displacement by means of externally surface forces on the body. Suppose that the equilibrium position of the body has some prescribed form given by a function $u_0(x)\in H^1(\O)^N$, called an objective function. We can
			ask the question of how to determine the function $g_\e$ such that $u_\e(g_\e)$ is as close
			to the objective function $u_0$.\\
			{\clr By introducing the two energy functional defined by
				\begin{equation}
					\begin{aligned}
						\text{$L^2$-energy}\quad			&\GJ_{1,\e}(g_\e)=\|u_\e(g_\e)-u_0\|^2_{L^2(\O)}+{\gamma\over 2}\|g_\e\|^2_{L^2(\Gamma_1)},\\
						\text{Dirichlet-energy}\quad & \GJ_{2,\e}(g_\e)=\|\nabla u_\e(g_\e)-\nabla u_0\|^2_{L^2(\O)}+{\gamma\over 2}\|g_\e\|^2_{L^2(\Gamma_1)}.
					\end{aligned}
				\end{equation}
				The optimal control problem can be formulated as a minimization problem of the cost functional: Determine a function $h_\e\in L^2(\Gamma_1)^N$ such that
				\begin{equation}
					\GJ_{\alpha,\e}(h_\e)=\inf_{g_\e\in L^2(\Gamma_1)^N}\GJ_{\alpha,\e}(g_\e).
			\end{equation}}	
			
			
			
			
			\begin{proof}
				Observe that 
				$$\Gm_{\e,\d}\leq \GJ_{\e,\d}(0)< +\infty\quad\text{which imply}\quad \Gm_{\e,\d}\in[0,+\infty).$$
				%    since $\theta_\e=0\in L^2(\O^1_\e)^N$.\\
				So, there exist a minimizing sequence $\{\theta_{n,\e}\}_n$ in $L^2(\O^1_\e)^N$ such that
				$$\Gm_{\e,\d}=\liminf_{n\to+\infty}\GJ_{\e,\d}(\theta_{n,\e}).$$
				Without loss of generality, we have $\GJ_{\e,\d}(\theta_{n,\e})\leq \GJ_{\e,\d}(0)$, which imply
				\begin{equation}\label{01}
					\|\theta_{n,\e}\|_{L^2(\O^1_\e)}\leq C.
				\end{equation}
				The constant is independent of $n$. So, there exist a subsequence (still denoted by $n$) and $\Theta_\e\in L^2(\O^1_\e)^N$ such that
				$$\theta_{n,\e} \rightharpoonup \Theta_\e,\quad \text{weakly in $L^2(\O^1_\e)^N$}.$$
				Since $L^2$-norm is weak lower semi-continuous, we have
				\begin{equation}\label{02}
					\|\Theta_\e\|_{L^2(\O^1_\e)}\leq \liminf_{n\to+\infty}\|\theta_{n,\e}\|_{L^2(\O^1_\e)}.
				\end{equation}
				Let us set $u_{n,\e}=u_\e(\theta_{n,\e})$ and $u_{0,\e}=u_{\e}(0)$, then using the $L^2$-norm bound of $\theta_{n,\e}$, we have that there exist a constant independent of $n$ such that
				$$\|u_{n,\e}\|_{H^1(\O)}\le C.$$
				So, there exist for the same subsequence as above, we have there exist $u_\e\in\GU$ such that
				$$u_{n,\e} \rightharpoonup u_\e,\quad\text{weakly in $H^1(\O)^N$}.$$
				Now, we show that $u_\e=u_\e(\Theta_\e)$. Since, each $u_{n,\e}$ is a solution of \eqref{Main02} with $\theta_\e=\theta_{n,\e}$. So, we want to prove $u_\e$ is a solution of \eqref{Main02} corresponding to $\theta_\e=\Theta_\e$, for that it is sufficient to prove that
				$$\lim_{n\to+\infty}\int_{\O^1_\e}\theta_{n,\e}\psi=\int_{\O^1_\e}\Theta_\e\psi,\quad\text{for $\psi\in L^2(\O^1_\e)^N$},$$
				which is a straightforward consequence of \eqref{01}.\\
				Again using the semi-lower continuity of $H^1$-norm, we obtain
				$$\|u_\e-u_0\|^2_{H^1(\O)}\leq \liminf_{n\to +\infty}\|u_{n,\e}-u_0\|^2_{H^1(\O)},$$
				which along with \eqref{02} give 
				$$\GJ_{\e,\d}(\Theta_\e)\leq \liminf_{n\to+\infty}\GJ_{\e,\d}(\theta_{n,\e})=\Gm_{\e,\d}.$$
				% {\color{red} In order to have minimization over V, we need to have $\Theta_\varepsilon\in \GV$. The above inequality shows that $\Theta_\varepsilon\in V.$}
				So, $\Theta_\e$ minimizes the energy functional \eqref{MainE01} and the uniqueness follows from the strict convexity of the $L^2$-cost functional. 
				%    {\clm We also have
					%    $$\|\Theta_\e\|^2_{L^2(\O^1_\e)}\leq {2\over \gamma}\GJ_{\e,\d}(0)\leq {2K_0\over \gamma}.$$}
				This completes the proof.
			\end{proof}	
			
			
			
			
			
			\begin{thebibliography}{}
				\bibitem{Donato} Cioranescu, Doina, and Patrizia Donato. An introduction to homogenization. Oxford university press, 1999.
				\bibitem{Olenik}
				Oleïnik, Olga Arsenievna, A. S. Shamaev, and G. A. Yosifian. {\em Mathematical problems in elasticity and homogenization.} Elsevier, 2009.
				
				
				%											
				%	
				%	\bibitem{CiarNecas} Ciarlet, Philippe G., and Jindřich Nečas. "Injectivity and self-contact in nonlinear elasticity." Archive for Rational Mechanics and Analysis 97 (1987): 171-188.
				
				%	\bibitem{BGRods} Blanchard, Dominique, and Georges Griso. "Asymptotic behavior of structures made of straight rods." Journal of Elasticity 108 (2012): 85-118.
				%	
				%	
				%	\bibitem{BGRod} D. Blanchard, G. Griso. Decomposition of deformations of thin rods: application to nonlinear elasticity. Analysis and Applications 7 (01), 21-71, (2009).
				
				
				%	\bibitem{GThinS} Griso, Georges. "Decompositions of displacements of thin structures." Journal de mathématiques pures et appliquées 89, no. 2 (2008): 199-223.
				
				%	\bibitem{ciarletS} Ciarlet, Philippe G. Mathematical elasticity: Volume III: Theory of shells. North-Holland, 2000.
				
				%	\bibitem{gries} Koch, Hannah Carlotta, David Schmelzeisen, and Thomas Gries. "4D textiles made by additive manufacturing on pre-stressed textiles—An overview." In Actuators, vol. 10, no. 2, p. 31. MDPI, 2021.
				%										
				%										
				
				%										
				
				%										
				%										
				%	\bibitem{CDG2}
				%	Cioranescu, Doina, Alain Damlamian, and Georges Griso. \underline{Periodic unfolding and homogenization.} Comptes Rendus Mathematique 335.1 (2002): 99-104.
				%										
				\bibitem{CDG1}
				Cioranescu, Doina, Alain Damlamian, and Georges Griso. \underline{The periodic unfolding method in homogenization.} SIAM Journal on Mathematical Analysis 40.4 (2008): 1585-1620.
				
				\bibitem{tartar1}
				Tartar, Luc. \underline{The general theory of homogenization: a personalized introduction.} Vol. 7. Springer Science and Business Media, 2009.
				
				\bibitem{asymptotic1}
				Bensoussan, Alain, Jacques-Louis Lions, and George Papanicolaou. \underline{Asymptotic analysis for periodic structures.} Vol. 374. American Mathematical Soc., 2011.
				
				\bibitem{Homogenization1}
				Cioranescu, Doina, and Patrizia Donato. \underline{An introduction to homogenization.} Vol. 17. Oxford: Oxford university press, 1999.
				
				\bibitem{twoscale1}
				Nguetseng, Gabriel. \underline{A general convergence result for a functional related to the theory of homogenization.} SIAM Journal on Mathematical Analysis 20.3 (1989): 608-623.
				
				%	\bibitem{twoscale2}
				%	Allaire, Grégoire. \underline{Homogenization and two-scale convergence.} SIAM Journal on Mathematical Analysis 23.6 (1992): 1482-1518.
				
				
				%										
				\bibitem{multi1}
				Panasenko, Grigorii Petrovich. \underline{Multi-scale modelling for structures and composites.} Vol. 3. Dordrecht: Springer, 2005.
				
				%	\bibitem{friesecke02}
				%	Friesecke, Gero, Richard D. James, and Stefan Müller. {\em A theorem on geometric rigidity and the derivation of nonlinear plate theory from three dimensional elasticity.} Communications on Pure and Applied Mathematics: A Journal Issued by the Courant Institute of Mathematical Sciences 55.11 (2002): 1461-1506.
				
				%	\bibitem{masoG}
				%	Dal Maso, Gianni. {\em An introduction to $\Gamma$-convergence. Vol. 8.} Springer Science and Business Media, 2012.
				
				%	\bibitem{ABGamma} A. Braides. {\em Gamma-convergence for beginners. Oxford lecture series in mathematics and its applications}. Oxford University Press, 2002. ISBN 9780198507840.
				%	
				%	\bibitem{Friesecke06} G. Friesecke, R.D. James and S. M{\"{u}}ller. {\em A Hierarchy of Plate Models Derived from Nonlinear Elasticity by Gamma-Convergence}. Archive for Rational Mechanics and Analysis 180 (2), 2006, 183-236.
				%	
				%	\bibitem{LemMu01} Lecumberry, Myriam, and Stefan Müller. "Stability of slender bodies under compression and validity of the von Kármán theory." Archive for rational mechanics and analysis 193.2 (2009): 255-310.
				%	\bibitem{MuPak01} Müller, Stefan, and Mohammad Reza Pakzad. "Regularity properties of isometric immersions." Mathematische Zeitschrift 251 (2005): 313-331.
				
				%	\bibitem{FJMKarman} G. Friesecke, R.D. James and S. M{\"u}ler. {\em The F{\"o}ppl-von K\'{a}rm\'{a}n plate theory as a low energy $\Gamma$-limit of nonlinear elasticity}. Comptes Rendus Mathematique 335, no. 2 (2002), 201-206.
				%												
				\bibitem{elasticity} 
				Ciarlet, Philippe G. {\em Three-dimensional elasticity}. Elsevier, 1988.
				
				\bibitem{Lions} Lions, J. L. (1971). Optimal control of systems governed by partial differential equations.
				
				%	\bibitem{CiarletShell1}
				%	Ciarlet, Philippe G., and Cristinel Mardare. \underline{An introduction to shell theory.} Differential geometry: theory and applications 9 (2008): 94-184.
				%	
				%	\bibitem{BGG1}
				%	Blanchard, Dominique, Antonio Gaudiello, and Georges Griso. {\em Junction of a periodic family of elastic rods with a 3d plate, Part I.} Journal de mathématiques pures et appliquées 88.1 (2007): 1-33.
				%	
				%	\bibitem{BGG2}
				%	Blanchard, Dominique, Antonio Gaudiello, and Georges Griso. {\em Junction of a periodic family of elastic rods with a thin plate, Part II.} Journal de mathématiques pures et appliquées 88.2 (2007): 149-190.
				
				\bibitem{CDG}  Cioranescu, D., Damlamian,  A. and Griso G.: {\em The Periodic Unfolding Method. Theory and Applications to Partial Differential Problems.} Springer, Singapore (2018).  10.1007/978-981-13-3032-2.
				
				%	\bibitem{DGJunction1}
				%	Blanchard, Dominique, and Georges Griso. \underline{Asymptotic behavior of a structure made by a plate and a straight rod.} Chinese Annals of Mathematics, Series B 34 (2013): 399-434.
				
				%	\bibitem{NV} S. Neukamm, I. Vel\v{c}i\'{c}. {\em Derivation of the homogenized von K\'{a}rm\'{a}n plate theory from 3d elasticity}. Mathematical Models and Methods in Applied Sciences 23 (2013), 2701-2748.
				%													
				%	\bibitem{Neu1}
				%	Böhnlein, Klaus, Stefan Neukamm, David Padilla-Garza, and Oliver Sander. {\em A homogenized bending theory for prestrained plates.} Journal of Nonlinear Science 33, no. 1 (2023): 1-90.
				
				%	\bibitem{Velcic1}
				%	Velčić, Igor. "On the general homogenization of von Kármán plate equations from three-dimensional nonlinear elasticity." Analysis and Applications 15, no. 01 (2017): 1-49.
				
				%													\bibitem{NeuRP}
				%													Bauer, Robert, Stefan Neukamm, and Mathias Schäffner. "Derivation of a homogenized bending–torsion theory for rods with micro-heterogeneous prestrain." Journal of Elasticity 141, no. 1 (2020): 109-145.
				
				%	\bibitem{GOW}
				%	Griso, G., Orlik, J. and Wackerle S.
				%	\newblock {\em Asymptotic behavior for textiles.}
				%	\newblock { SIAM Journal of Mathematical Analysis}, 52(2):1639--1689, 2020.
				
				%	\bibitem{GFOW}
				%	Falconi, R., Griso, G., Orlik, J., Wackerle S.
				%	\newblock {\em Asymptotic behavior for textiles with loose contact}.
				%	\newblock {Mathematical Methods in the Applied Sciences (2022)}.
				%	
				%	\bibitem{GFOW2}	
				%	Falconi, R., Griso, G., Orlik, J.
				%	\newblock {\em 	Periodic unfolding for lattice structures}.
				%	\newblock {Ricerche di Matematica}, 1-35 (2023).
				
				%\bibitem{GFOW3}	
				%Falconi, R., Griso, G., Orlik, J.
				%\newblock {\em 	Periodic unfolding for anisotropically bounded sequences}.
				%\newblock {Asymptotic Analysis} 132 (3-4), 383-407,  (2023).
				
				
				%	\bibitem{hauck2}
				%	Griso, Georges, Michael Hauck, and Julia Orlik. {\em Asymptotic analysis for periodic perforated shells.} ESAIM: Mathematical Modelling and Numerical Analysis 55.1 (2021): 1-36.
				
				%	\bibitem{larysa1}
				%	Griso, G., Khilkova, L., Orlik, J. and Sivak O.
				%	{\em Homogenization of Perforated Elastic Structures.}
				%	\newblock { Journal of Elasticity}, 141:181--225, 2020.	
				
				%	\bibitem{GOW2}
				%	Griso, G., Orlik, J. and Wackerle Stephan.
				%	\newblock {\em Asymptotic behavior for textiles in von-Kármán regime},
				%	{ Journal de Mathématiques Pures et Appliquées
					%	}, 2020.	
				%	%										
				%	%										
				%	%										\bibitem{BGRod} D. Blanchard, G. Griso. {\em Decomposition of deformations of thin rods: application to nonlinear elasticity.} Analysis and Applications 7 (01), 21-71, (2009).
				%	%										
				%	\bibitem{Shell1} D. Blanchard, G. Griso. {\em Decomposition of the deformations of a thin shell. Asymptotic behavior of the Green-St Venant's strain tensor}.
				%	Journal of Elasticity 101 (2), 2010, 179-205.
				%										
				%	\bibitem{GKL} G. Griso. {\em Decomposition of plate displacements via Kirchhoff-Love  displacements}.   https://hal.science/hal-03938946.
				%											
				
				
				%											\bibitem{berdichevsky}
				%											Berdichevsky, Victor, and V. L. Berdichevsky. \underline{Variational principles.} Springer Berlin Heidelberg, 2009.
				%											
				%											\bibitem{LecMuller} Lecumberry, Myriam, and Stefan Müller. {\em Stability of slender bodies under compression and validity of the von Kármán theory.} Archive for rational mechanics and analysis 193, no. 2 (2009): 255-310.
				%										
				%											\bibitem{NPE} Antman, Stuart S. {\em Problems in nonlinear elasticity.} Nonlinear Problems of Elasticity (2005): 513-584.
				%											
				
				%										
				%	\bibitem{pretrain01} Eckart, Carl. "The thermodynamics of irreversible processes. III. Relativistic theory of the simple fluid." Physical review 58, no. 10 (1940): 919.
				%									
				%	\bibitem{prestrain02} Lee, Erastus H. "Elastic-plastic deformation at finite strains." (1969): 1-6.
				
				%	\bibitem{IniStress01} Riccobelli, Davide, Abramo Agosti, and Pasquale Ciarletta. "On the existence of elastic minimizers for initially stressed materials." Philosophical Transactions of the Royal Society A 377.2144 (2019): 20180074.
				
				%	\bibitem{IniStress02} Shams, Moniba, Michel Destrade, and Ray W. Ogden. "Initial stresses in elastic solids: constitutive laws and acoustoelasticity." Wave Motion 48.7 (2011): 552-567.
				
				\bibitem{Amartya} Amartya Chakrabortty, Georges Griso, and Julia Orlik. "Dimension reduction and homogenization of composite plate with matrix pre-strain." Asymptotic Analysis Preprint (2024): 1-56.
				%	
				%	\bibitem{GFOW4} Falconi, Riccardo, Georges Griso, and Julia Orlik. "Asymptotic behavior for nonlinear textiles with glued yarns." (2023).
				
				
				
				
				%%%%%%%%%%%%%%%%%%%%%%%%%%%%%%%%%%%%%%%%%%%%%%%%%%%%%%%%%%%%%%%%%%%%%%%%%%%%%%%%%%%%%%%%%%%%%%%%%%%%%%%%%%%											
				
				%										%		
			\end{thebibliography}
